%% Theory of Informational Relations — Monograph v12.4 Relational Zero and Graph Reconciliation
%% Dependency-ordered publication architecture
%% Adrian Lipa
%% Version 12.4 Relational Zero and Graph Reconciliation — 28 September 2026
%% Historical v12.0 content-migration baseline — 30 August 2026

\documentclass[11pt,a4paper,twoside]{book}
\usepackage[utf8]{inputenc}
\usepackage[T1]{fontenc}
\usepackage{lmodern}
\usepackage{amsmath,amssymb,amsthm,bm,mathtools}
\usepackage{mathrsfs}
\usepackage{slashed}
\usepackage{graphicx}
\usepackage{xurl}
\usepackage[protrusion=true,expansion=false]{microtype}
\usepackage[hidelinks,pdfpagelayout=TwoPageRight,pdftitle={Theory of Informational Relations: Foundations, Emergent Geometry, and Phenomenological Tests — Version 12.4},pdfauthor={Adrian Lipa},pdfsubject={Audited repository synchronization of primitive informational relations, emergent geometry, flavour structure, particle and gauge sectors, and falsification tests},pdfkeywords={Theory of Informational Relations, TIR, information geometry, emergent geometry, geometric phase, flavour mixing, Standard Model phenomenology, falsification}]{hyperref}
\pdfstringdefDisableCommands{%
  \def\({}%
  \def\){}%
  \def\pi{pi}%
}
\usepackage{makeidx}
\usepackage{verbatim}
\usepackage{longtable}
\usepackage{booktabs}
\usepackage[left=2.5cm,right=2.5cm,top=3cm,bottom=3cm,headheight=14pt]{geometry}
\setlength{\emergencystretch}{3em}
\makeindex
\theoremstyle{definition}
\newtheorem{definition}{Definition}[chapter]
\newtheorem{theorem}{Theorem}[chapter]
\newtheorem{lemma}{Lemma}[chapter]
\newcommand{\claimstatus}[1]{\par\noindent\textbf{Claim status: #1}\par}
\newcommand{\datasetstatus}[1]{\par\noindent\textbf{Data status: #1}\par}
\newcommand{\vTwelveStatus}[3]{\par\noindent\textbf{Claim Class: #1 \quad Timing: #2 \quad Verdict: #3}\par}
\newcommand{\OI}{\kappa}
\newcommand{\Lam}{\Lambda}
\newcommand{\dd}{\,\mathrm{d}}
\newcommand{\mev}{\,\text{MeV}}
\newcommand{\gev}{\,\text{GeV}}
\newcommand{\eV}{\,\text{eV}}
\newcommand{\PDG}[1]{\ensuremath{#1^{\text{PDG}}}}
\newcommand{\pdgsup}{\ensuremath{^{\text{PDG}}}}
\newcommand{\Eproton}{E_{\text{proton}}}
\newcommand{\Eplanck}{E_{P}}
\newcommand{\keV}{\,\text{keV}}
\newcommand{\tev}{\,\text{TeV}}
\newcommand{\SM}{\text{SM}}

\begin{document}
\raggedbottom
\frontmatter
\title{Theory of Informational Relations\\\large Foundations, Emergent Geometry, and Phenomenological Tests}
\author{Adrian Lipa}
\date{Independent Researcher, Doncaster, UK\\Version 12.4 Relational Zero and Graph Reconciliation -- 28 September 2026}
\maketitle
% GENERATED LOCAL-BUILD FLAT PART — DO NOT MERGE

% ---- BEGIN TIR/monograph/frontmatter/publication_frontmatter_v12_0.tex ----
\chapter*{Abstract}
\addcontentsline{toc}{chapter}{Abstract}

This monograph presents the Theory of Informational Relations (TIR) in dependency order: primitive informational relations, emergent geometry, information/phase/flavour structure, particle and gauge sectors, and finally extensions, evidence and falsification gates. Version 12.4 adds the relational-zero foundation and a current-graph reconciliation through 28 September 2026. It preserves v12.0--v12.3 as historical provenance while separating the pre-object relational zero from the later point support and refreshing the cross-repository dependency boundary against current TIR, IDT, RFC, and Secret-of-a-Half surfaces.

The foundational spine begins with the relational zero \(\mathfrak Z_{\rm rel}=(\varnothing,\varnothing)\), proceeds to the minimal nonempty point support, then to the minimal binary distinction and its exchange-fixed half coordinate, passes through binary information and the complex two-state carrier, and reaches the real traceless-Hermitian relational generator space used by the spatial construction. The geometric branch then develops Euclidean structure, the minimal tetrahedral cell, connection/holonomy transport, the connection-lifted $SE(3)$ surface and discrete solder/torsion source. The current repository further supplies conditional Cartan refinement, endpoint-compatible zero-torsion and Levi-Civita selection, a leading-loop second-metric-jet selection rule, and an executable global three-manifold certifier. The production global relational-complex input and production inter-leaf matching field remain explicit open inputs.

The flavour branch introduces the three-flavour carrier before the normalization it supports. The canonical TIR normalization
\[
\kappa=\frac{\ln2}{24\pi}
\]
is owned by the dedicated flavour-mixing normalization theorem and is referenced downstream rather than independently re-derived in multiple publication locations. The discrete constants $(L_3,L_4,L_5)=(7,2,5)$ now additionally have an exact TIR-internal Platonic finite-group closure conditional on the explicit tetrahedral-root extension rule; the older Collatz/twin-prime path is retained as an independent arithmetic crosscheck. This structural closure does not by itself establish a physical law.

The Pass-4 flavour update closes the Jarlskog invariant exactly downstream of the retained CKM matrix inputs, while leaving the first-principles forcing of the CKM mixing weights conditional/open. The golden-mean primitive-orbit asymptotic and XF-9 closure-force/Gram-kernel crosswalk are also synchronized with explicit arithmetic-prime and RH firewalls. Phenomenological sectors retain their source provenance, empirical residuals and frozen failures. The publication status of each promoted result is represented by three independent axes: Claim Class, Timing and Verdict. Repository theorem or software PASS states do not silently promote physical observables.

The complete source, theorem surfaces, validation receipts and version history are maintained in the accompanying repository.

\section*{Keywords}
Theory of Informational Relations; information geometry; emergent geometry; geometric phase; flavour mixing; $SU(3)$; tetrahedral geometry; holonomy; Cartan geometry; Standard Model observables; falsification; reproducibility.

\section*{Claim classes}
\begin{description}
  \item[A -- established result:] a standard mathematical identity or external experimental result supported by the cited literature.
  \item[B -- TIR structural layer:] a TIR structural law, model identification, or internally derived structural quantity with explicit upstream dependencies.
  \item[C -- retrospective assignment:] a sector formula assessed after some or all of the comparison data were known.
  \item[D -- diagnostic or no-go:] an internal audit, retained failure, falsification witness, or restricted impossibility result.
  \item[E -- prospective prediction:] a frozen formula, observable and decision rule specified before the named future dataset is inspected.
  \item[F -- external anchor:] a measured scale, convention or conversion input supplied to a construction.
\end{description}

\section*{Timing and verdict axes}
Timing is recorded as \texttt{RETROSPECTIVE}, \texttt{PROSPECTIVE}, \texttt{EXTERNAL}, or an empty axis where timing is inapplicable. Verdict is one of \texttt{PASS}, \texttt{COMPATIBLE}, \texttt{TENSION}, \texttt{FAIL}, \texttt{OPEN}, or \texttt{QUARANTINED}. The canonical normalization map is maintained in \path{TIR/monograph/v12/STATUS_TAXONOMY.md}.

\section*{Version status}
Version 12.0 remains the dependency-ordered content-migration baseline, Version 12.1 the audited repository-synchronization layer, Version 12.2 the GREMLIN Pass-4 branch consolidation, and Version 12.3 the 18 September semantic-frontier reconciliation. Version 12.4 corrects the primitive zero semantics and refreshes the publication/graph crosswalk. Production/physical gates remain separated from formal/derivational status.

\section*{Current open-boundary summary}
The true current frontier is narrower than the v12.1 headline list. Local Einstein/ADM constraints, evolution and local Einstein-form closure are closed on the declared RFC/TIR selection rules; hypercharge relative uniqueness is closed on the declared field content and TIR anchor; the neutrino source conflict is repaired; and coefficient magnitude evaluation is exact once a parent packet is selected. Remaining gates concern source-owned production spatial/matching/shared-atlas realization and coverage, the coefficient transition-parent selector, continuum gauge normalization/running, common electroweak transport, Higgs action binding, a scheme/scale quark-mass map, meson and strong-CP source closures, cosmological scale/critical-density binding, physical Collatz--FS metric/source coupling, and native Li/Weil/global critical-axis positivity. The Riemann hypothesis remains open.

\section*{Declarations}
This manuscript is an independent research work. The repository records source provenance, validation history and publication dependencies. A successful software or document build certifies the stated source and validation contracts only; it does not by itself provide external experimental validation.

\cleardoublepage

% ---- END TIR/monograph/frontmatter/publication_frontmatter_v12_0.tex ----


\tableofcontents
\mainmatter

\part{Primitive Informational Relations}
% GENERATED LOCAL-BUILD FLAT PART — DO NOT MERGE

% ---- BEGIN TIR/monograph/v12/chapters/ch01_first_distinction_relational_kernel.tex ----
% !TEX root = ../../tir_monograph_v12.tex
\chapter{First Distinction and the Relational Kernel}
\label{ch:v12_first_distinction}

\section{Dependency before temporal order}

The monograph begins before any object is admitted. Write
\begin{equation}
\boxed{X\prec_{\rm dep}Y}
\end{equation}
when the admitted structure represented by $Y$ uses $X$ as an upstream parent in the derivation graph. Temporal ordering receives its own dynamical treatment in the companion time branch; the primitive relation used here is therefore a dependency order.

TIR defines existential admissibility relationally. A relational presentation is
\begin{equation}
\mathfrak R=(X,\mathcal D),
\end{equation}
with support set $X$ and realized distinction/relation set $\mathcal D$. The pre-object zero layer is
\begin{equation}
\boxed{
\mathfrak Z_{\rm rel}=(\varnothing,\varnothing).
}
\label{eq:v12_relational_zero}
\end{equation}
It is not a point and not a singleton. It is zero realized relational structure. Therefore the support set contains no element.

The logical status is exact but scoped. Once TIR defines existence through realized relation or distinction, the absence of an existent object in the empty relational presentation follows by definition. The relational-existence rule itself is the declared TIR ontological entry condition.

Leaving relational zero requires nonempty support. The least nonempty support candidate is the point
\begin{equation}
\boxed{
\mathcal P=\{p\},
\qquad
|\mathcal P|=1.
}
\label{eq:v12_minimal_point_support}
\end{equation}
Thus
\begin{equation}
\boxed{
\mathfrak Z_{\rm rel}\neq\mathcal P.
}
\end{equation}
The unresolved point $(\mathcal P,\varnothing)$ supplies support but not yet a realized relational distinction.
\vTwelveStatus{B}{--}{PASS}

\section{Primitive kernel}

The A0 relational-existence root and the eight A1--A8 TIR axioms are grouped by role rather than by historical order:
\begin{description}
\item[A0 -- relational existence / zero:] defines the empty pre-object relational presentation $\mathfrak Z_{\rm rel}$ and separates zero from the later point support.
\item[A1 -- point minimality:] supplies the minimal non-empty support candidate $\mathcal P$.
\item[A2 -- quantum point:] supplies a normalized Hilbert-space representation for the minimal carrier.
\item[A3 -- information primacy:] supplies distinguishable informational relations and their normalized information carrier.
\item[A4 -- spherical geometric efficiency:] selects the sphere for the declared isotropic boundary/volume functional.
\item[A5 -- arithmetic measures geometry:] reads geometric closure through discrete invariants such as winding or multiplicity.
\item[A6 -- natural numbers from complex phase closure:] assigns natural closure indices to admitted complex phase structure.
\item[A7 -- universal symmetry:] supplies symmetry-governed relational laws and the primitive exchange action.
\item[A8 -- paradox stabilization:] lifts context-dependent incompatible projections into a higher-order closure carrier.
\end{description}

This partition creates a circularity firewall: minimal nontrivial distinction fixes binarity before exchange symmetry is introduced; A7 then acts on the already-binary relation. Arithmetic closure indices enter later through A5--A6.

\section{Minimal first distinction and primitive exchange orbit}

Let $\Omega_{\mathcal P}$ be the state/aspect domain carried by the point support. A distinction is a partition $\Pi$ of this domain. Nontriviality requires more than one nonempty block,
\begin{equation}
|\Pi|>1.
\end{equation}
Hence the least nontrivial distinction has exactly two outcomes,
\begin{equation}
\boxed{
|\Pi_1|=2,
\qquad
\Pi_1=\{N,S\}.
}
\label{eq:v12_minimal_binary_partition}
\end{equation}

\begin{theorem}[Minimal binary distinction]
Every nontrivial partition has at least two nonempty blocks. Therefore the minimal nontrivial distinction is binary.
\end{theorem}

The labels $N,S$ are relational poles or aspects of the minimal carrier. They are not assumed to be two pre-existing spatial objects.

Only after this binary distinction exists does the primitive exchange symmetry act:
\begin{equation}
\boxed{
J:N\leftrightarrow S,
\qquad
J^2=\mathrm{id}.
}
\label{eq:v12_primitive_involution}
\end{equation}
Choosing $x=N$ gives $Jx=S\neq x$ and
\begin{equation}
\boxed{
\mathcal O_J(x)=\{x,Jx\}=\{N,S\}.
}
\label{eq:v12_primitive_binary_orbit}
\end{equation}

Thus binarity is not inferred from an order-two involution. The involution is the exchange symmetry of the already-minimal binary distinction.
\vTwelveStatus{B}{--}{PASS}

\section{Unique exchange-invariant share}

Attach normalized nonnegative relational weights
\begin{equation}
w_N,w_S\ge0,
\qquad
w_N+w_S=1.
\end{equation}
Exchange invariance gives
\begin{equation}
(w_N,w_S)=(w_S,w_N),
\end{equation}
therefore
\begin{equation}
w_N=w_S.
\end{equation}
Normalization fixes the unique symmetric assignment,
\begin{equation}
\boxed{
(w_N,w_S)=\left(\frac12,\frac12\right).
}
\label{eq:v12_unique_half_share}
\end{equation}

\begin{theorem}[Unique exchange-invariant share]
The unique normalized nonnegative weight assignment invariant under exchange of the two primitive poles is $(1/2,1/2)$.
\end{theorem}

Introducing the scalar coordinate
\[
u=w_S,
\qquad
w_N=1-u,
\]
turns the pole exchange into
\begin{equation}
J(u)=1-u.
\end{equation}
Its unique fixed point is
\begin{equation}
\boxed{\operatorname{Fix}(J)=\left\{\frac12\right\}.}
\label{eq:v12_half_fixed_point}
\end{equation}
Thus the normalized exchange-invariant vector and the half coordinate are two descriptions of the same primitive balance.

\section{First information value}

For the normalized binary relation, use the Shannon carrier
\begin{equation}
H_2(u)=-(1-u)\ln(1-u)-u\ln u.
\end{equation}
At the exchange-fixed share,
\begin{equation}
\boxed{H_2(1/2)=\ln2.}
\label{eq:v12_first_ln2}
\end{equation}
This yields the first exact scalar dependency chain
\begin{equation}
\boxed{
\mathfrak Z_{\rm rel}
\rightarrow
\mathcal P
\rightarrow
\{N,S\}
\xrightarrow{\rm exchange\ invariance}
\frac12
\xrightarrow{H_2}
\ln2.
}
\label{eq:v12_primitive_scalar_chain}
\end{equation}
Its minimal TIR parent set is A0, A1, the minimal-nontrivial-distinction rule, A3 and A7. A2 enters only at the coherent quantum lift.

\section{Quantum lift}

Represent the two poles by an orthonormal basis
\[
|N\rangle,\qquad |S\rangle.
\]
Their minimal complex span is
\begin{equation}
\boxed{\mathcal H_{NS}\cong\mathbb C^2.}
\label{eq:v12_binary_c2}
\end{equation}
The balanced coherent family is
\begin{equation}
\boxed{
|\psi_{1/2}(\varphi)\rangle
=\frac{|N\rangle+e^{i\varphi}|S\rangle}{\sqrt2},
\qquad
\varphi\in\mathbb R/2\pi\mathbb Z.
}
\label{eq:v12_balanced_spinor}
\end{equation}
The relative phase is therefore an additional coherent coordinate above the same scalar half point.

The primitive output packet is
\begin{equation}
\boxed{
\mathfrak Z_{\rm rel}
\prec_{\rm dep}
\mathcal P
\prec_{\rm dep}
\{N,S\}
\prec_{\rm dep}
\frac12
\prec_{\rm dep}
\ln2
\prec_{\rm dep}
\mathbb C^2.
}
\label{eq:v12_primitive_packet}
\end{equation}

\section{Branch separation}

Projectivization of the two-state carrier and the spatial geometry developed in Parts I--II proceed from the common primitive packet. The three-flavour/Standard-Model branch proceeds from the same root after the geometric carrier is established. The sibling time programme receives the primitive packet through its own typed interface.

The dependency architecture is therefore
\[
\boxed{
\mathcal C_0
\prec_{\rm dep}
\left\{
\mathcal G_X^{\rm TIR},
\mathcal B_{\rm SM}^{\rm TIR},
\mathcal B_T^{\rm IDT}
\right\},
}
\]
where $\mathcal C_0$ denotes the admitted primitive information/quantum packet.

The spatial and temporal branches later meet through the declared closure surface
\begin{equation}
\mathfrak S_X^{\rm TIR}\otimes\mathfrak T^{\rm IDT}
\longrightarrow
\mathfrak M_{XT},
\end{equation}
and the matter branch then couples to the resulting spacetime carrier. The formal completion order of these joins is recorded in Chapter~\ref{ch:v12_completion_frontier}.

\section{First-chapter invariant}

The foundational narrative used by the rest of v12 is therefore
\begin{equation}
\boxed{
\text{relational zero}
\to\text{minimal point support}
\to\text{minimal nontrivial distinction}
\to\text{two poles}
\to\text{exchange symmetry}
\to\frac12
\to\ln2
\to\mathbb C^2.
}
\end{equation}
The next chapter resolves the geometry and information content of the half coordinate before any flavour normalization is introduced.

% ---- END TIR/monograph/v12/chapters/ch01_first_distinction_relational_kernel.tex ----
\clearpage

% ---- BEGIN TIR/monograph/v12/chapters/ch02_half_binary_information_ln2.tex ----
% !TEX root = ../../tir_monograph_v12.tex
\chapter[Half Coordinate, Binary Information and ln2]{The Half Coordinate, Binary Information and $\ln2$}
\label{ch:v12_half_ln2}

\section{Binary probability coordinate}

The primitive pair from Chapter~\ref{ch:v12_first_distinction} admits the normalized probability coordinate
\begin{equation}
\mathbf p(u)=(1-u,u),
\qquad
0\le u\le1.
\end{equation}
Pole exchange is
\begin{equation}
J(u)=1-u.
\end{equation}
The unique fixed point is
\begin{equation}
\boxed{u_\star=\frac12.}
\label{eq:v12_half_point}
\end{equation}
At this point
\begin{equation}
p_N=p_S=\frac12,
\qquad
\boxed{H_2=\ln2.}
\label{eq:v12_half_entropy}
\end{equation}
The scalar half and the first symmetric information value are therefore fixed before the flavour and particle branches enter.
\vTwelveStatus{B}{--}{PASS}

\section{Coherent lift of the half point}

A normalized pure state of the two-state carrier can be written, after removal of the global phase, as
\begin{equation}
|\psi(u,\varphi)\rangle
=\sqrt{1-u}\,|N\rangle
+e^{i\varphi}\sqrt{u}\,|S\rangle,
\qquad
\varphi\in\mathbb R/2\pi\mathbb Z.
\label{eq:v12_pure_binary_state}
\end{equation}
Projection to the scalar probability coordinate is
\begin{equation}
\pi:\ |\psi(u,\varphi)\rangle\longmapsto(1-u,u).
\end{equation}
Over the half point,
\begin{equation}
\boxed{
\pi^{-1}\!\left(\frac12,\frac12\right)
=
\left\{
\frac{|N\rangle+e^{i\varphi}|S\rangle}{\sqrt2}
:\varphi\in S^1
\right\}.
}
\label{eq:v12_half_fiber}
\end{equation}
Thus the scalar half coordinate lifts to a one-dimensional relative-phase fiber
\begin{equation}
\boxed{\frac12\longrightarrow U(1)\cong S^1.}
\end{equation}
This fiber carries coherent information that is absent from the probability projection.

\section{Bloch realization of the half fiber}

For Eq.~\eqref{eq:v12_pure_binary_state}, the Bloch vector is
\begin{equation}
\mathbf r(u,\varphi)
=\left(
2\sqrt{u(1-u)}\cos\varphi,
2\sqrt{u(1-u)}\sin\varphi,
1-2u
\right).
\end{equation}
Its norm satisfies
\begin{equation}
|\mathbf r|^2
=4u(1-u)+(1-2u)^2
=1.
\end{equation}
Hence pure states lie on $S^2$. At the half point,
\begin{equation}
\boxed{
\mathbf r\left(\frac12,\varphi\right)
=(\cos\varphi,\sin\varphi,0),
}
\label{eq:v12_half_equator}
\end{equation}
which is the Bloch equator.

The primitive two-pole relation therefore becomes
\begin{equation}
\boxed{
N\text{ pole}
\quad\longleftrightarrow\quad
S^1_{\rm half}\text{ equator}
\quad\longleftrightarrow\quad
S\text{ pole}.
}
\end{equation}
The projective geometry of this carrier is developed systematically in Chapter~\ref{ch:v12_complex_carrier}.

\section{Pole exchange on the phase fiber}

Let the exchange operator act by
\begin{equation}
X|N\rangle=|S\rangle,
\qquad
X|S\rangle=|N\rangle.
\end{equation}
On the balanced family,
\begin{equation}
X|\psi_{1/2}(\varphi)\rangle
\sim
|\psi_{1/2}(-\varphi)\rangle,
\end{equation}
where $\sim$ denotes projective equivalence after removal of a global phase. Hence
\begin{equation}
\boxed{\varphi\mapsto-\varphi\pmod{2\pi}.}
\label{eq:v12_half_phase_exchange}
\end{equation}
The exchange-fixed projective phases obey
\begin{equation}
\varphi=-\varphi\pmod{2\pi},
\end{equation}
so
\begin{equation}
\boxed{\varphi\in\{0,\pi\}.}
\label{eq:v12_half_fixed_phases}
\end{equation}
These are the two exchange eigen-directions on the equator.

\section{Phase closure and arithmetic indices}

For a closed phase loop,
\begin{equation}
\boxed{
n=\frac1{2\pi}\oint d\varphi\in\mathbb Z.
}
\label{eq:v12_winding}
\end{equation}
Equivalently, with $z=e^{i\varphi}$,
\begin{equation}
z^n=1.
\end{equation}
For a primitive $n$-fold closure,
\begin{equation}
\varphi=\frac{2\pi k}{n}.
\end{equation}
This supplies the operational A5--A6 chain
\begin{equation}
\boxed{
S^1
\longleftrightarrow
U(1)
\longrightarrow
\text{winding/closure invariant}
\longrightarrow
n\in\mathbb N.
}
\end{equation}
The ordering is important: the phase geometry is admitted before the arithmetic closure index is used as a TIR physical label.

\section{Information geometry at the half seam}

The binary probability projection sees the Bloch vertical coordinate
\begin{equation}
z_B=1-2u.
\end{equation}
At $u=1/2$,
\begin{equation}
z_B=0.
\end{equation}
All balanced relative phases therefore project to the same scalar probability point,
\begin{equation}
\boxed{
\pi:S^1_{\rm equator}\to\left\{\frac12\right\}.
}
\end{equation}
This is the exact interface between binary information and the coherent phase degree of freedom: the scalar coordinate fixes equal modulus, while the $U(1)$ fiber carries the relative phase.

\section{Dynamical interfaces}

A strongly continuous unitary flow on the half fiber admits a self-adjoint generator under the standard unitary-generator theorem. In the temporal branch, a parametrization
\begin{equation}
\varphi=\varphi(\tau)
\end{equation}
therefore supplies the entry point to phase dynamics. Multiple Bloch distinction axes similarly give noncommuting Pauli generators and the standard Robertson--Heisenberg structure.

These are downstream interfaces of the carrier established here. The present chapter fixes the static packet
\begin{equation}
\boxed{
\frac12
\leftrightarrow
\ln2
\leftrightarrow
S^1_{\rm phase}
\subset
\mathbb{CP}^1\cong S^2.
}
\end{equation}

\section{Ownership boundary}

Chapter 2 owns the binary half, $\ln2$, and its coherent phase fiber. Chapter 8 later introduces the three-flavour carrier, and Chapter 9 is the sole v12 publication owner of the complete normalization
\[
\kappa=\frac{\ln2}{24\pi}.
\]
This dependency order keeps the binary information numerator upstream of the flavour-mixing phase measure that normalizes it.

% ---- END TIR/monograph/v12/chapters/ch02_half_binary_information_ln2.tex ----
\clearpage

% ---- BEGIN TIR/monograph/v12/chapters/ch03_complex_carrier_information_geometry.tex ----
% !TEX root = ../../tir_monograph_v12.tex
\chapter{Complex Relational Carrier and Information Geometry}
\label{ch:v12_complex_carrier}

\section{Two-state carrier}

The primitive quantum lift gives
\begin{equation}
\mathcal H_2\cong\mathbb C^2.
\end{equation}
Choose the pole basis $|N\rangle,|S\rangle$. A normalized pure state can be parameterized by
\begin{equation}
\boxed{
|\psi(\vartheta,\varphi)\rangle
=\cos\frac\vartheta2\,|N\rangle
+e^{i\varphi}\sin\frac\vartheta2\,|S\rangle.
}
\label{eq:v12_spinor_spherical}
\end{equation}
The associated probability coordinate is
\begin{equation}
u=\sin^2\frac\vartheta2,
\qquad
1-u=\cos^2\frac\vartheta2.
\end{equation}
Hence the half seam $u=1/2$ is the equatorial condition
\begin{equation}
\vartheta=\frac\pi2.
\end{equation}

\section{Projective state space}

Physical pure-state rays remove global complex phase, giving
\begin{equation}
\boxed{\mathbb{CP}^1\cong S^2.}
\label{eq:v12_cp1_s2}
\end{equation}
The standard Fubini--Study line element for normalized rays may be written
\begin{equation}
 ds_{\rm FS}^2
=\langle d\psi|d\psi\rangle
-|\langle\psi|d\psi\rangle|^2
\label{eq:v12_fs_general}
\end{equation}
and, for Eq.~\eqref{eq:v12_spinor_spherical}, becomes
\begin{equation}
\boxed{
 ds_{\rm FS}^2
=\frac14\left(d\vartheta^2+\sin^2\vartheta\,d\varphi^2\right).
}
\label{eq:v12_fs_qubit}
\end{equation}
Thus the qubit projective metric is one quarter of the unit Bloch-sphere metric. Its total area is
\begin{equation}
\boxed{\operatorname{Area}_{FS}(\mathbb{CP}^1)=\pi.}
\end{equation}
This normalized projective area later provides an exact geometric measure for tetrahedral crosswalks; physical length/area calibration remains separately typed.

\section{Density-matrix representation}

For a normalized pure state,
\begin{equation}
\rho_\psi=|\psi\rangle\langle\psi|
\end{equation}
can be written in Bloch form
\begin{equation}
\boxed{
\rho_\psi
=\frac12\left(I+\mathbf r\cdot\boldsymbol\sigma\right),
\qquad
|\mathbf r|=1.
}
\label{eq:v12_bloch_density}
\end{equation}
The coefficient vector is
\begin{equation}
\mathbf r
=(\sin\vartheta\cos\varphi,
\sin\vartheta\sin\varphi,
\cos\vartheta).
\end{equation}
For a general normalized qubit state, $|\mathbf r|\le1$ and the trace-one Hermitian state space has affine hull
\begin{equation}
\boxed{
\mathcal A_2
=\{\rho=\rho^\dagger:\operatorname{Tr}\rho=1\}
=\frac12I+\operatorname{Herm}_0(2).
}
\label{eq:v12_affine_qubit_hull}
\end{equation}

\section{Information coordinate inside projective geometry}

The binary probability coordinate is related to the Bloch vertical coordinate by
\begin{equation}
r_z=1-2u=\cos\vartheta.
\end{equation}
The exchange-fixed information point therefore obeys
\begin{equation}
u=\frac12
\iff
r_z=0
\iff
\vartheta=\frac\pi2.
\end{equation}
The phase coordinate $\varphi$ remains free along the equator. This makes the information/projective relation explicit:
\begin{equation}
\boxed{
\text{binary share }u
\longleftrightarrow
\text{Bloch latitude},
\qquad
\text{relative phase }\varphi
\longleftrightarrow
\text{Bloch azimuth}.
}
\end{equation}

At the half seam the entropy is maximal,
\begin{equation}
H_2(1/2)=\ln2,
\end{equation}
while the coherent family still carries a full $S^1$ phase orbit. The scalar information coordinate and the projective phase coordinate are therefore complementary coordinates of the same two-state carrier.

\section{Operator coordinates}

The real vector space of Hermitian operators on $\mathbb C^2$ admits the unique Pauli decomposition
\begin{equation}
\boxed{
A=a_0I+a_x\sigma_x+a_y\sigma_y+a_z\sigma_z,
\qquad a_\mu\in\mathbb R.
}
\label{eq:v12_pauli_decomposition}
\end{equation}
with
\begin{equation}
\dim_{\mathbb R}\operatorname{Herm}(2)=4.
\end{equation}
The trace component is the uniform direction; the distinction-generating part lies in
\begin{equation}
\operatorname{Herm}_0(2)
=\operatorname{span}_{\mathbb R}\{\sigma_x,\sigma_y,\sigma_z\}.
\end{equation}
Chapter~\ref{ch:v12_c2_herm0_r3} turns this observation into the explicit three-real-dimensional relational generator carrier and its affine spatial relation object.

\section{Rotation covariance}

For $U\in SU(2)$,
\begin{equation}
A\mapsto UAU^\dagger
\end{equation}
preserves trace and the Hilbert--Schmidt pairing. On the Pauli coefficient vector the standard adjoint action induces
\begin{equation}
\boxed{
\operatorname{Ad}:SU(2)\to SO(3),
\qquad
SU(2)/\{\pm I\}\cong SO(3).
}
\label{eq:v12_su2_so3}
\end{equation}
The Bloch sphere, traceless operator sphere and rotation group are therefore coordinated by one algebra rather than introduced as unrelated copies of $S^2$ and $\mathbb R^3$.

\section{Fubini--Study tetrahedral preview}

Four unit Bloch vectors satisfying
\begin{equation}
\mathbf n_a\cdot\mathbf n_b=-\frac13
\qquad(a\neq b)
\end{equation}
form a regular tetrahedral frame. Their Bloch central angle obeys
\begin{equation}
\cos\chi=-\frac13.
\end{equation}
For the corresponding geodesic $\mathbb{CP}^1$ face, the spherical interior angle is $2\pi/3$ and the Fubini--Study area is
\begin{equation}
\boxed{a_{FS}^{\rm face}=\frac\pi4.}
\end{equation}
Four faces therefore cover total Fubini--Study area
\begin{equation}
\boxed{a_{FS}^{\rm tet}=\pi,}
\end{equation}
matching the full normalized $\mathbb{CP}^1$ area. Chapter~\ref{ch:v12_tetrahedral_closure} returns to this frame after the spatial carrier and Euclidean dimension gate have been established.

\section{Dependency output}

The chapter exports the exact mathematical packet
\begin{equation}
\boxed{
\mathbb C^2
\to
\mathbb{CP}^1\cong S^2
\to
\rho=\frac12(I+\mathbf r\cdot\boldsymbol\sigma)
\to
\operatorname{Herm}(2)
\supset
\operatorname{Herm}_0(2).
}
\end{equation}
The next chapter isolates the traceless translation/generator space and constructs the canonical affine relation
\[
\mathcal E_{xy}=2(\rho_y-\rho_x),
\]
which is the bridge from projective quantum-state geometry to the local three-real-dimensional relational carrier.

% ---- END TIR/monograph/v12/chapters/ch03_complex_carrier_information_geometry.tex ----
\clearpage




\part{Emergence of Geometry}
% GENERATED LOCAL-BUILD FLAT PART — DO NOT MERGE

% ---- BEGIN TIR/monograph/v12/chapters/ch04_c2_to_herm0_r3.tex ----
% !TEX root = ../../tir_monograph_v12.tex
\chapter[C2 to Herm0(2) and R3]{From $\mathbb C^2$ to $\operatorname{Herm}_0(2)\cong\mathbb R^3$}
\label{ch:v12_c2_herm0_r3}

\section{Traceless relational generator space}

Every Hermitian operator on the binary carrier has the Pauli decomposition
\begin{equation}
A=a_0I+\mathbf a\cdot\boldsymbol\sigma,
\qquad
\mathbf a=(a_x,a_y,a_z)\in\mathbb R^3.
\end{equation}
The uniform component $a_0I$ acts equally on the two primitive alternatives. The distinction-generating translation space is therefore
\begin{equation}
\boxed{
\mathfrak g_{\rm rel}
:=\operatorname{Herm}_0(2)
=\{A=A^\dagger:\operatorname{Tr}A=0\}
=\operatorname{span}_{\mathbb R}\{\sigma_x,\sigma_y,\sigma_z\}.
}
\label{eq:v12_herm0_def}
\end{equation}
Hence
\begin{equation}
\boxed{
\dim_{\mathbb R}\mathfrak g_{\rm rel}=3,
\qquad
\operatorname{Herm}_0(2)\cong\mathbb R^3.
}
\label{eq:v12_herm0_r3}
\end{equation}
\vTwelveStatus{A}{--}{PASS}
This is an exact linear-algebraic property of the admitted two-state complex carrier.

\section{Canonical metric}

Define the Hilbert--Schmidt pairing on traceless Hermitian generators by
\begin{equation}
\boxed{
\langle A,B\rangle_{\rm HS}
=\frac12\operatorname{Tr}(AB).
}
\label{eq:v12_hs_metric}
\end{equation}
The Pauli basis obeys
\begin{equation}
\frac12\operatorname{Tr}(\sigma_i\sigma_j)=\delta_{ij}.
\end{equation}
Thus, for
\[
A=\mathbf a\cdot\boldsymbol\sigma,
\qquad
B=\mathbf b\cdot\boldsymbol\sigma,
\]
one obtains
\begin{equation}
\boxed{
\langle A,B\rangle_{\rm HS}=\mathbf a\cdot\mathbf b,
\qquad
\|A\|^2=\frac12\operatorname{Tr}(A^2)=|\mathbf a|^2.
}
\label{eq:v12_hs_euclidean}
\end{equation}
The generator space therefore carries a positive-definite Euclidean quadratic form canonically inherited from the operator algebra.
\vTwelveStatus{A}{--}{PASS}

\section{Unit relational directions}

For
\begin{equation}
A=\mathbf n\cdot\boldsymbol\sigma,
\qquad
|\mathbf n|=1,
\end{equation}
the Pauli algebra gives
\begin{equation}
A^2=I.
\end{equation}
Conversely, a traceless Hermitian $2\times2$ operator with $A^2=I$ has unit coefficient vector. Hence
\begin{equation}
\boxed{
\{A\in\mathfrak g_{\rm rel}:\|A\|=1\}
\cong S^2.
}
\label{eq:v12_generator_sphere}
\end{equation}
The Bloch pure-state sphere and the unit generator sphere therefore live in the same coefficient carrier.

\section[SU(2) covariance]{$SU(2)$ covariance}

Conjugation by $U\in SU(2)$ acts as
\begin{equation}
A\mapsto UAU^\dagger.
\end{equation}
It preserves tracelessness, Eq.~\eqref{eq:v12_hs_metric}, and the norm. On coefficient vectors this is the adjoint rotation
\begin{equation}
\boxed{
\operatorname{Ad}:SU(2)\to SO(3),
\qquad
\ker\operatorname{Ad}=\{\pm I\}.
}
\end{equation}
Thus the same algebra supplies the three-real-dimensional carrier, its Euclidean metric and its rotation group.

\section{Trace-one affine state space}

Let a primitive locus $x$ carry a normalized two-level state
\begin{equation}
\rho_x=\rho_x^\dagger,
\qquad
\operatorname{Tr}\rho_x=1.
\end{equation}
The affine hull of normalized states is
\begin{equation}
\boxed{
\mathcal A_2
=\frac12I+\operatorname{Herm}_0(2),
\qquad
\dim_{\mathbb R}\mathcal A_2=3.
}
\label{eq:v12_affine_state_space}
\end{equation}
Its translation space is exactly $\operatorname{Herm}_0(2)$.

\section{Canonical affine relation}

For two primitive loci $x,y$, define
\begin{equation}
\boxed{
\mathcal E_{xy}:=2(\rho_y-\rho_x).
}
\label{eq:v12_affine_relation}
\end{equation}
If
\[
\rho_x=\frac12(I+\mathbf r_x\cdot\boldsymbol\sigma),
\qquad
\rho_y=\frac12(I+\mathbf r_y\cdot\boldsymbol\sigma),
\]
then
\begin{equation}
\boxed{
\mathcal E_{xy}
=(\mathbf r_y-\mathbf r_x)\cdot\boldsymbol\sigma
\in\operatorname{Herm}_0(2).
}
\end{equation}
The normalization factor $2$ removes the Bloch $1/2$ convention and gives the direct coefficient displacement.

The affine relation satisfies the exact endpoint laws
\begin{align}
\mathcal E_{xx}&=0,\\
\mathcal E_{yx}&=-\mathcal E_{xy},\\
\boxed{\mathcal E_{xz}}&=\boxed{\mathcal E_{xy}+\mathcal E_{yz}}.
\label{eq:v12_endpoint_composition}
\end{align}
Equivalently,
\begin{equation}
\mathcal E_{xy}+\mathcal E_{yz}+\mathcal E_{zx}=0.
\end{equation}
The local endpoint-composition rule is therefore inherited directly from affine state differences.
\vTwelveStatus{A}{--}{PASS}

\section{Metric of relational displacement}

Using Eq.~\eqref{eq:v12_hs_metric},
\begin{equation}
\boxed{
\|\mathcal E_{xy}\|^2
=\frac12\operatorname{Tr}(\mathcal E_{xy}^2)
=|\mathbf r_y-\mathbf r_x|^2.
}
\label{eq:v12_affine_relation_metric}
\end{equation}
A common local unitary frame transformation
\begin{equation}
\rho_x\mapsto U\rho_xU^\dagger
\end{equation}
gives
\begin{equation}
\mathcal E_{xy}\mapsto U\mathcal E_{xy}U^\dagger,
\end{equation}
so the relation, metric and composition law are all $SU(2)$-covariant with effective $SO(3)$ action on coefficient vectors.

\section{Spatial promotion gate}

The TIR spatial bridge is now sharply typed:
\begin{equation}
\boxed{
\texttt{RELATION\_AS\_CANONICAL\_QUANTUM\_STATE\_DIFFERENCE}:
\qquad
x\to y
\ \widehat{=}\ 
\mathcal E_{xy}=2(\rho_y-\rho_x).
}
\label{eq:v12_relation_promotion_rule}
\end{equation}
Under this identification, the local physical relation carrier is
\begin{equation}
\boxed{
T_x\Sigma\;\widehat{=}\;\operatorname{Herm}_0(2)\cong\mathbb R^3
}
\label{eq:v12_spatial_promotion}
\end{equation}
with metric inherited from Eq.~\eqref{eq:v12_hs_metric} and additive local displacement inherited from Eq.~\eqref{eq:v12_endpoint_composition}.
\vTwelveStatus{B}{--}{OPEN}
The exact algebraic consequences of the rule are established above; the foundational promotion of physical spatial relation to the canonical quantum-state difference remains the named bridge whose uniqueness is addressed by the next chapter.

\section{Tetrahedral affine frame preview}

Four affinely independent pure qubit states span the full trace-one affine hull. The regular tetrahedral set
\begin{equation}
\sum_{a=1}^4\mathbf n_a=0,
\qquad
\mathbf n_a\cdot\mathbf n_b=-\frac13\quad(a\neq b)
\end{equation}
provides a minimal four-point affine frame in this three-dimensional carrier. It later appears simultaneously as the regular spatial cell and the qubit informationally complete tetrahedral frame.

The output of Part I is therefore
\begin{equation}
\boxed{
\mathbb C^2
\to
\mathcal A_2=\frac12I+\operatorname{Herm}_0(2)
\to
\mathcal E_{xy}\in\operatorname{Herm}_0(2)\cong\mathbb R^3
}
\end{equation}
with exact affine composition, Euclidean metric and rotation covariance. Part II now asks when this carrier is uniquely promoted to local physical space and how its minimal finite cell closes.

% ---- END TIR/monograph/v12/chapters/ch04_c2_to_herm0_r3.tex ----
\clearpage

% ---- BEGIN TIR/monograph/v12/chapters/ch05_euclidean_geometry_dimension_gate.tex ----
% !TEX root = ../../tir_monograph_v12.tex
\chapter{Euclidean Geometry and the Spatial Dimension Gate}
\label{ch:v12_euclidean_dimension}

\section{Canonical affine relation carrier}

The normalized two-level state carrier has affine hull
\[
\mathcal A_2
=
\{\rho=\rho^\dagger:\operatorname{Tr}\rho=1\}
=
\frac12I+\operatorname{Herm}_0(2).
\]
Write
\[
V:=\operatorname{Herm}_0(2)
=
\operatorname{span}_{\mathbb R}\{\sigma_x,\sigma_y,\sigma_z\}.
\]
Since \(\dim_{\mathbb R}\operatorname{Herm}(2)=4\) and the trace-one condition removes one
real affine degree of freedom,
\[
\boxed{\dim_{\mathbb R}V=3.}
\]

For an ordered pair of admitted point states \((\rho_x,\rho_y)\), the affine torsor
has a unique endpoint-carrying displacement
\[
\boxed{\mathfrak e_{xy}:=\rho_y-\rho_x\in V},
\qquad
\rho_x+\mathfrak e_{xy}=\rho_y.
\]
The generator-normalized relation used below is
\[
\boxed{\mathcal E_{xy}:=2\mathfrak e_{xy}.}
\]
It satisfies the exact endpoint identities
\[
\mathcal E_{yx}=-\mathcal E_{xy},
\qquad
\boxed{\mathcal E_{xz}=\mathcal E_{xy}+\mathcal E_{yz}},
\]
and therefore
\[
\mathcal E_{xy}+\mathcal E_{yz}+\mathcal E_{zx}=0.
\]
This endpoint-carrying universal property is the canonical local relation exported
by the v12 spatial branch.

\section{Invariant Euclidean metric}

On \(V\), define
\[
\boxed{
\langle A,B\rangle
=
\frac12\operatorname{Tr}(AB).
}
\]
The Pauli basis obeys
\[
\frac12\operatorname{Tr}(\sigma_i\sigma_j)=\delta_{ij}.
\]
Hence, for
\[
A=\mathbf a\cdot\boldsymbol\sigma,
\qquad
B=\mathbf b\cdot\boldsymbol\sigma,
\]
one has
\[
\boxed{
\langle A,B\rangle=\mathbf a\cdot\mathbf b,
\qquad
\|A\|^2=\frac12\operatorname{Tr}(A^2)=|\mathbf a|^2.
}
\]
The positive quadratic form is therefore Euclidean in Pauli coefficients.

The Pythagorean relation follows immediately. If
\[
\langle A,B\rangle=0,
\]
then
\[
\begin{aligned}
\|A+B\|^2
&=\langle A+B,A+B\rangle\\
&=\|A\|^2+2\langle A,B\rangle+\|B\|^2,
\end{aligned}
\]
so
\[
\boxed{\|A+B\|^2=\|A\|^2+\|B\|^2.}
\]
Thus the local affine relation carrier already contains the Euclidean norm,
orthogonality and Pythagorean closure used by the later spatial construction.

\section{Rotational covariance}

For \(U\in SU(2)\),
\[
A\mapsto UAU^\dagger
\]
preserves trace and the Hilbert--Schmidt metric.  The induced action on coefficient
vectors is
\[
\boxed{
\operatorname{Ad}:SU(2)\longrightarrow SO(3),
\qquad
SU(2)/\{\pm I\}\cong SO(3).
}
\]
The unit relation directions therefore form
\[
\boxed{
\{A\in V:\|A\|=1\}\cong S^2.
}
\]

\section{Rank-three stabilization}

At a relational locus \(x\), define the outgoing relation span
\[
W_x
:=
\operatorname{span}_{\mathbb R}
\{\mathcal E_{xy}:y\sim x\}
\subseteq V.
\]
Assume a populated primitive patch,
\[
W_x\neq\{0\},
\]
and full local relational isotropy,
\[
R(W_x)=W_x
\qquad
\text{for every }R\in SO(3)
\]
from the adjoint action above.

The defining real representation of \(SO(3)\) on \(\mathbb R^3\) is irreducible.
Its invariant linear subspaces are therefore \(\{0\}\) and the full carrier.
Consequently
\[
\boxed{
W_x\neq0,\quad SO(3)W_x=W_x
\Longrightarrow
W_x=V
\Longrightarrow
\dim W_x=3.
}
\]
Rank one selects an axis and rank two selects a plane; the full primitive isotropy
stabilizes the three-dimensional carrier.

\section{Minimal faithful realization crosscheck}

The same dimension is independently fixed by the minimal faithful real orthogonal
realization of the primitive rotation group. A continuous \(SO(3)\) action on a
one-dimensional real orthogonal carrier has trivial connected image.  A
two-dimensional orthogonal image lies in \(SO(2)\), while \(SO(3)\) is non-Abelian
and perfect.  The defining three-dimensional representation is faithful.

Therefore the minimal faithful real orthogonal carrier has
\[
\boxed{\dim_{\mathbb R}V_x=3}
\]
and is unique up to orthogonal equivalence.  Its invariant positive metric is
unique up to an overall positive scale.  The canonical dimensionless
normalization used here is exactly
\[
h(A,B)=\frac12\operatorname{Tr}(AB).
\]

\section{Continuum interface}

The discrete spatial branch is now source-bound by
\[
\boxed{
(\rho_x,\rho_y)
\longmapsto
\mathcal E_{xy}=2(\rho_y-\rho_x)
\in\operatorname{Herm}_0(2)\cong\mathbb R^3.
}
\]
The remaining geometric completion gate is the regular gluing/refinement step
that turns neighboring local copies of \(V\) into a smooth coframe/tangent-bundle
carrier.  Chapter~\ref{ch:v12_connections_holonomy} supplies the discrete
connection, solder and torsion objects required for that refinement.

\section{Source and validation receipt}

The canonical theorem surfaces are
\begin{itemize}
\item \path{TIR/subrepos/the-space-of-geometry/foundations/CANONICAL_SPATIAL_RELATION_EXTRACTION_V0_1.md};
\item \path{TIR/foundations/TIR_RANK3_ISOTROPY_STABILIZATION_V0_1.md};
\item \path{TIR/subrepos/the-space-of-geometry/foundations/SPATIAL_PROMOTION_UNIQUENESS_V0_1.md}.
\end{itemize}

The deterministic v12 receipt is generated by
\path{TIR/validation/tir_v12_ch05_euclidean_spatial_gate_v0_1.py}.

% ---- END TIR/monograph/v12/chapters/ch05_euclidean_geometry_dimension_gate.tex ----
\clearpage

% ---- BEGIN TIR/monograph/v12/chapters/ch06_tetrahedral_closure.tex ----
% !TEX root = ../../tir_monograph_v12.tex
\chapter{Minimal Finite Cell and Tetrahedral Closure}
\label{ch:v12_tetrahedral_closure}

\section{Minimal finite isotropy}

Let
\[
\mathbf n_1,\ldots,\mathbf n_m\in\mathbb R^3,
\qquad
|\mathbf n_a|=1,
\]
be equal-weight local relation directions in the carrier constructed in
Chapter~\ref{ch:v12_euclidean_dimension}.  Define
\[
\mathbf M:=\sum_{a=1}^{m}\mathbf n_a,
\qquad
Q:=\sum_{a=1}^{m}\mathbf n_a\mathbf n_a^T.
\]
The finite isotropy conditions are
\[
\boxed{\mathbf M=0},
\qquad
\boxed{Q=\frac m3 I_3}.
\]
The second condition has rank three.  For \(m\le2\) the relation span is
rank-deficient.  If \(m=3\), the zero-moment equation places the third vector in
the span of the first two, so the rank is at most two.  Hence
\[
\boxed{m\ge4.}
\]

For \(m=4\), place the four vectors into
\[
N=(\mathbf n_1\ \mathbf n_2\ \mathbf n_3\ \mathbf n_4).
\]
Then
\[
NN^T=\frac43I_3,
\qquad
N\mathbf1=0.
\]
The Gram matrix \(G=N^TN\) is therefore
\[
\boxed{
G
=
\frac43I_4-\frac13\mathbf1\mathbf1^T,
}
\]
so
\[
\boxed{
\mathbf n_a\cdot\mathbf n_b
=
-\frac13
\quad(a\neq b).
}
\]
The minimal equal-weight finite full-isotropy cell is consequently the regular
tetrahedral frame.

A convenient exact representative is
\[
\mathbf n_1=\frac1{\sqrt3}(1,1,1),\quad
\mathbf n_2=\frac1{\sqrt3}(1,-1,-1),
\]
\[
\mathbf n_3=\frac1{\sqrt3}(-1,1,-1),\quad
\mathbf n_4=\frac1{\sqrt3}(-1,-1,1).
\]

\section{Exact shape data}

For distinct vertices,
\[
\|\mathbf n_a-\mathbf n_b\|^2
=
2-2\left(-\frac13\right)
=
\frac83,
\]
hence the normalized edge length is
\[
\boxed{\hat a=\sqrt{\frac83}}.
\]
The Euclidean volume of the regular tetrahedron is
\[
\boxed{
\hat V_{\Delta^3}
=
\frac{\hat a^3}{6\sqrt2}
=
\frac{8}{9\sqrt3}.
}
\]

\section{Independent qubit informational-completeness route}

A generic qubit density operator has three independent real Bloch coordinates.
An \(m\)-outcome normalized probability vector has at most \(m-1\) independent
real entries.  Informational completeness therefore requires
\[
m-1\ge3,
\qquad
\boxed{m\ge4}.
\]

Using the same tetrahedral directions, define
\[
P_a=\frac12(I+\mathbf n_a\cdot\boldsymbol\sigma),
\qquad
E_a=\frac12P_a
=\frac14(I+\mathbf n_a\cdot\boldsymbol\sigma).
\]
Since \(\sum_a\mathbf n_a=0\),
\[
\boxed{\sum_{a=1}^{4}E_a=I}.
\]
For \(a\neq b\),
\[
\boxed{\operatorname{Tr}(P_aP_b)=\frac13}.
\]
For
\[
\rho=\frac12(I+\mathbf r\cdot\boldsymbol\sigma)
\]
the outcome probabilities are
\[
p_a=\operatorname{Tr}(\rho E_a)
=\frac14(1+\mathbf r\cdot\mathbf n_a).
\]
Using
\[
\sum_a\mathbf n_a\mathbf n_a^T=\frac43I_3
\]
gives the exact reconstruction
\[
\boxed{
\mathbf r
=
3\sum_{a=1}^{4}p_a\mathbf n_a.
}
\]
Thus the minimal symmetric informationally complete qubit frame and the minimal
finite isotropic relation frame carry the same regular tetrahedral Gram data.

\section{Congruence-class closure}

Let \(N\) be an ordered spatial tetrahedral frame and \(M\) an ordered SIC
tetrahedral frame, with
\[
N^TN=M^TM=G,
\qquad
NN^T=MM^T=\frac43I_3.
\]
Define
\[
\boxed{Q_{\Delta}:=\frac34MN^T.}
\]
Then
\[
Q_{\Delta}N=M,
\qquad
Q_{\Delta}Q_{\Delta}^T=I_3,
\]
so
\[
\boxed{Q_{\Delta}\in O(3).}
\]
After compatible orientation of the ordered frames the representative may be
chosen in \(SO(3)\).  The two routes therefore occupy one exact tetrahedral
orthogonal-congruence class in
\(\operatorname{Herm}_0(2)\cong\mathbb R^3\).

The common promoted object at this stage is this Gram/congruence class and its
orthogonal shape invariants.  Semantic role binding between an informational
outcome label and a spatial adjacency label remains a separate downstream
binding.

\section{Fubini--Study shape crosswalk}

The Bloch central angle between tetrahedral directions obeys
\[
\cos\chi=-\frac13.
\]
For one spherical tetrahedral face the spherical cosine law gives the interior
angle
\[
\alpha=\frac{2\pi}{3}.
\]
The spherical excess on the unit Bloch sphere is \(\pi\).  Since the qubit
Fubini--Study metric satisfies
\[
ds_{FS}^2=\frac14ds_{S^2}^2,
\]
one face has
\[
\boxed{a_{FS}^{face}=\frac\pi4},
\]
and all four faces carry
\[
\boxed{a_{FS}^{tet}=\pi}.
\]
Therefore the dimensionless dual-shape coefficient is
\[
\boxed{
C_{\Delta/FS}
=
\frac{\hat V_{\Delta^3}}{a_{FS}^{tet}}
=
\frac{8}{9\sqrt3\,\pi}.
}
\]

\section{Legacy tetrahedral semantics}

The historical v11 tetrahedron appendix carries the sector-specific
operator labels
\(\{\hat I,\hat M,\hat A_\mu,\hat A_\tau\}\) and its own six edge assignments.
In v12 that object remains sector-level provenance.  The canonical geometric
tetrahedron of this chapter is the regular Gram class derived above; any
role-binding map between the two tetrahedral descriptions must preserve their
typed edge data and receives its own source receipt.

\section{Source and validation receipt}

The canonical theorem surfaces are
\begin{itemize}
\item \path{TIR/foundations/TIR_MINIMAL_ISOTROPIC_TETRAHEDRAL_CELL_V0_1.md};
\item \path{TIR/foundations/TIR_QUBIT_TETRAHEDRAL_INFORMATIONAL_COMPLETENESS_V0_1.md};
\item \path{TIR/integration/TIR_TETRAHEDRAL_CONGRUENCE_CLASS_CLOSURE_V0_1.md}.
\end{itemize}
The deterministic v12 receipt is generated by
\path{TIR/validation/tir_v12_ch06_tetrahedral_closure_v0_1.py}.

% ---- END TIR/monograph/v12/chapters/ch06_tetrahedral_closure.tex ----
\clearpage

% ---- BEGIN TIR/monograph/v12/chapters/ch07_connections_holonomy_se3_solder_torsion.tex ----
% !TEX root = ../../tir_monograph_v12.tex
\chapter[Connections, Holonomy, SE(3), Solder and Torsion]{Connections, $W_{ij}$, Holonomy, $SE(3)$, Solder and Torsion}
\label{ch:v12_connections_holonomy}

\paragraph{v12.1 synchronization.}
The original v12 migration closed this chapter at the discrete Universal-Loop
source and treated Cartan refinement as the next theorem.  The current
repository has since closed the local/refining Gates A2--A4 and implemented the
Gate-A5 global three-manifold certifier.  This chapter now records those results
without promoting the still-missing production global relational-complex input.

\section{Typed connection transport}

Let \(G\) be the structure group of an active sector and \(A_R\) its connection
in a representation \(R\).  The common transport object is
\[
\boxed{
W_{ij}^{(G,R)}
=\mathcal P\exp\!\left(\int_{\gamma_{ij}}A_R\right).
}
\]
For unitary transport,
\[
\boxed{W_{ji}=W_{ij}^{-1}=W_{ij}^{\dagger}},
\]
and under local frame changes
\[
\boxed{W_{ij}\mapsto G_iW_{ij}G_j^{-1}}.
\]

The v12 typing is
\[
\boxed{
W_{ij}^{WT}\in U(1),
\qquad W_{ij}^{X}\in SU(2),
\qquad W_{ij}^{c}\in SU(3).
}
\]
For the spatial branch,
\[
\boxed{R_{ij}:=\operatorname{Ad}_{W_{ij}^{X}}\in SO(3)}.
\]
A composable closed path carries
\[
\boxed{W_\gamma=W_{i_0i_1}W_{i_1i_2}\cdots W_{i_{m-1}i_0}}.
\]
Its conjugacy class is the gauge-covariant finite loop datum.

\section{Atlas transitions and the holonomy firewall}

Let a local affine chart \(a\) have anchor \(\mathbf r_a\) and orthonormal
frame \(Q_a\in SO(3)\), with
\[
x_a=Q_a^T(\mathbf r-\mathbf r_a).
\]
For overlapping charts,
\[
\boxed{x_b=Q_b^TQ_a\,x_a+Q_b^T(\mathbf r_a-\mathbf r_b)}.
\]
Therefore
\[
\boxed{
G_{ba}^{atlas}
=(R_{ba},t_{ba})
=\left(Q_b^TQ_a,Q_b^T(\mathbf r_a-\mathbf r_b)\right)
\in SE(3).
}
\]
For three charts,
\[
G_{cb}^{atlas}G_{ba}^{atlas}=G_{ca}^{atlas},
\qquad
G_{ab}^{atlas}=(G_{ba}^{atlas})^{-1}.
\]
A closed sequence of pure atlas transitions therefore gives
\[
\boxed{G_C^{atlas}=e_{SE(3)}}.
\]
This exact cocycle is the zero-holonomy baseline.  Path-dependent connection
transport is carried by the separately typed \(W_{ij}\) layer.

\section{Discrete solder object and endpoint defect}

For every oriented spatial edge define
\[
\mathcal E_{xy}=E^a_{xy}\sigma_a\in\operatorname{Herm}_0(2),
\qquad
\ell_{xy}^2=\frac12\operatorname{Tr}(\mathcal E_{xy}^2).
\]
The reversed edge obeys the covariant orientation rule
\[
\boxed{
\mathcal E_{yx}
=-W_{yx}^{X}\mathcal E_{xy}(W_{yx}^{X})^\dagger.
}
\]
For the oriented triangle \(x\to y\to z\to x\), define
\[
\boxed{
\mathcal T_{xyz}
=\mathcal E_{xy}
+W_{xy}^{X}\mathcal E_{yz}(W_{xy}^{X})^\dagger
+W_{xz}^{X}\mathcal E_{zx}(W_{xz}^{X})^\dagger.
}
\]
The transported endpoint defect is
\[
\boxed{
\mathcal C_{xyz}
=\mathcal E_{xz}
-\left(
\mathcal E_{xy}
+W_{xy}^{X}\mathcal E_{yz}(W_{xy}^{X})^\dagger
\right).
}
\]
The existing relational endpoint identity is
\[
\boxed{\mathcal T_{xyz}=-\mathcal C_{xyz}}.
\]

\section{Gate A: affine-loop torsion source}

Writing
\[
\mathcal E_{xy}=\mathbf e_{xy}\cdot\boldsymbol\sigma,
\qquad
R_{xy}=\operatorname{Ad}(W_{xy}^{X}),
\]
the connection-lifted affine edge is
\[
\boxed{G_{xy}^{\nabla}=(R_{xy},\mathbf e_{xy})\in SE(3)}.
\]
On the rotationally consistent triangle sector
\[
R_{xz}=R_{xy}R_{yz},
\]
the closed affine loop has
\[
\boxed{R_C=I},
\qquad
\boxed{
\mathbf t_C
=\operatorname{vec}(\mathcal T_{xyz})
=-\operatorname{vec}(\mathcal C_{xyz}).
}
\]
The scalar loop witness is
\[
\boxed{
\tau_C
=\|\mathbf t_C\|
=\sqrt{\frac12\operatorname{Tr}(\mathcal T_{xyz}^2)}.
}
\]
The source theorem is
\path{TIR/foundations/TIR_UNIVERSAL_LOOP_TORSION_SOURCE_BINDING_V0_1.md}.
\vTwelveStatus{B}{--}{PASS}

\section{Gate A2: Cartan refinement}

For an admitted shape-regular smooth family of shrinking relational triangles,
with edge size \(\ell\to0\), the covariant edge realization gives
\[
\mathcal E_{pq}
=e^a{}_{\mu}(p)\Delta x^{\mu}\sigma_a+O(\ell^2).
\]
With oriented area bivector \(\Sigma^{\mu\nu}_{\triangle}=O(\ell^2)\),
the discrete solder object satisfies
\[
\boxed{
\mathcal T_{\triangle}
=\frac12T^a{}_{\mu\nu}(x)
\Sigma^{\mu\nu}_{\triangle}\sigma_a
+O(\ell^3),
}
\]
where
\[
\boxed{T^a=de^a+\omega^a{}_b\wedge e^b}.
\]
The rotational loop satisfies
\[
\boxed{
R_{\triangle}
=I+\frac12\Omega_{\mu\nu}(x)
\Sigma^{\mu\nu}_{\triangle}+O(\ell^3),
}
\]
with
\[
\boxed{
\Omega^a{}_b
=d\omega^a{}_b+\omega^a{}_c\wedge\omega^c{}_b.
}
\]
Thus the corresponding area-normalized quantities converge to the Cartan
torsion and curvature components under the declared refinement assumptions.
The full affine-loop translation differs from the discrete solder torsion only
by \(O(\ell^3)\) on a smooth curved connection, so its area-normalized limit
carries the same torsion coefficient.

Canonical source:
\path{TIR/foundations/TIR_CARTAN_CONTINUUM_REFINEMENT_V0_1.md}.
\vTwelveStatus{B}{--}{PASS}

This is a conditional local/refining theorem.  It does not assert that the
actual global TIR relational complex already supplies the required production
refining family.

\section{Gate A3: zero torsion and Levi-Civita selection}

When the direct \(x\to z\) relation and the connection-composed
\(x\to y\to z\) description are admitted as two representations of the same
primitive ordered endpoint relation in the same comparison frame, intrinsic
affine-displacement uniqueness selects
\[
\boxed{\mathcal C_{xyz}=0}.
\]
Hence
\[
\mathcal T_{xyz}=0,
\]
and Gate A2 carries the endpoint-compatible regular refinement to
\[
\boxed{T^a=0}.
\]
Because the spatial transport is induced through
\[
SU(2)\xrightarrow{\operatorname{Ad}}SO(3),
\]
it preserves the spatial metric, so \(Dh=0\).  The standard uniqueness theorem
for a torsion-free metric-compatible connection then gives
\[
\boxed{D=D^{\rm LC}}.
\]
Zero torsion does not force zero curvature; the sector
\[
T^a=0,
\qquad
\Omega^a{}_b\ne0
\]
remains admitted.

Canonical source:
\path{TIR/foundations/TIR_ZERO_TORSION_LEVI_CIVITA_SELECTION_V0_1.md}.
\vTwelveStatus{B}{--}{PASS}

The broader context-lifted torsional Cartan branch remains typed separately.

\section{Gate A4: leading-loop locality and metric jet}

For a regular shrinking loop with \(A_C=\Theta(\epsilon^2)\), the leading
rotational coefficient is
\[
\boxed{
\frac{R_C-I}{A_C}\longrightarrow\Omega.
}
\]
Curvature-derivative terms enter at higher powers of the refinement scale.
The TIR Leading Refinement Rule selects the lowest finite nonzero covariant
normalized loop coefficient as the primitive continuum carrier.  On the A3
Levi-Civita sector, this gives the leading metric operator class
\[
\boxed{
\mathcal E_{\mu\nu}
=\mathcal E_{\mu\nu}(g,\partial g,\partial^2g),
}
\]
while \(\nabla\Omega,\nabla^2\Omega,\ldots\) remain typed as higher-jet
correction sectors.

Canonical source:
\path{TIR/foundations/TIR_LEADING_LOOP_LOCALITY_METRIC_JET_V0_1.md}.
\vTwelveStatus{B}{--}{PASS}

\section{Gate A5: global three-manifold certificate}

For a supplied finite pure tetrahedral complex, Gate A5 checks the closed
combinatorial three-manifold conditions through tetrahedral incidence and
spherical vertex links.  A passing certificate gives a combinatorial
three-manifold, and the standard Moise bridge supplies compatible smooth
realization in dimension three.  On such a carrier, compatible local
\(SO(3)\)-related coframes glue the metric globally, and the local A3
Levi-Civita restrictions glue by uniqueness.

Canonical source:
\path{TIR/foundations/TIR_GLOBAL_3MANIFOLD_SMOOTH_CERTIFICATE_V0_1.md}.

The certifier and theorem are implemented.  The actual source-owned TIR global
incidence complex has not yet been supplied as a frozen production input.
\vTwelveStatus{B}{--}{OPEN}

The production-input/freeze contract is
\path{TIR/foundations/TIR_GLOBAL_SPATIAL_COMPLEX_INPUT_CONTRACT_V0_1.md}.
The correct firewall is therefore
\[
\boxed{
\text{A5 certifier PASS}
\not\Rightarrow
\text{production global-complex PASS}.
}
\]

\section{Current spatial dependency line}

The synchronized spatial chain is
\[
\boxed{
\begin{aligned}
\mathbb C^2
&\to\operatorname{Herm}_0(2)\cong\mathbb R^3
\to\mathcal E_{xy}
\to W_{ij}^{X}/R_{ij}\\
&\to\mathcal T_{xyz}
\to\text{A2 Cartan refinement}
\to\text{A3 }T^a=0/D^{LC}\\
&\to\text{A4 leading metric jet}
\to\text{A5 global manifold certificate}.
\end{aligned}
}
\]
The remaining spatial/global gate is the frozen production incidence carrier,
followed by the inter-leaf temporal/spacetime join described in the completion
frontier.

\section{Source and validation receipt}

The canonical source chain now includes
\begin{itemize}
\item \path{TIR/foundations/TIR_WIJ_HOLONOMY_CROSSWALK_V0_1.md};
\item \path{TIR/foundations/TIR_SE3_ANCHOR_SOURCE_BINDING_V0_1.md};
\item \path{TIR/foundations/TIR_DISCRETE_SOLDER_FORM_V0_1.md};
\item \path{TIR/foundations/TIR_UNIVERSAL_LOOP_TORSION_SOURCE_BINDING_V0_1.md};
\item \path{TIR/foundations/TIR_CARTAN_CONTINUUM_REFINEMENT_V0_1.md};
\item \path{TIR/foundations/TIR_ZERO_TORSION_LEVI_CIVITA_SELECTION_V0_1.md};
\item \path{TIR/foundations/TIR_LEADING_LOOP_LOCALITY_METRIC_JET_V0_1.md};
\item \path{TIR/foundations/TIR_GLOBAL_3MANIFOLD_SMOOTH_CERTIFICATE_V0_1.md};
\item \path{TIR/foundations/TIR_GLOBAL_SPATIAL_COMPLEX_INPUT_CONTRACT_V0_1.md}.
\end{itemize}
The original deterministic v12 receipt remains generated by
\path{TIR/validation/tir_v12_ch07_holonomy_se3_torsion_v0_1.py}; the v12.1
repository-synchronization gate additionally checks the A2--A5 ownership
markers.

% ---- END TIR/monograph/v12/chapters/ch07_connections_holonomy_se3_solder_torsion.tex ----
\clearpage





\part{Information, Phase and Flavour}
% GENERATED LOCAL-BUILD FLAT PART — DO NOT MERGE

% ---- BEGIN TIR/monograph/v12/chapters/ch08_three_flavour_carrier_su3.tex ----
% !TEX root = ../../tir_monograph_v12.tex
% V12_SOURCE_MAP:
% - TIR/foundations/TIR_KAPPA_FLAVOUR_MIXING_NORMALIZATION_V0_1.md
% - TIR/monograph/appendices/appC_su3_primer.tex
% - TIR/monograph/chapters/ch16_pmns_mixing.tex
% - TIR/monograph/chapters/ch18_ckm_matrix.tex
\chapter[Three-Flavour Carrier]{Three-Flavour Carrier and $SU(3)_F$}
\label{ch:v12_three_flavour_su3}

\paragraph{Established context.}
For a complex three-component carrier, unitary basis changes are described by
$U(3)$, while determinant-one unitary transformations form $SU(3)$.  The Lie
algebra dimension
\[
\dim_{\mathbb R}\mathfrak{su}(N)=N^2-1
\]
is standard.  TIR uses this group-theoretic structure as the family-mixing
carrier for its three-flavour branch.

\vTwelveStatus{B}{--}{PASS}

\section{Three-flavour carrier}
The family branch is represented by
\begin{equation}
\boxed{V_F\cong\mathbb C^3.}
\label{eq:v12-flavour-carrier}
\end{equation}
Choose an ordered flavour basis
\[
\{|f_1\rangle,|f_2\rangle,|f_3\rangle\}.
\]
A normalized family state is
\[
|\psi_F\rangle
=\sum_{i=1}^{3}c_i|f_i\rangle,
\qquad
\sum_{i=1}^{3}|c_i|^2=1.
\]
Full determinant-one family mixing acts through
\begin{equation}
\boxed{U_F\in SU(3)_F.}
\label{eq:v12-su3f-carrier}
\end{equation}
The carrier dimension therefore fixes the flavour multiplicity
\begin{equation}
\boxed{N_F=3.}
\label{eq:v12-flavour-multiplicity}
\end{equation}

\section{Eight independent mixing directions}
The associated mixing algebra is
\[
\mathfrak{su}(3)_F,
\]
with
\begin{equation}
\boxed{
\dim_{\mathbb R}\mathfrak{su}(3)_F
=3^2-1
=8.
}
\label{eq:v12-su3f-dimension}
\end{equation}
If $\{T_a\}_{a=1}^{8}$ is any basis of generators, a local family
transformation may be written
\[
U_F(\boldsymbol\theta)
=\exp\!\left(i\sum_{a=1}^{8}\theta_aT_a\right).
\]
The number eight is basis-independent: it is the dimension of the Lie algebra,
not a property of a particular generator convention.

The TIR symmetric-pair crosswalk supplies the decomposition
\begin{equation}
\mathfrak{su}(3)_F
=\mathfrak{so}(3)\oplus\mathfrak p,
\qquad
8=3+5,
\label{eq:v12-su3-symmetric-pair}
\end{equation}
which gives an independent carrier-level accounting of the same eight real
mixing directions.

\section{Generator--flavour incidence set}
For later normalization bookkeeping define
\begin{equation}
\mathcal C_{\rm mix}
=\{1,\ldots,8\}_{\rm generator}
\times\{1,2,3\}_{\rm flavour}.
\label{eq:v12-mixing-incidence-set}
\end{equation}
Its cardinality is
\begin{equation}
|\mathcal C_{\rm mix}|
=8\cdot3
=24.
\label{eq:v12-mixing-incidence-cardinality}
\end{equation}
At this stage the integer $24$ is only the exact incidence count generated by
the already specified carrier and algebra.  Chapter~\ref{ch:v12_kappa_normalization}
adds the half-turn phase measure and the binary-information numerator to obtain
the TIR normalization.

\section{Downstream mixing matrices}
The PMNS and CKM constructions are treated in
Chapter~\ref{ch:v12_ckm_pmns}.  Their numerical angles, phases and empirical
residuals are downstream sector assignments.  The carrier statements
\eqref{eq:v12-flavour-carrier}--\eqref{eq:v12-su3f-dimension} are the structural
parents shared by those sectors and by the normalization chapter that follows.

% ---- END TIR/monograph/v12/chapters/ch08_three_flavour_carrier_su3.tex ----
\clearpage

% ---- BEGIN TIR/monograph/v12/chapters/ch09_kappa_normalization.tex ----
% !TEX root = ../../tir_monograph_v12.tex
% V12_CANONICAL_KAPPA_OWNER
% Canonical theorem source:
% TIR/foundations/TIR_KAPPA_FLAVOUR_MIXING_NORMALIZATION_V0_1.md
% Validator:
% TIR/validation/tir_kappa_flavour_mixing_normalization_v0_1.py
\chapter[The Kappa Normalization]{The \texorpdfstring{$\kappa=\ln2/(24\pi)$}{kappa = ln2/(24 pi)} Normalization}
\label{ch:v12_kappa_normalization}

\paragraph{Established context.}
The binary-information identity $H_2(1/2)=\ln2$ is standard Shannon
information theory \cite{shannon1948,CoverThomas2006}.  The angular closure of
one full turn is $2\pi$, and $\dim\mathfrak{su}(3)=8$ is the standard
special-unitary Lie-algebra dimension.  TIR combines these established
ingredients with the three-flavour carrier derived in the preceding chapter.

\vTwelveStatus{B}{--}{PASS}

\section{Binary information numerator}
For a binary distribution $(1-p,p)$,
\[
H_2(p)=-p\ln p-(1-p)\ln(1-p).
\]
At the balanced relational coordinate,
\begin{equation}
\boxed{
\sigma_\star=\frac12,
\qquad
I_\star=H_2(1/2)=\ln2.
}
\label{eq:v12-kappa-binary-information}
\end{equation}
This supplies the informational numerator.

\section{Mixing-channel multiplicity}
From Chapter~\ref{ch:v12_three_flavour_su3},
\[
N_F=3,
\qquad
\dim_{\mathbb R}\mathfrak{su}(3)_F=8.
\]
The generator--flavour incidence set therefore has
\begin{equation}
\boxed{
N_{\rm mix}
=N_F\dim\mathfrak{su}(3)_F
=3\cdot8
=24.
}
\label{eq:v12-kappa-mixing-count}
\end{equation}
Equivalently,
\[
N_{\rm mix}
=N_F(N_F^2-1)
=3(3^2-1)
=24.
\]

\section{Half-turn phase unit}
The primitive half coordinate acts on the standard angular closure as
\begin{equation}
\boxed{
\Delta\phi_{1/2}
=\frac12(2\pi)
=\pi.
}
\label{eq:v12-kappa-half-turn}
\end{equation}
Assigning one such phase unit to each primitive generator--flavour incidence
channel gives the total mixing-phase measure
\begin{equation}
\boxed{
\Phi_{\rm mix}
=N_{\rm mix}\Delta\phi_{1/2}
=24\pi.
}
\label{eq:v12-kappa-mixing-phase-measure}
\end{equation}
Thus
\begin{equation}
\boxed{
24\pi
=\underbrace{(3^2-1)}_{8\ \mathrm{mixing\ directions}}
\underbrace{3}_{\mathrm{flavours}}
\underbrace{\pi}_{\mathrm{half\!\!-turn\ phase}}.
}
\label{eq:v12-kappa-denominator-factorization}
\end{equation}

\section{Structural normalization}
The TIR information-per-mixing-phase normalization is
\begin{equation}
\boxed{
\kappa
=\frac{I_\star}{\Phi_{\rm mix}}
=\frac{\ln2}{N_F(N_F^2-1)\pi}
=\frac{\ln2}{3(3^2-1)\pi}
=\frac{\ln2}{24\pi}.
}
\label{eq:v12-kappa-canonical}
\end{equation}
Numerically,
\[
\kappa
=0.009193150006360484\ldots .
\]
Equation~\eqref{eq:v12-kappa-canonical} is the canonical v12 publication
location of the full derivation.  Downstream chapters refer to this equation and
to its theorem/validator provenance.

\section{Independent tetrahedral integer crosscheck}
The spatial branch obtains the regular tetrahedral finite cell with abstract
automorphism group
\[
\operatorname{Aut}(\Delta^3)\cong S_4,
\qquad
|S_4|=24.
\]
Hence
\begin{equation}
\boxed{
3\,\dim\mathfrak{su}(3)_F
=3\cdot8
=24
=|S_4|.
}
\label{eq:v12-kappa-tetra-crosscheck}
\end{equation}
The flavour-mixing incidence route supplies the normalization dependency; the
spatial route supplies an independent finite-symmetry occurrence of the same
integer.

\section{Exact phase-rate consequence}
Let
\[
\omega\equiv\frac{d\phi}{dt}=2\pi f,
\qquad
d\mathcal I\equiv\kappa\,d\phi.
\]
Then
\begin{equation}
\boxed{
\Gamma_{\mathcal I}
\equiv\frac{d\mathcal I}{dt}
=\kappa\omega
=\frac{\ln2}{12}f.
}
\label{eq:v12-kappa-phase-rate}
\end{equation}
One complete angular cycle carries
\begin{equation}
\boxed{
\Delta\mathcal I_{\rm cycle}
=2\pi\kappa
=\frac{\ln2}{12}.
}
\label{eq:v12-kappa-information-cycle}
\end{equation}

\section{Constraint manifold}
For
\[
\mathbf q=(\kappa,\omega,f,\Gamma_{\mathcal I})
\]
and constraints
\[
C_1=\kappa-\frac{\ln2}{24\pi}=0,
\qquad
C_2=\omega-2\pi f=0,
\qquad
C_3=\Gamma_{\mathcal I}-\kappa\omega=0,
\]
the constraint Jacobian has rank three.  The declared subsystem is therefore
one-dimensional and can be parameterized by
\begin{equation}
\boxed{
\mathbf q(f)=
\left(
\frac{\ln2}{24\pi},
2\pi f,
f,
\frac{\ln2}{12}f
\right).
}
\label{eq:v12-kappa-constraint-curve}
\end{equation}

\section{Reproducibility}
The theorem source is
\path{TIR/foundations/TIR_KAPPA_FLAVOUR_MIXING_NORMALIZATION_V0_1.md}.
Its deterministic numerical/algebraic validator is
\path{TIR/validation/tir_kappa_flavour_mixing_normalization_v0_1.py}.

% ---- END TIR/monograph/v12/chapters/ch09_kappa_normalization.tex ----
\clearpage

% ---- BEGIN TIR/monograph/v12/chapters/ch10_discrete_structural_labels.tex ----
% !TEX root = ../../tir_monograph_v12.tex
% V12_SOURCE_MAP:
% - TIR/monograph/chapters/ch03_l_constants.tex
% - TIR/monograph/chapters/ch04_quark_primes.tex
% - TIR/STRUCTURAL_CHOICES.md
% - TIR/foundations/TIR_PLATONIC_L_CONSTANTS_CLOSURE_V0_1.md
% Validators:
% - TIR/validation/tir_v12_discrete_labels_audit_v0_1.py
% - TIR/validation/tir_platonic_l_constants_closure_v0_1.py
\chapter[Discrete Structural Labels]{Discrete Structural Labels -- $L_i$, Primes and Provenance}
\label{ch:v12_discrete_labels}

\paragraph{Established context.}
The Collatz map, prime-number sequence, Platonic classification and finite groups
\(A_4,S_4,A_5\) are standard mathematical objects.  TIR uses selected finite
properties of those objects inside an explicitly typed structural layer.  The
v12.1 synchronization separates three levels: standard mathematical facts, the
TIR selection/closure rule, and any downstream physical interpretation.

\section{Platonic index alphabet}

For a convex regular polyhedron with Schlaefli symbol \(\{p,q\}\),
\[
\frac1p+\frac1q>\frac12,
\qquad p,q\ge3.
\]
The integer solutions are exactly
\[
\boxed{
\{3,3\},\{3,4\},\{3,5\},\{4,3\},\{5,3\}.
}
\]
Thus the complete index alphabet appearing in the convex Platonic
classification is
\[
\boxed{\{3,4,5\}}.
\]
The triangular-face branch is
\[
\{3,3\},\qquad\{3,4\},\qquad\{3,5\},
\]
with orientation-preserving rotational groups
\[
G_3\simeq A_4,
\qquad
G_4\simeq S_4,
\qquad
G_5\simeq A_5,
\]
and orders
\[
|A_4|=12,
\qquad |S_4|=24,
\qquad |A_5|=60.
\]
Equivalently, the spherical triangle groups \((2,3,q)\) have
\[
|G_q|
=\frac{2}{\frac12+\frac13+\frac1q-1},
\qquad q\in\{3,4,5\}.
\]
At \(q=6\), the spherical excess vanishes; hence the finite spherical branch
terminates at \(q=5\).

\section{Common tetrahedral root and subgroup indices}

The tetrahedral rotational group satisfies
\[
A_4\triangleleft S_4,
\]
so
\[
\boxed{[S_4:A_4]=\frac{24}{12}=2}.
\]
A point stabilizer in the natural five-point action of \(A_5\) is isomorphic to
\(A_4\), giving
\[
\boxed{[A_5:A_4]=\frac{60}{12}=5}.
\]
Define the two extension coset spaces over the tetrahedral root,
\[
X_4:=S_4/A_4,
\qquad
X_5:=A_5/A_4.
\]
Then
\[
|X_4|=2,
\qquad
|X_5|=5.
\]
TIR identifies
\[
\boxed{L_4:=|X_4|=2},
\qquad
\boxed{L_5:=|X_5|=5}.
\]

\section[TIR root-extension closure and L3]{TIR root-extension closure and $L_3$}

The TIR structural step is explicit:

\begin{quote}
For the tetrahedral root, the complete non-self-dual Platonic extension carrier
is the disjoint union of all admissible extension coset spaces.
\end{quote}

Therefore
\[
X_3^{\rm closure}:=X_4\sqcup X_5
=S_4/A_4\sqcup A_5/A_4.
\]
Since the union is disjoint,
\[
|X_3^{\rm closure}|=|X_4|+|X_5|=2+5=7,
\]
and TIR defines
\[
\boxed{L_3:=|X_3^{\rm closure}|=7}.
\]
Thus
\[
\boxed{
(L_3,L_4,L_5)
=\left(
|S_4/A_4|+|A_5/A_4|,
|S_4/A_4|,
|A_5/A_4|
\right)
=(7,2,5).
}
\]
\vTwelveStatus{B}{--}{PASS}

The standard finite-group facts are not fitted.  The root-extension closure
rule is a TIR structural definition; the triple follows exactly once that rule
is admitted.  The finite-group derivation alone does not establish a physical
law, a dynamical role for the labels, or a derivation from a specific Kaehler
manifold.

Canonical theorem:
\path{TIR/foundations/TIR_PLATONIC_L_CONSTANTS_CLOSURE_V0_1.md}.

Deterministic validator:
\path{TIR/validation/tir_platonic_l_constants_closure_v0_1.py}.

\section[Independent finite Collatz crosscheck for L3]{Independent finite Collatz crosscheck for $L_3$}

Define the usual Collatz step on positive integers by
\begin{equation}
C(n)=
\begin{cases}
3n+1, & n\ \text{odd},\\
n/2, & n\ \text{even}.
\end{cases}
\label{eq:v12-collatz-map}
\end{equation}
For the single seed required here, direct iteration gives the finite orbit
\begin{equation}
\boxed{
3\to10\to5\to16\to8\to4\to2\to1.
}
\label{eq:v12-collatz-three-orbit}
\end{equation}
There are exactly seven steps before the first occurrence of \(1\).  Hence the
pre-existing arithmetic provenance gives independently
\begin{equation}
\boxed{\operatorname{depth}_{\rm Collatz}(3)=7=L_3.}
\label{eq:v12-L3}
\end{equation}
\vTwelveStatus{B}{--}{PASS}
The validator evaluates only this displayed finite orbit; no global assertion
about all positive Collatz seeds enters the calculation.

\section[Independent twin-prime crosschecks for L4 and L5]{Independent twin-prime crosschecks for $L_4$ and $L_5$}

From the prime pair \((3,5)\), the older arithmetic route gives
\begin{equation}
\boxed{5-3=2=L_4,}
\qquad
\boxed{5=L_5.}
\label{eq:v12-L4-L5}
\end{equation}
\vTwelveStatus{B}{--}{PASS}
These identities are retained as independent crosschecks of the values obtained
from the Platonic subgroup-index construction, not as inputs used to fit that
construction.

Useful downstream combinations include
\[
L_3+L_4=9,
\qquad
L_3-1=6,
\qquad
\frac{L_4}{L_5}=\frac25,
\qquad
L_4^{L_5}=2^5=32.
\]
Additional exact checks connecting the Platonic index alphabet and the three
rotation-group orders are
\[
3^2+4^2=5^2,
\qquad
3+4+5=12=|A_4|,
\]
\[
2(3+4+5)=24=|S_4|,
\qquad
3\cdot4\cdot5=60=|A_5|.
\]
These are consistency checks, not extra derivation premises.

\section{Prime candidate set for quark labels}

The stronger derivation of the \(L_i\) values does not automatically derive the
quark-prime map.  That map retains its own provenance and identifiability gate.
Reserve the first prime \(p_1=2\) for the binary structural branch and define the
ordered candidate list
\begin{equation}
\mathcal P_Q
=(p_2,p_3,p_4,p_5,p_6,p_7)
=(3,5,7,11,13,17).
\label{eq:v12-prime-candidate-list}
\end{equation}
Let the ordered quark-slot list be
\begin{equation}
\mathcal F_Q=(u,d,s,c,b,t).
\label{eq:v12-quark-slot-list}
\end{equation}

\section{Arithmetic constraints and residual ambiguity}

The historical assignment satisfies
\begin{equation}
(u,d,s,c,b,t)=(3,5,7,11,13,17),
\label{eq:v12-quark-prime-assignment}
\end{equation}
together with
\[
d-u=2,
\qquad
s=L_3=7,
\qquad
|s-c|=4,
\qquad
|b-t|=4.
\]
An exhaustive scan over all
\[
6!=720
\]
bijections between \(\mathcal F_Q\) and \(\mathcal P_Q\) shows that these
arithmetic conditions alone leave exactly two survivors:
\begin{align}
(u,d,s,c,b,t)&=(3,5,7,11,13,17),\\
(u,d,s,c,b,t)&=(3,5,7,11,17,13).
\end{align}
Thus the displayed arithmetic relations leave one residual heavy-pair exchange
\(b\leftrightarrow t\).
\vTwelveStatus{D}{--}{PASS}

\section{Typed ordering rule}

The v12 structural selection is the explicit order-preserving map
\begin{equation}
\boxed{
q:\mathcal F_Q\to\mathcal P_Q,
\qquad
q((\mathcal F_Q)_i)=(\mathcal P_Q)_i,
\quad i=1,\ldots,6.
}
\label{eq:v12-quark-prime-typed-rule}
\end{equation}
Equivalently,
\[
q(u)=3,\quad q(d)=5,\quad q(s)=7,\quad q(c)=11,\quad q(b)=13,\quad q(t)=17.
\]
Within the declared ordered candidate set and ordered flavour slots, this rule
has exactly one assignment.  The uniqueness is uniqueness of the typed
selection rule; a scheme- and scale-defined physical quark-mass map remains a
separate open gate.
\vTwelveStatus{B}{--}{PASS}

\section{Reproducibility receipt}

The deterministic audit
\path{TIR/validation/tir_v12_discrete_labels_audit_v0_1.py}
checks the finite Collatz orbit and enumerates all \(720\) prime assignments.
The independent validator
\path{TIR/validation/tir_platonic_l_constants_closure_v0_1.py}
enumerates the five Platonic Schlaefli pairs, constructs the finite permutation
groups and cosets, verifies the subgroup indices and reproduces
\((L_3,L_4,L_5)=(7,2,5)\) without using the older arithmetic values as inputs.

% ---- END TIR/monograph/v12/chapters/ch10_discrete_structural_labels.tex ----
\clearpage

% ---- BEGIN TIR/monograph/v12/chapters/ch11_action_architecture_coefficient_forcing.tex ----
% !TEX root = ../../tir_monograph_v12.tex
% V12_SOURCE_MAP:
% - TIR/monograph/chapters/ch05_electron_action.tex
% - TIR/monograph/chapters/ch06_generation_release.tex
% - TIR/foundations/TIR_COEFFICIENT_ROLE_ORIENTATION_FORCING_V0_1.md
% - TIR/foundations/TIR_COEFFICIENT_MAGNITUDE_PARENT_EVALUATION_V0_1.md
% - TIR/TIR_COMPLETION_FRONTIER_V0_7.md
% Validators:
% - TIR/validation/tir_coefficient_role_orientation_forcing_v0_1.py
% - TIR/validation/tir_coefficient_magnitude_parent_evaluation_v0_1.py
\chapter[Action Architecture and Coefficient Forcing]{Action Architecture and the Coefficient-Forcing Problem}
\label{ch:v12_action_coefficients}

\section{Intrinsic coefficient generator}
The current TIR coefficient architecture is organized by the integer state
\begin{equation}
\boxed{(h,a,b,c)\in\mathbb Z^4}
\label{eq:v12-coefficient-state}
\end{equation}
and generator
\begin{equation}
\boxed{
G(h,a,b,c)
=\frac h2
+a\kappa
+b\frac{\kappa}{L_3}
+c\frac{\kappa^2}{2}.
}
\label{eq:v12-coefficient-generator}
\end{equation}
The normalization $\kappa$ is the canonical quantity derived in
Chapter~\ref{ch:v12_kappa_normalization}; the discrete label $L_3$ is fixed in
Chapter~\ref{ch:v12_discrete_labels}.

\section{Typed role routing}
Four source roles are distinguished:
\[
\mathcal R
=\{R_h,R_a,R_b,R_c\},
\]
with
\begin{align}
R_h&=\text{projective-half-spin},\\
R_a&=\text{generation-release},\\
R_b&=\text{return-axis},\\
R_c&=\text{curvature-holonomy}.
\end{align}
The corresponding generator basis is
\begin{equation}
\mathcal B
=\left\{
\frac12,
\kappa,
\frac{\kappa}{L_3},
\frac{\kappa^2}{2}
\right\}.
\label{eq:v12-generator-basis}
\end{equation}
The canonical role router gives
\begin{equation}
\boxed{
R_h\leftrightarrow h,
\qquad
R_a\leftrightarrow a,
\qquad
R_b\leftrightarrow b,
\qquad
R_c\leftrightarrow c.
}
\label{eq:v12-role-slot-map}
\end{equation}
The four parent signatures are pairwise distinct. Exhaustive permutation of the
four coefficient slots therefore leaves exactly one role-preserving
permutation: the identity.

\vTwelveStatus{B}{--}{PASS}
The deterministic validator is
\path{TIR/validation/tir_coefficient_role_orientation_forcing_v0_1.py}.

\section{Gradient orientation}
For the runtime source operator
\begin{equation}
s(\alpha_s,x)=\tanh(-\alpha_s x),
\qquad
x=\frac{\partial V}{\partial\phi},
\qquad
\alpha_s>0,
\label{eq:v12-gradient-source}
\end{equation}
monotonicity and oddness of $\tanh$ imply
\begin{equation}
\boxed{
\operatorname{sgn}s(\alpha_s,x)
=\operatorname{sgn}(-x)
}
\label{eq:v12-gradient-sign}
\end{equation}
for every finite $x$ and positive $\alpha_s$. The orientation class can
therefore be represented by
\begin{equation}
\boxed{
\chi_{\nabla}
:=\operatorname{sgn}\!\left(-\frac{\partial V}{\partial\phi}\right)
\in\{-1,0,+1\}.
}
\label{eq:v12-gradient-orientation}
\end{equation}

\vTwelveStatus{B}{--}{PASS}

\section{Directed-orbit and chiral orientation}
The directed-orbit source is
\begin{equation}
\chi_O=
\begin{cases}
+1,&\text{attractor}\to\text{satellite},\\
-1,&\text{satellite}\to\text{attractor},\\
0,&\text{neutral},
\end{cases}
\label{eq:v12-orbit-orientation}
\end{equation}
while the chiral-path source is
\begin{equation}
\chi_H=
\begin{cases}
+1,&H_+,\\
-1,&H_-,\\
0,&\text{neutral}.
\end{cases}
\label{eq:v12-chiral-orientation}
\end{equation}
For an atomic/orbital state define the active source-sign set
\[
S=\{\chi_\nabla,\chi_O,\chi_H\}\setminus\{0\}.
\]
If $S$ is nonempty and contains a single sign value, the compatible orientation
is unique:
\begin{equation}
\boxed{
\chi_{\rm TIR}=\chi_\nabla=\chi_O=\chi_H
}
\label{eq:v12-consensus-orientation}
\end{equation}
for every nonzero participating source. States with disagreeing active source
signs are routed to the conflict surface.

\vTwelveStatus{B}{--}{PASS}

\section{What has been forced}
Combining role routing and orientation gives
\begin{equation}
\boxed{
\text{pre-coefficient state}
\longrightarrow
\begin{cases}
R_h\to h,\\
R_a\to a,\\
R_b\to b,\\
R_c\to c,
\end{cases}
\qquad
(\chi_\nabla,\chi_O,\chi_H)
\longrightarrow\chi_{\rm TIR}.
}
\label{eq:v12-coefficient-forcing-chain}
\end{equation}
Thus slot identity and orientation/sign are fixed whenever the declared source
operators are mutually consistent.

\section{Magnitude parent evaluation}
A source audit shows that the historical integration lineage already declares
magnitude-parent expressions for the three charged-lepton action/release
packets. The upstream flavour and Platonic closures supply
\begin{equation}
\boxed{N_F=3,\qquad (L_3,L_4,L_5)=(7,2,5).}
\label{eq:v12-magnitude-upstream-integers}
\end{equation}
Let $I$ denote the unit count, $I=1$, and let $X_i$ denote the typed parent
whose cardinality is $L_i$. Componentwise evaluation of the declared packets
\begin{align}
P_e&=(I,F,I,I),\\
P_{e\mu}&=(0,X_5,X_4,X_3+I),\\
P_{\mu\tau}&=(0,F,I,X_3)
\end{align}
gives
\begin{equation}
\boxed{M(P_e)=(1,3,1,1),}
\label{eq:v12-magnitude-electron}
\end{equation}
\begin{equation}
\boxed{M(P_{e\mu})=(0,5,2,8),}
\label{eq:v12-magnitude-emu}
\end{equation}
and
\begin{equation}
\boxed{M(P_{\mu\tau})=(0,3,1,7).}
\label{eq:v12-magnitude-mutau}
\end{equation}
The evaluation is unique once a parent packet is supplied. It consumes no
measured lepton mass, Yukawa coupling, fitted action or residual-to-target
quantity.

\vTwelveStatus{B}{--}{PASS}
The canonical theorem is
\path{TIR/foundations/TIR_COEFFICIENT_MAGNITUDE_PARENT_EVALUATION_V0_1.md},
and the deterministic validator is
\path{TIR/validation/tir_coefficient_magnitude_parent_evaluation_v0_1.py}.

\section{Residual selector frontier}
The preceding PASS closes arithmetic evaluation of the already-declared parent
expressions; it does not prove that a physical transition uniquely selects those
expressions. The old lineage still marks the transition-level semantic mapping
as \texttt{PROJECT\_MODEL\_ASSIGNMENT}. The remaining coefficient theorem is
therefore the narrower selector
\begin{equation}
\boxed{
\text{coefficient-free transition state}
\longrightarrow
(P_h,P_a,P_b,P_c).
}
\label{eq:v12-coefficient-parent-selector}
\end{equation}
The selector must be evaluated before reading a recovered coefficient tuple or
any mass/Yukawa target.

\vTwelveStatus{B}{--}{OPEN}
This gate is named \texttt{COEFFICIENT\_TRANSITION\_PARENT\_SELECTOR} in
\path{TIR/TIR_COMPLETION_FRONTIER_V0_7.md}.

\section{Circularity firewall}
The coefficient-forcing map admits coefficient-free geometric and routing
parents. Its declared input surface excludes observed particle masses,
observed Yukawa couplings, an already recovered coefficient tuple, a runtime
envelope generated from that tuple, and residual-to-target information. The
retrospective stopping-length pattern remains candidate evidence and is not
promoted by the magnitude parent-evaluation theorem.

\section{Interface to particle sectors}
Historical formulas such as the electron action and the charged-lepton release
operators are migrated to Chapter~\ref{ch:v12_leptons_neutrinos}. There they
are treated as sector constructions over the basis
\eqref{eq:v12-generator-basis}, with their own empirical timing and verdict.
The present chapter owns the upstream generator, typed slot routing, orientation
forcing, exact parent-expression evaluation and the residual transition-parent
selector frontier.

% ---- END TIR/monograph/v12/chapters/ch11_action_architecture_coefficient_forcing.tex ----
\clearpage





\part{Particle and Gauge Sectors}
% GENERATED LOCAL-BUILD FLAT PART — DO NOT MERGE

% ---- BEGIN TIR/monograph/v12/chapters/ch12_charged_leptons_neutrinos.tex ----
% !TEX root = ../../tir_monograph_v12.tex
% V12_SOURCE_MAP:
% - TIR/monograph/chapters/ch05_electron_action.tex
% - TIR/monograph/chapters/ch06_generation_release.tex
% - TIR/monograph/chapters/ch07_lepton_spectrum.tex
% - TIR/monograph/chapters/ch17_neutrino_masses.tex
% - TIR/monograph/chapters/ch30_parameter_summary.tex
% - TIR/monograph/chapters/ch32_pdg2026_validation_addendum.tex
% - TIR/foundations/TIR_COEFFICIENT_MAGNITUDE_PARENT_EVALUATION_V0_1.md
% Validator: TIR/validation/tir_v12_neutrino_action_consistency_v0_1.py
\chapter{Charged Leptons and Neutrinos}
\label{ch:v12_leptons_neutrinos}

\paragraph{Established context.}
Charged-lepton masses and neutrino oscillation observables are externally
measured quantities. The TIR relations collected here are sector constructions
built over the upstream normalization and coefficient architecture. Current
external-data comparison is centralized in Chapter~\ref{ch:v12_evidence_matrix}
rather than repeated in this chapter.

\section{Exponential mass map}
For the mass sectors considered here, TIR uses the dimensionless action map
\begin{equation}
\boxed{
m=E_P\exp\!\left(-\frac{S}{\kappa}\right),
}
\label{eq:v12-exponential-mass-map}
\end{equation}
where $E_P$ is an explicit external scale and $\kappa$ is the structural
normalization of Chapter~\ref{ch:v12_kappa_normalization}.

\vTwelveStatus{B}{--}{OPEN}
The parent-count arithmetic for the three historical charged-lepton coefficient
packets is now closed exactly in
Chapter~\ref{ch:v12_action_coefficients}. What remains open upstream is the
coefficient-free transition-parent selector, followed by the physical binding
and empirical test of the action-to-mass map. A parent-evaluation PASS does not
change the charged-lepton precision verdict below.

\section{Charged-lepton action family}
The historical electron action is
\begin{equation}
\boxed{
S_e
=\frac12-3\kappa+\frac{\kappa}{L_3}-\frac{\kappa^2}{2}.
}
\label{eq:v12-electron-action}
\end{equation}
The two generation-release operators are
\begin{align}
R_{e\mu}
&=5\kappa+\frac{2\kappa}{L_3}+\frac{L_3+1}{2}\kappa^2,
\label{eq:v12-electron-muon-release}\\
R_{\mu\tau}
&=3\kappa-\frac{\kappa}{L_3}-\frac{L_3}{2}\kappa^2.
\label{eq:v12-muon-tau-release}
\end{align}
The derived action chain is
\begin{equation}
\boxed{
S_\mu=S_e-R_{e\mu},
\qquad
S_\tau=S_\mu-R_{\mu\tau}.
}
\label{eq:v12-charged-lepton-actions}
\end{equation}
Using the canonical v12 values of $\kappa$ and $L_3=7$ gives
\begin{align}
S_e&=0.473691600121\ldots,\\
S_\mu&=0.424761179774\ldots,\\
S_\tau&=0.398790835923\ldots .
\end{align}
Inserted into Eq.~\eqref{eq:v12-exponential-mass-map}, the historical formula
family evaluates to
\begin{align}
m_e^{\rm TIR}&=0.511641997\ldots\ \mev,\\
m_\mu^{\rm TIR}&=104.831758\ldots\ \mev,\\
m_\tau^{\rm TIR}&=1767.503883\ldots\ \mev.
\end{align}
These numbers identify the exact formula version being discussed; the later
v11 summary contains a different retrospectively refined numerical triplet.
The current publication comparison therefore lives only in the unified evidence
matrix.

\vTwelveStatus{C}{RETROSPECTIVE}{FAIL}
The 2026 precision audit classifies the charged-lepton sector as arithmetic
coherent but not precision-level agreement with experiment.

\section{Neutrino tetrahedral action branch}
The historical neutrino branch starts from
\begin{equation}
S_{\rm bare}
=\frac{1+L_4/L_3}{2}
=\frac9{14},
\label{eq:v12-neutrino-bare-action}
\end{equation}
and the projected tetrahedral face factor
\begin{equation}
A_{\rm face}
=\frac{(L_4/L_3)^2}{2}
=\frac2{49}.
\label{eq:v12-neutrino-face-area}
\end{equation}
Define the action offset
\begin{equation}
\boxed{
dS
=\kappa A_{\rm face}(1-\kappa).
}
\label{eq:v12-neutrino-offset}
\end{equation}

\section{Formula--value reconciliation}
The v11 neutrino chapter prints
\[
S_1=S_{\rm bare}+\kappa dS,
\]
but its published mass values are generated by the single-offset relation
\begin{equation}
\boxed{
S_1^{\rm rec}=S_{\rm bare}+dS.
}
\label{eq:v12-neutrino-reconciled-S1}
\end{equation}
The deterministic diagnostic
\path{TIR/validation/tir_v12_neutrino_action_consistency_v0_1.py}
checks both equations. Literal evaluation of the printed double-$\kappa$ rule
moves all three printed masses by approximately $4.088\%$, whereas
Eq.~\eqref{eq:v12-neutrino-reconciled-S1} reproduces the historical mass triplet
to the displayed rounding precision.

\vTwelveStatus{D}{RETROSPECTIVE}{PASS}
This is a formula/value consistency diagnosis and reconstruction of the
historical calculation path. It does not promote the reconstructed relation to
an upstream structural theorem.

\section{Neutrino ratio pattern}
The declared mass-ratio labels are
\begin{equation}
\boxed{
r_i\in\{1,L_4,L_3+L_4+1\}=\{1,2,10\}.
}
\label{eq:v12-neutrino-ratios}
\end{equation}
With the reconciled historical action base,
\begin{equation}
S_i=S_1^{\rm rec}-\kappa\ln r_i,
\label{eq:v12-neutrino-actions}
\end{equation}
so the exponential map gives exactly the ratio inheritance
\begin{equation}
\frac{m_i}{m_1}=r_i.
\label{eq:v12-neutrino-ratio-inheritance}
\end{equation}
Numerically this historical formula version gives
\begin{align}
m_1&=0.005009994\ldots\ \eV,\\
m_2&=0.010019988\ldots\ \eV,\\
m_3&=0.050099939\ldots\ \eV,
\end{align}
and therefore
\begin{align}
\Delta m_{21}^2
&=7.5300117\times10^{-5}\ \eV^2,\\
\Delta m_{31}^2
&=2.4849039\times10^{-3}\ \eV^2.
\end{align}

\vTwelveStatus{C}{RETROSPECTIVE}{COMPATIBLE}
The mass-splitting comparison remains a retrospective compatibility surface.
Absolute mass values remain a model output whose external test status is tracked
separately in Chapter~\ref{ch:v12_evidence_matrix}.

\section{Separation from flavour mixing}
PMNS angles and CP phase are not folded into the present mass-action chapter.
They are owned by Chapter~\ref{ch:v12_ckm_pmns}, where the three-flavour carrier
of Chapter~\ref{ch:v12_three_flavour_su3} is already available upstream. This
keeps neutrino mass-scale construction distinct from the flavour-mixing
observable matrix.

% ---- END TIR/monograph/v12/chapters/ch12_charged_leptons_neutrinos.tex ----
\clearpage

% ---- BEGIN TIR/monograph/v12/chapters/ch13_ckm_pmns_flavour_mixing.tex ----
% !TEX root = ../../tir_monograph_v12.tex
% V12_SOURCE_MAP:
% - TIR/monograph/chapters/ch16_pmns_mixing.tex
% - TIR/monograph/chapters/ch18_ckm_matrix.tex
% - TIR/monograph/chapters/ch19_ckm_cp_violation.tex
% - TIR/validation/pdg2026/TIR_PDG2026_VALIDATION_MATRIX_V2.md
% Validators:
% - TIR/validation/tir_v12_flavour_mixing_internal_consistency_v0_1.py
% - TIR/validation/tir_ckm_jarlskog_invariant_closure_v0_1.py
\chapter{CKM/PMNS and Flavour Mixing}
\label{ch:v12_ckm_pmns}

\paragraph{Established context.}
Lepton mixing is conventionally described by the PMNS framework and quark
mixing by the CKM framework.  This chapter collects the TIR arithmetic/geometric
assignments only after the common three-flavour carrier and its $SU(3)_F$
structure have been established in Chapter~\ref{ch:v12_three_flavour_su3}.
External-data comparison is owned by Chapter~\ref{ch:v12_evidence_matrix}.

\section{PMNS structural assignments}
The historical TIR PMNS map uses the discrete labels of
Chapter~\ref{ch:v12_discrete_labels}.  The retained formula family is
\begin{align}
\sin\theta_{13}
&=\frac1{L_3}=\frac17,
\label{eq:v12-pmns-theta13}\\
\sin^2\theta_{12}
&=\frac{L_4}{L_3+L_4}
+\left(\frac{L_4}{L_3}\right)^2,
\label{eq:v12-pmns-theta12}\\
\sin^2\theta_{23}
&=\frac12+\frac{L_4}{L_3^2},
\label{eq:v12-pmns-theta23}\\
\delta_{\rm CP}^{\rm PMNS}
&=\pi\left(
1+\frac{L_4}{L_3}
+\left(\frac{L_4}{L_3}\right)^2
\right).
\label{eq:v12-pmns-delta}
\end{align}
For $(L_3,L_4)=(7,2)$ these evaluate to
\begin{align}
\sin^2\theta_{12}&=0.3038548753\ldots,\\
\sin^2\theta_{23}&=0.5408163265\ldots,\\
\sin^2\theta_{13}&=\frac1{49}=0.0204081633\ldots,\\
\delta_{\rm CP}^{\rm PMNS}&=246.122449^\circ\ldots .
\end{align}

\subsection{Current PMNS verdicts}
The PDG-2026 validation matrix assigns different evidential states to the four
observables; v12 therefore does not compress them into one sector-wide claim.

{\setlength{\tabcolsep}{3pt}
\begin{longtable}{p{0.32\textwidth}p{0.13\textwidth}p{0.22\textwidth}p{0.20\textwidth}}
\toprule
Observable & Claim Class & Timing & Verdict\\
\midrule
$\sin^2\theta_{12}$ & C & RETROSPECTIVE & COMPATIBLE\\
$\sin^2\theta_{23}$ & C & RETROSPECTIVE & OPEN\\
$\sin^2\theta_{13}=1/49$ & C & RETROSPECTIVE & TENSION\\
$\delta_{\rm CP}^{\rm PMNS}$ & C & RETROSPECTIVE & OPEN\\
\bottomrule
\end{longtable}}
The reactor-angle tension is retained without retuning the formula.

\section{CKM primary assignments}
The TIR Cabibbo parameter is
\begin{equation}
\boxed{
\lambda
=\frac{L_4}{L_3+L_4}
+\frac{L_4}{L_3}\kappa
=0.2248488365\ldots .
}
\label{eq:v12-ckm-lambda}
\end{equation}
The two small primary magnitudes used downstream are
\begin{align}
|V_{cb}|
&=\frac{(L_4/L_3)^2}{2}
=\frac2{49}
=0.0408163265\ldots,
\label{eq:v12-ckm-vcb}\\
|V_{ub}|
&=\frac{(L_4/L_3)^2L_4}{(L_3+L_4)L_5}
=\frac8{2205}
=0.0036281179\ldots .
\label{eq:v12-ckm-vub}
\end{align}
The corresponding Wolfenstein coefficient is
\begin{equation}
A=\frac{|V_{cb}|}{\lambda^2}
=0.8073328542\ldots .
\label{eq:v12-ckm-A}
\end{equation}

The displayed historical CKM magnitude matrix is
\begin{equation}
|V_{\rm CKM}|_{\rm hist}
=\begin{pmatrix}
0.97439 & 0.22485 & 0.00363\\
0.22470 & 0.97357 & 0.04082\\
0.00886 & 0.04001 & 0.99916
\end{pmatrix}.
\label{eq:v12-ckm-magnitude-matrix}
\end{equation}
At the printed precision its row- and column-norm residuals are below
$7\times10^{-6}$.  Full complex unitarity and phase-convention statements are
kept distinct from this rounded magnitude check.

\vTwelveStatus{C}{RETROSPECTIVE}{COMPATIBLE}
The 2026 validation matrix reports all nine displayed magnitudes within roughly
$1.71\sigma$ of its selected reference snapshot, while retaining postdiction
provenance.

\section{CKM CP phase and exact Jarlskog invariant}
The direct TIR phase assignment is
\begin{equation}
\boxed{
\delta_{\rm CKM}
=\arccos\frac{L_4}{L_5}
=\arccos\frac25
=66.4218215^\circ\ldots .
}
\label{eq:v12-ckm-delta}
\end{equation}

Once the three CKM mixing amplitudes and this phase are fixed, the Jarlskog
invariant is not an independent parameter.  Let
\begin{equation}
s_{12}=\lambda,
\qquad
s_{23}=|V_{cb}|,
\qquad
s_{13}=|V_{ub}|,
\qquad
c_{ij}=\sqrt{1-s_{ij}^2}.
\label{eq:v12-ckm-sines}
\end{equation}
For the standard three-angle one-phase CKM parameterization,
\begin{equation}
\boxed{
J_{\rm CP}^{\rm CKM}
=\operatorname{Im}
\left(V_{ud}V_{cs}V_{us}^{\ast}V_{cd}^{\ast}\right)
=s_{12}s_{23}s_{13}c_{12}c_{23}c_{13}^{2}\sin\delta_{\rm CKM}.
}
\label{eq:v12-ckm-jarlskog}
\end{equation}
Substituting Eqs.~\eqref{eq:v12-ckm-lambda}--\eqref{eq:v12-ckm-delta}
gives
\begin{equation}
\boxed{
J_{\rm CP}^{\rm CKM}
=2.97106576668\times10^{-5}.
}
\label{eq:v12-ckm-jarlskog-value}
\end{equation}
The dedicated closure validator constructs the full complex CKM matrix,
checks $VV^\dagger=I$ and $\det V=1$ to floating-point roundoff, and verifies
that all nine nontrivial rephasing-invariant quartets carry the same
$|J_{\rm CP}^{\rm CKM}|$ up to orientation sign.

\vTwelveStatus{C}{RETROSPECTIVE}{COMPATIBLE}
The external-data timing classification remains retrospective because the
underlying mixing-magnitude formulas retain their historical provenance.  The
Jarlskog reduction itself is an exact matrix identity conditional on those CKM
inputs.

\section{Historical independent Jarlskog assignment: deprecated}
Earlier TIR versions used the independent expression
\begin{equation}
J_{\rm old}
=\kappa^2\frac{L_4}{L_5}
\left(1-\frac{(L_4/L_5)^2}{2}\right)
=3.11011545905\times10^{-5}.
\label{eq:v12-ckm-j-old}
\end{equation}
It is not the exact Jarlskog invariant of the current CKM matrix and is no
longer retained as an independent CP assignment.  Its gap from the exact matrix
invariant is
\begin{equation}
\boxed{
\frac{|J_{\rm old}-J_{\rm CP}^{\rm CKM}|}{J_{\rm old}}
=4.47088522\%.
}
\label{eq:v12-ckm-j-old-gap}
\end{equation}

Writing
\[
c=\frac{L_4}{L_5}=\frac25,
\]
the exact phase factor is
\[
\sin\delta_{\rm CKM}=\sqrt{1-c^2},
\]
whereas $1-c^2/2$ is only the first nontrivial truncation of the Taylor series
of $\sqrt{1-c^2}$.  In addition, the historical assignment omitted the full
mixing Jacobian
\begin{equation}
\mathcal M
=\frac{s_{12}s_{23}s_{13}c_{12}c_{23}c_{13}^2}{\kappa^2c}
=0.95892344664\ldots .
\label{eq:v12-ckm-mixing-jacobian}
\end{equation}
Thus the exact current relation may be written as
\begin{equation}
\boxed{
J_{\rm CP}^{\rm CKM}
=\kappa^2c\,\mathcal M\sqrt{1-c^2}.
}
\label{eq:v12-ckm-j-factorized}
\end{equation}
The former independent assignment effectively replaced the product
$\mathcal M\sqrt{1-c^2}$ by $1-c^2/2$.

The archived Stage-32 phase backsolve also used a denominator proportional to
$s_{12}s_{23}s_{13}c_{12}c_{23}c_{13}$; the exact Jarlskog identity contains
$c_{13}^2$.  Because $c_{13}\simeq1$, the numerical effect is tiny, but the
formal correction is retained.

\section{Remaining CKM upstream gate}
The Jarlskog invariant is now closed downstream of the current CKM inputs.  A
full first-principles CKM closure still requires an independent forcing theorem
for the functional forms
\begin{equation}
\boxed{
\begin{aligned}
s_{12}
&=\frac{L_4}{L_3+L_4}+\frac{L_4}{L_3}\kappa,\\
s_{23}
&=\frac12\left(\frac{L_4}{L_3}\right)^2,\\
s_{13}
&=\frac12\left(\frac{L_4}{L_3}\right)^2
\frac{L_4}{L_3+L_4}\frac{L_4}{L_5}.
\end{aligned}
}
\label{eq:v12-ckm-upstream-open}
\end{equation}
The new invariant closure does not silently promote these three historical
structural assignments.

\section{GREMLIN Pass-4 selector audit}

The preserved 600-cell/McKay/A2 branch narrows the upstream problem without
changing its evidence class.  The exact mathematical layer now contains:

\begin{itemize}
\item tetrahedral incidence counts \(\nu_F=2\) and \(\nu_E=5\), with
      the structural ratio alphabet \((2/7,2/9,2/5)\);
\item a local centred \(A_2\) three-channel carrier and an exact
      non-adjacent-plane principal cosine \(2/9\);
\item exact finite-operator readouts for \(s_{23}=2/49\) and the
      composite-root factorization \(s_{13}=s_{23}(2/9)(2/5)\);
\item a conditional projective even-sector selector for the rational base
      weights;
\item a half-kernel boundary compression \(9\to7\) and a Kronecker-sum
      realization of the historical scalar readout \(2/9+(2/7)\kappa\);
\item an exact operator-identifiability no-go: the scalar trace alone does not
      determine a unique generator;
\item an exact angle--sine firewall: additivity of a Lie-algebra angle does not
      imply additivity of the sine coordinate.
\end{itemize}

These results are kept as structural/conditional mechanisms.  They do not turn
the historically developed CKM formula into a prospective prediction, and they
do not establish that the 600-cell is a literal physical lattice.  The
remaining physical theorem is a coefficient-free flavour/readout map that
selects the admitted channel operators and their observable sine coordinates
without consuming CKM targets.

\vTwelveStatus{B}{--}{OPEN}

The Pass-4 provenance chain is recorded in
\path{TIR/integration/TIR_CKM_600CELL_FLAG_INCIDENCE_CANDIDATE_V0_1.md}
through
\path{TIR/integration/TIR_CKM_SU3F_ANGLE_SINE_GENERATOR_FIREWALL_V0_8.md}.

\section{Reproducibility}
The retained arithmetic/internal-consistency audit is
\path{TIR/validation/tir_v12_flavour_mixing_internal_consistency_v0_1.py}.
The exact CKM/Jarlskog closure audit is
\path{TIR/validation/tir_ckm_jarlskog_invariant_closure_v0_1.py}.
The theorem surface is
\path{TIR/integration/TIR_CKM_JARLSKOG_INVARIANT_CLOSURE_V0_1.md}.
Current external comparison is read from the unified evidence layer rather than
from the historical tables embedded in the v11 sector chapters.

% ---- END TIR/monograph/v12/chapters/ch13_ckm_pmns_flavour_mixing.tex ----
\clearpage

% ---- BEGIN TIR/monograph/v12/chapters/ch14_hadrons_baryons_mesons.tex ----
% !TEX root = ../../tir_monograph_v12.tex
\chapter{Hadrons — Baryons and Mesons}
\label{ch:v12_hadrons}

\section{Source and status discipline}

This chapter consolidates the historical baryon and meson source family. Standard $SU(3)$ flavour relations are separated from TIR coefficient assignments and from empirical comparison. Current observable verdicts are owned by Chapter~\ref{ch:v12_evidence_matrix}.

Primary legacy sources are
\path{ch08_gmo_formula.tex} through \path{ch15_heavy_mesons.tex}. The v12 arithmetic audit is
\path{TIR/validation/tir_v12_hadron_formula_audit_v0_2.py}.

\section{Baryon octet: the standard GMO carrier}

For the baryon octet the standard Gell--Mann--Okubo organization is
\begin{equation}
M = a+bY+c\left[I(I+1)-\frac{Y^2}{4}\right],
\label{eq:v12_gmo_standard}
\end{equation}
with the familiar octet relation
\begin{equation}
\boxed{
\frac12(M_N+M_\Xi)=\frac34M_\Lambda+\frac14M_\Sigma.
}
\label{eq:v12_gmo_relation}
\end{equation}
\vTwelveStatus{A}{--}{PASS}
The algebraic relation belongs to the established $SU(3)$ flavour framework \cite{GellMann1962,okubo1962}.

The historical TIR parameterization writes
\begin{equation}
M=M_0+\alpha+\beta Y+\gamma\left[I(I+1)-\frac{Y^2}{4}\right].
\label{eq:v12_tir_gmo}
\end{equation}
Its coefficient family is retained as a retrospective model surface.

\section{Octet reference scale: v12 derivation audit}

The legacy source defines
\begin{equation}
S_0=\frac{s+u}{L_3}\,\kappa
\label{eq:v12_legacy_s0}
\end{equation}
and cites the exponential architecture
\begin{equation}
M_0=E_P\exp(-S_0/\kappa).
\label{eq:v12_legacy_m0_exp}
\end{equation}
Equations~\eqref{eq:v12_legacy_s0}--\eqref{eq:v12_legacy_m0_exp} imply
\begin{equation}
\frac{S_0}{\kappa}=\frac{s+u}{L_3}=\frac{10}{7},
\end{equation}
so literal evaluation gives
\begin{equation}
E_Pe^{-10/7}\simeq2.926\times10^{21}\ \mathrm{MeV}.
\end{equation}

The printed $925.95$ MeV value instead comes from a distinct proton-anchored linear map,
\begin{equation}
\boxed{
M_0^{\rm lin}
=E_{\rm proton}
\left(1-\frac{s+u}{L_3}\kappa\right)
\simeq925.9496\ \mathrm{MeV}.
}
\label{eq:v12_octet_m0_linear}
\end{equation}
The corresponding exact proton-anchored exponential would be
\begin{equation}
E_{\rm proton}\exp\!\left[-\frac{s+u}{L_3}\kappa\right]
\simeq926.0302\ \mathrm{MeV}.
\end{equation}

The v12 audit therefore preserves Eq.~\eqref{eq:v12_octet_m0_linear} as the formula that generated the historical table while routing its derivational parentage to repair.
\vTwelveStatus{D}{RETROSPECTIVE}{QUARANTINED}

\section{Octet coefficient family}

The historical coefficient assignments are
\begin{align}
\alpha
&=E_{\rm proton}\kappa
\left(q_s^2-q_u^2-q_d^2+L_3\right),
\label{eq:v12_octet_alpha}\\
\beta
&=-\alpha\frac{L_3^2-1}{L_3^2},
\label{eq:v12_octet_beta}\\
\gamma
&=\alpha\frac{L_4(L_3-L_4)}{L_3^2}.
\label{eq:v12_octet_gamma}
\end{align}
For the canonical labels these give approximately
\[
\alpha=189.765\ \mathrm{MeV},\qquad
\beta=-185.892\ \mathrm{MeV},\qquad
\gamma=38.728\ \mathrm{MeV}.
\]

The legacy chapter attributes Eqs.~\eqref{eq:v12_octet_beta}--\eqref{eq:v12_octet_gamma} to a ``J--KJ'' eigenvalue construction. Version 12 records a completion gate for the explicit operator/matrix, its domain and its eigenvalue calculation. Until that surface is supplied, the coefficient equations are typed retrospective assignments rather than an operator-level theorem.
\vTwelveStatus{C}{RETROSPECTIVE}{QUARANTINED}

With Eqs.~\eqref{eq:v12_octet_m0_linear}--\eqref{eq:v12_octet_gamma}, the constructed masses satisfy Eq.~\eqref{eq:v12_gmo_relation} identically. This is an internal algebraic consistency check of the chosen GMO parameterization. Empirical independence is controlled separately through provenance.

\section{Baryon decuplet}

The historical decuplet uses
\begin{equation}
M=M_0'+\alpha_{10}+\beta_{10}Y,
\label{eq:v12_decuplet}
\end{equation}
with
\begin{align}
M_0'&=E_{\rm proton}\left[1+(s-u)\kappa\right],\\
\alpha_{10}
&=E_{\rm proton}\kappa
\frac{L_3^2+L_4^2+L_5^2+L_3+L_4+L_5+L_4}{L_4},\\
\beta_{10}
&=-E_{\rm proton}\kappa\frac{L_3L_5}{L_4}.
\end{align}
The displayed values are approximately
\[
M_0'=972.72\ \mathrm{MeV},\qquad
\alpha_{10}=405.41\ \mathrm{MeV},\qquad
\beta_{10}=-150.95\ \mathrm{MeV}.
\]
They generate the equal-spaced sequence reported historically for $\Delta$, $\Sigma^*$, $\Xi^*$ and $\Omega$.

The proton mass enters explicitly as an external anchor and the coefficient numerators are discrete structural assignments. The frozen 2026 evidence matrix therefore treats this as postdiction with a proton anchor and candidate compatibility.
\vTwelveStatus{C}{RETROSPECTIVE}{COMPATIBLE}

\section{Pseudoscalar mesons: retained arithmetic failure}

The legacy pion and kaon source prints
\begin{align}
\frac{\Delta S_\pi}{\kappa}&=\frac6\pi,
& m_\pi&=E_Pe^{-6/\pi},\\
\frac{\Delta S_K}{\kappa}&=\zeta(2)-1,
& m_K&=E_Pe^{-(\zeta(2)-1)}.
\end{align}
Literal evaluation with the displayed $E_P=1.22089\times10^{22}$ MeV gives
\begin{align}
E_Pe^{-6/\pi}&\simeq1.808\times10^{21}\ \mathrm{MeV},\\
E_Pe^{-(\zeta(2)-1)}&\simeq6.406\times10^{21}\ \mathrm{MeV}.
\end{align}
These values establish the formula/value failure recorded by the PDG-2026 audit.
\vTwelveStatus{D}{RETROSPECTIVE}{FAIL}

The $\eta$--$\eta'$ source introduces the mixing matrix
\begin{equation}
\mathcal M^2=
\begin{pmatrix}
m_8^2&M_{81}^2\\
M_{81}^2&m_1^2
\end{pmatrix},
\end{equation}
while the numerical entries and their independent derivation remain the active provenance gate for the displayed eigenvalues.
\vTwelveStatus{C}{RETROSPECTIVE}{QUARANTINED}

\section{Heavy and vector mesons}

The historical heavy/vector chapter lists near-identical ``predicted'' and PDG masses for open- and hidden-flavour states while referring to a zeta-operator spectrum. Version 12 requires the operator, input map and reproducible calculation that generates those values before the table can carry evidential weight.
\vTwelveStatus{C}{RETROSPECTIVE}{QUARANTINED}

\section{Hadronic evidence boundary}

The v12 hadronic hierarchy is therefore
\[
\boxed{\begin{aligned}
\text{standard }SU(3)\text{ mass algebra}
&\to \text{TIR coefficient/anchor assignments}\\
&\to \text{reproducible mass calculation}\\
&\to \text{Chapter 19 verdict}.
\end{aligned}}
\]

The audit receipt is
\path{TIR/validation/TIR_V12_HADRON_FORMULA_AUDIT_V0_2.json}. It records \texttt{PASS\_WITH\_QUARANTINES}: the diagnostic checks themselves pass, while the identified legacy formula/provenance surfaces retain their explicit v12 quarantine/failure states. Version 0.2 supersedes the v0.1 validator metric while preserving the same audited scientific witnesses.

% ---- END TIR/monograph/v12/chapters/ch14_hadrons_baryons_mesons.tex ----
\clearpage

% ---- BEGIN TIR/monograph/v12/chapters/ch15_gauge_hypercharge_anomalies.tex ----
% !TEX root = ../../tir_monograph_v12.tex
\chapter[Gauge Structure and Anomalies]{Gauge Structure, Hypercharge and Anomaly Cancellation}
\label{ch:v12_gauge_hypercharge}

\section{Typed gauge transport parent}

The current TIR transport layer uses the common connection-holonomy family
\[
W_{ij}^{(G,R)}
=\mathcal P\exp\!\left(\int_{\gamma_{ij}}A_R\right),
\]
with sector-typed instances including
\[
W_{ij}^{WT}\in U(1),
\qquad
W_{ij}^{X}\in SU(2),
\qquad
W_{ij}^{c}\in SU(3).
\]
The source binding and transformation laws are recorded in
\path{TIR/foundations/TIR_WIJ_HOLONOMY_CROSSWALK_V0_1.md}.

This chapter audits the separate representation/hypercharge layer used by the Standard Model branch.

\section{Hypercharge assignment map}

The legacy TIR formulas are
\begin{align}
Y_Q&=\frac{q_s}{L_3L_4N_c}=\frac16,\\
Y_{uR}&=\frac{L_4}{N_c}=\frac23,\\
Y_{dR}&=-\frac{L_4}{N_cL_4}=-\frac13,\\
Y_{eR}&=-\frac{L_3-L_4}{L_5}=-1,\\
Y_L&=-\frac12\qquad\text{(external input in the legacy surface)}.
\end{align}
For
\[
(q_s,L_3,L_4,L_5,N_c)=(7,7,2,5,3)
\]
the arithmetic map reproduces the conventional one-generation Standard Model hypercharges exactly.

\vTwelveStatus{B}{--}{OPEN}
The arithmetic map is reproducible. Version 12 keeps the derivation/uniqueness gate explicit because $Y_L$ enters as an external parent and several formulas simplify by cancellation of the same structural labels appearing in numerator and denominator.

\section{Anomaly convention}

For anomaly sums, all fermions are represented as left-handed Weyl fields. A right-handed field of hypercharge $Y_R$ is therefore represented by its left-handed charge conjugate with hypercharge $-Y_R$. With this convention the four local one-generation anomaly coefficients are
\begin{align}
\mathcal A_{33Y}
&=2Y_Q-Y_{uR}-Y_{dR},\\
\mathcal A_{22Y}
&=N_cY_Q+Y_L,\\
\mathcal A_{YYY}
&=N_c\left(2Y_Q^3-Y_{uR}^3-Y_{dR}^3\right)
 +\left(2Y_L^3-Y_{eR}^3\right),\\
\mathcal A_{\mathrm{grav}\,Y}
&=N_c\left(2Y_Q-Y_{uR}-Y_{dR}\right)
 +\left(2Y_L-Y_{eR}\right).
\end{align}

Substituting
\[
Y_Q=\frac16,\quad
Y_{uR}=\frac23,\quad
Y_{dR}=-\frac13,\quad
Y_L=-\frac12,\quad
Y_{eR}=-1
\]
gives
\begin{align}
\mathcal A_{33Y}&=0,\\
\mathcal A_{22Y}&=0,\\
\mathcal A_{YYY}&=0,\\
\mathcal A_{\mathrm{grav}\,Y}&=0.
\end{align}

\vTwelveStatus{A}{--}{PASS}
These cancellations are the standard anomaly consistency of the displayed Standard Model representation content. Their exact rational audit is reproduced by
\path{TIR/validation/tir_v12_hypercharge_anomaly_audit_v0_1.py}.

\section[Global SU(2) consistency]{Global $SU(2)$ consistency}

For one generation, the left-handed $SU(2)$ doublets consist of three colour copies of $Q_L$ plus one lepton doublet $L_L$:
\[
N_{\rm doublet}=N_c+1=4.
\]
The count is even, satisfying the Witten global $SU(2)$ consistency condition.
\vTwelveStatus{A}{--}{PASS}

\section{Baryon-current anomaly and legacy A5 repair}

The legacy Chapter 28 labels
\begin{equation}
3Y_Q-Y_{uR}-Y_{dR}=\frac16
\label{eq:v12_legacy_a5}
\end{equation}
as a fifth condition $[SU(2)]^2U(1)_B$. Equation~\eqref{eq:v12_legacy_a5} is built from hypercharges. The baryon-current coefficient instead uses baryon charge $B=1/3$ for each coloured quark doublet. In the conventional $T(\mathbf2)=1/2$ normalization,
\begin{equation}
\boxed{
\mathcal A_{SU(2)^2B}
=N_c\left(\frac13\right)T(\mathbf2)
=3\left(\frac13\right)\left(\frac12\right)
=\frac12.
}
\end{equation}

The nonzero baryon-current anomaly is the expected electroweak anomalous current structure; the v12 repair is the charge/type correction of the legacy A5 expression.
\vTwelveStatus{D}{RETROSPECTIVE}{QUARANTINED}

\section{Result hierarchy}

The v12 gauge/anomaly dependency is
\[
\boxed{\begin{aligned}
\text{typed }W_{ij}\text{ transport}
&\to \text{representation + hypercharge assignment}\\
&\to \text{left-Weyl anomaly sums}\\
&\to \text{current evidence/status layer}.
\end{aligned}}
\]

The exact standard anomaly cancellations and Witten count receive formal \texttt{PASS}. The TIR hypercharge source map remains a structural \texttt{OPEN} gate for derivation/uniqueness, and the legacy A5 label is quarantined with an explicit corrected baryon-current coefficient.

Audit receipt:
\path{TIR/validation/TIR_V12_HYPERCHARGE_ANOMALY_AUDIT_V0_1.json}.

% ---- END TIR/monograph/v12/chapters/ch15_gauge_hypercharge_anomalies.tex ----
\clearpage

% ---- BEGIN TIR/monograph/v12/chapters/ch16_electroweak_higgs.tex ----
% !TEX root = ../../tir_monograph_v12.tex
\chapter{Electroweak and Higgs Sector}
\label{ch:v12_electroweak_higgs}

\section{Structural scale and coupling candidates}

The historical electroweak branch supplies the structural scale
\begin{equation}
v_0
=E_{\rm proton}\left(L_3^2L_5+L_3L_4+L_4\right).
\label{eq:v12_ew_v0}
\end{equation}
For $(L_3,L_4,L_5)=(7,2,5)$ the integer factor is $261$, giving
\begin{equation}
\boxed{v_0=244.888992\ \mathrm{GeV}}.
\end{equation}
The proton mass is an explicit external anchor in this relation.
\vTwelveStatus{B}{RETROSPECTIVE}{OPEN}

The weak-angle candidate is
\begin{equation}
\boxed{
\sin^2\theta_W^{(0)}
=\frac{L_4}{L_3+L_4}+\kappa
=\frac29+\kappa
=0.2314153722\ldots .
}
\label{eq:v12_sin2w0}
\end{equation}
Its precision comparison requires a declared renormalization scheme and scale.
\vTwelveStatus{B}{RETROSPECTIVE}{TENSION}

The separate $SU(2)$ coupling candidate is
\begin{equation}
\boxed{
g_0=\frac{L_4}{L_3}+\frac{L_4}{L_5}
=\frac{24}{35}
=0.6857142857\ldots .
}
\label{eq:v12_g0}
\end{equation}

\section[Tree-level W/Z map]{Tree-level $W/Z$ map}

Using Eqs.~\eqref{eq:v12_ew_v0}--\eqref{eq:v12_g0}, the historical tree-level construction is
\begin{align}
M_W^{(0)}&=\frac{g_0v_0}{2}
=83.9619401\ldots\ \mathrm{GeV},\\
M_Z^{(0)}&=\frac{M_W^{(0)}}{\sqrt{1-\sin^2\theta_W^{(0)}}}
=95.7715725\ldots\ \mathrm{GeV}.
\end{align}
The frozen PDG-2026 comparison assigns precision verdict \texttt{FAIL} to both pole-mass readings.
\vTwelveStatus{C}{RETROSPECTIVE}{FAIL}

The active closure target is the joint transport
\begin{equation}
\boxed{
(g_0,\theta_W^{(0)},v_0)
\xrightarrow{\mathcal R_{EW}(\mu,\mathrm{scheme})}
(g(\mu),\theta_W(\mu),v(\mu))
\longrightarrow
(M_W^{\rm pole},M_Z^{\rm pole}).
}
\label{eq:v12_rew}
\end{equation}
This transport is a shared sector gate rather than an observable-specific correction.

\section{Fine-structure relation}

The independent TIR relation is
\begin{equation}
\boxed{
\alpha_{\rm TIR}^{-1}
=(L_3L_4)^2-L_3^2-L_4L_5+L_4^2\kappa
=137.0367726000\ldots .
}
\label{eq:v12_alpha_inv}
\end{equation}
Its interpretation as the zero-momentum fine-structure constant carries the frozen precision verdict \texttt{FAIL}.
\vTwelveStatus{C}{RETROSPECTIVE}{FAIL}

\section{Electroweak charge-closure gate}

At tree level the standard electroweak relation between the electromagnetic and $SU(2)$ couplings is
\begin{equation}
e=g\sin\theta_W,
\qquad
\alpha=\frac{e^2}{4\pi}.
\label{eq:v12_ew_charge_relation}
\end{equation}
Applying the raw structural pair $(g_0,\theta_W^{(0)})$ gives
\begin{equation}
e_{g\theta}=g_0\sin\theta_W^{(0)}
=0.3298673257\ldots,
\end{equation}
and therefore
\begin{equation}
\boxed{
\alpha_{g\theta}^{-1}
=\frac{4\pi}{e_{g\theta}^2}
=115.4865120\ldots .
}
\label{eq:v12_alpha_from_gtheta}
\end{equation}
By contrast, Eq.~\eqref{eq:v12_alpha_inv} gives $137.0367726\ldots$. Equivalently, combining the TIR $\alpha$ relation with Eq.~\eqref{eq:v12_sin2w0} requires
\begin{equation}
\boxed{g_{\alpha\theta}=0.6294920778\ldots}
\end{equation}
rather than the raw $g_0=24/35$.

The three raw structural relations therefore define distinct pre-transport quantities. Their common electroweak binding is an explicit completion problem:
\begin{equation}
\boxed{
(g_0,\theta_W^{(0)},\alpha_0,v_0)
\xrightarrow{\mathcal R_{EW}(\mu,\mathrm{scheme})}
(g,e,\theta_W,\alpha,v)_{\mu,\mathrm{scheme}}
}
\label{eq:v12_ew_full_closure}
\end{equation}
subject to Eq.~\eqref{eq:v12_ew_charge_relation} at the declared comparison surface.
\vTwelveStatus{B}{--}{OPEN}

\section{Higgs mass relation}

The current Higgs assignment is
\begin{equation}
\boxed{
M_H=v_0\kappa(L_3^2+L_4+L_5)
=v_0\kappa\,56
=126.0728693\ldots\ \mathrm{GeV}.
}
\label{eq:v12_higgs_mass}
\end{equation}
The functional form was introduced after the earlier scalar relation had failed and after the Higgs target was known, so its evidence timing is retrospective. The frozen 2026 precision comparison assigns verdict \texttt{FAIL}.
\vTwelveStatus{C}{RETROSPECTIVE}{FAIL}

The structural completion target is a scalar-sector derivation of the discrete factor $L_3^2+L_4+L_5$ from the field/action geometry, followed by the same scheme/scale discipline used by the wider electroweak sector.

\section{v12 electroweak hierarchy}

The resulting dependency structure is
\[
\boxed{
\begin{array}{c}
(v_0,g_0,\theta_W^{(0)},\alpha_0)_{\rm structural}\\[2pt]
\downarrow\\[-2pt]
\mathcal R_{EW}(\mu,\mathrm{scheme})\\[2pt]
\downarrow\\[-2pt]
(g,e,\theta_W,\alpha,v)_{\mu,\mathrm{scheme}}\\[2pt]
\downarrow\\[-2pt]
(M_W^{\rm pole},M_Z^{\rm pole},M_H^{\rm observable})
\end{array}}
\]
with the current empirical verdicts maintained in Chapter~\ref{ch:v12_evidence_matrix}.

The audit receipt is
\path{TIR/validation/TIR_V12_ELECTROWEAK_HIGGS_AUDIT_V0_1.json}. It records \texttt{PASS\_WITH\_EW\_CLOSURE\_GATE}: the arithmetic reproduction checks pass and the raw electroweak charge-binding mismatch is promoted to an explicit completion gate.

% ---- END TIR/monograph/v12/chapters/ch16_electroweak_higgs.tex ----
\clearpage

% ---- BEGIN TIR/monograph/v12/chapters/ch17_strong_cp_neutron_edm.tex ----
% !TEX root = ../../tir_monograph_v12.tex
\chapter{Strong CP and Neutron EDM}
\label{ch:v12_strong_cp_edm}

\section{Frozen strong-CP map}

The reviewed v11 branch supplies the assignment
\begin{equation}
\boxed{
\theta_{\rm QCD}
=\kappa\left(\frac{L_4}{L_3}\right)^{L_3+L_4+L_5}
=\kappa\left(\frac27\right)^{14}.
}
\label{eq:v12_theta_qcd}
\end{equation}
The exponent $14=L_3+L_4+L_5$ and the suppression operator define the frozen retrospective strong-CP version.
\vTwelveStatus{C}{RETROSPECTIVE}{OPEN}

The current source-completion target is supplied by the non-Abelian colour-holonomy branch:
\begin{equation}
\boxed{
W_{ij}^{c}
\longrightarrow
U_\gamma
\longrightarrow
\text{topological CP phase / sector invariant}
\longrightarrow
\theta_{\rm QCD}.
}
\label{eq:v12_strongcp_source_axis}
\end{equation}
This gives the existing numerical assignment a precise upstream theorem target.

\section{Corrected arithmetic}

The exact numerical suppression in Eq.~\eqref{eq:v12_theta_qcd} is
\begin{equation}
\left(\frac27\right)^{14}
=2.4157243620710218\times10^{-8}.
\end{equation}
With
\[
\kappa=\frac{\ln2}{24\pi}=0.009193150006360484\ldots,
\]
the frozen phase is
\begin{equation}
\boxed{
\theta_{\rm QCD}
=2.220811643453839\times10^{-10}.
}
\label{eq:v12_theta_qcd_numeric}
\end{equation}
The v12 validator reproduces this value independently from the published decimal.
\vTwelveStatus{C}{RETROSPECTIVE}{PASS}
Here \texttt{PASS} is the arithmetic/formula evaluation of the frozen assignment; its physical consequence is tested separately below.

\section{Fixed hadronic conversion}

The publication snapshot uses the conversion coefficient
\begin{equation}
C_n=2.4\times10^{-16}\ e\,\mathrm{cm},
\end{equation}
with
\begin{equation}
\boxed{
d_n=C_n\theta_{\rm QCD}.
}
\label{eq:v12_nedm_map}
\end{equation}
Therefore
\begin{equation}
\boxed{
d_n^{\rm frozen}
=5.329947944289213\times10^{-26}\ e\,\mathrm{cm}.
}
\label{eq:v12_nedm_value}
\end{equation}
The coefficient is an explicit external/hadronic matching input to this version.

\section{Physical gate}

Against the manuscript bound
\begin{equation}
|d_n|<1.8\times10^{-26}\ e\,\mathrm{cm}
\qquad(90\%\ \mathrm{CL}),
\label{eq:v12_nedm_bound}
\end{equation}
the ratio is
\begin{equation}
\boxed{
\frac{d_n^{\rm frozen}}{d_n^{\rm bound}}
=2.96108219127\ldots .
}
\end{equation}
The frozen physical gate therefore has verdict
\vTwelveStatus{D}{RETROSPECTIVE}{FAIL}
The failure remains attached to the exact exponent, normalization and hadronic conversion that generated it.

\section{Versioned replacement rule}

A future strong-CP relation is promoted as a new version only after its upstream source map is frozen. The preferred structural route is Eq.~\eqref{eq:v12_strongcp_source_axis}, followed by a declared hadronic conversion and the same observable-level nEDM gate. This produces an append-only chain
\[
\boxed{
\text{colour holonomy source}
\to\theta_{\rm QCD}
\to\text{hadronic matching}
\to d_n
\to\text{experimental verdict}.
}
\]

\section{Audit receipt}

The arithmetic and empirical gate are reproduced by
\path{TIR/validation/tir_v12_strong_cp_nedm_audit_v0_1.py}.
Its receipt
\path{TIR/validation/TIR_V12_STRONG_CP_NEDM_AUDIT_V0_1.json}
reports \texttt{PASS\_ARITHMETIC\_PHYSICAL\_FAIL\_RETAINED}: the frozen calculation is reproducible and the physical \texttt{FAIL} remains visible in Chapter~\ref{ch:v12_evidence_matrix}.

% ---- END TIR/monograph/v12/chapters/ch17_strong_cp_neutron_edm.tex ----
\clearpage







\part{Extensions and Tests}
% GENERATED LOCAL-BUILD FLAT PART — DO NOT MERGE

% ---- BEGIN TIR/monograph/v12/chapters/ch18_cosmological_extension.tex ----
% !TEX root = ../../tir_monograph_v12.tex
\chapter{Cosmological Extension}
\label{ch:v12_cosmology}

\section{Legacy suppression law}

The cosmological branch proposes the dimensionless suppression
\begin{equation}
\boxed{
r_\Lambda
=\left(\frac{L_4}{L_3}\right)^D,
\qquad
D=L_3L_4^{L_5}.
}
\label{eq:v12_cosmo_suppression}
\end{equation}
For $(L_3,L_4,L_5)=(7,2,5)$,
\begin{equation}
D=7\times2^5=224.
\end{equation}
The discrete exponent is therefore reproduced exactly from the declared structural labels.
\vTwelveStatus{B}{--}{PASS}
Here \texttt{PASS} refers to evaluation of the declared discrete expression.

\section{Exact power audit}

Literal evaluation of Eq.~\eqref{eq:v12_cosmo_suppression} gives
\begin{equation}
\boxed{
\left(\frac27\right)^{224}
=1.345110818116154\times10^{-122},
}
\label{eq:v12_cosmo_power}
\end{equation}
with
\begin{equation}
\log_{10}r_\Lambda=-121.87124193446175.
\end{equation}
The historical intermediate $10^{-119.5}$ belongs to a superseded arithmetic snapshot and receives the formula-audit verdict
\vTwelveStatus{D}{RETROSPECTIVE}{FAIL}

\section{Dimensionful reconstruction gate}

A physical vacuum-energy density requires a declared dimensionful reference scale $M_\star$:
\begin{equation}
\boxed{
\rho_\Lambda^{\rm TIR}
=r_\Lambda M_\star^4.
}
\label{eq:v12_cosmo_dimensional}
\end{equation}
The v12 audit evaluates two standard Planck-scale conventions as explicit diagnostic reconstructions.

Using the unreduced Planck mass appearing elsewhere in the monograph,
\begin{equation}
M_P=1.22089\times10^{19}\ \mathrm{GeV},
\end{equation}
Eq.~\eqref{eq:v12_cosmo_dimensional} gives
\begin{equation}
\rho_\Lambda^{(M_P)}
=2.9885753618\times10^{-46}\ \mathrm{GeV}^4.
\end{equation}
Using the reduced Planck mass
\begin{equation}
\bar M_P=2.435\times10^{18}\ \mathrm{GeV}
\end{equation}
gives
\begin{equation}
\rho_\Lambda^{(\bar M_P)}
=4.7288324630\times10^{-49}\ \mathrm{GeV}^4.
\end{equation}

The historical displayed value $3.3\times10^{-47}\ \mathrm{GeV}^4$ would correspond, under the same fourth-power conversion, to
\begin{equation}
M_\star\simeq7.037833183\times10^{18}\ \mathrm{GeV}.
\end{equation}
Version 12 therefore promotes the scale convention to an explicit parent variable rather than carrying forward an implicit conversion.
\vTwelveStatus{D}{RETROSPECTIVE}{QUARANTINED}

\section{Density parameter gate}

The dimensionless density parameter is defined by
\begin{equation}
\Omega_\Lambda=\frac{\rho_\Lambda}{\rho_{\rm crit}}.
\label{eq:v12_omega_lambda}
\end{equation}
The legacy chapter prints $\Omega_\Lambda=0.685$. Version 12 requires the corresponding $\rho_{\rm crit}$ input or its derivation, unit convention, epoch and uncertainty model to be frozen alongside Eq.~\eqref{eq:v12_cosmo_dimensional} before the displayed number is promoted as an output of the cosmological branch.
\vTwelveStatus{C}{RETROSPECTIVE}{QUARANTINED}

\section{Typed interpretation of the exponent}

The exponent
\begin{equation}
D=L_3L_4^{L_5}=224
\end{equation}
is retained as the declared TIR composite label. Its factors are kept typed by their own upstream definitions. In particular, standard $SU(3)$ dimensions are carried by the representation or Lie-algebra object to which they refer, while $L_3=7$ retains its v12 finite-orbit provenance from Chapter~10. This removes the legacy type collision between the integer label $L_3$ and a group dimension.
\vTwelveStatus{D}{RETROSPECTIVE}{QUARANTINED}

\section{Cosmological completion surface}

The v12 dependency chain is
\[
\boxed{
(L_3,L_4,L_5)
\to D
\to r_\Lambda=(L_4/L_3)^D
\to M_\star^4
\to\rho_\Lambda
\to\rho_{\rm crit}
\to\Omega_\Lambda
\to\text{Chapter 19 verdict}.
}
\]
This chain makes the dimensionless suppression, dimensionful scale, cosmological reference density and empirical comparison separately auditable.

\section{Audit receipt}

The arithmetic and scale diagnostics are reproduced by
\path{TIR/validation/tir_v12_cosmology_formula_audit_v0_1.py}.
The receipt
\path{TIR/validation/TIR_V12_COSMOLOGY_FORMULA_AUDIT_V0_1.json}
reports \texttt{PASS\_WITH\_COSMOLOGY\_QUARANTINE}: the exact discrete power is reproduced, while the historical intermediate, dimensionful conversion and density-parameter path are routed to their explicit v12 repair gates.

% ---- END TIR/monograph/v12/chapters/ch18_cosmological_extension.tex ----
\clearpage

% ---- BEGIN TIR/monograph/v12/chapters/ch19_unified_evidence_matrix.tex ----
% !TEX root = ../../tir_monograph_v12.tex
\chapter{Unified Evidence Matrix}
\label{ch:v12_evidence_matrix}

\section{Evidence ownership}

Version 12 assigns one publication owner to the current result status of each observable: this chapter. Sector chapters own their equations, derivations and source provenance; historical comparison tables remain versioned evidence snapshots. Current status is normalized here as
\[
\boxed{(\text{Claim Class},\ \text{Timing},\ \text{Verdict})}
\]
using the taxonomy in \path{TIR/monograph/v12/STATUS_TAXONOMY.md}.

The principal frozen external-validation source for particle-sector rows is
\path{TIR/validation/pdg2026/TIR_PDG2026_VALIDATION_MATRIX_V2.md}, frozen 15 August 2026. The retained strong-CP/neutron-EDM failure is sourced from \path{TIR/FALSIFIABILITY.md}. Formula-level diagnostics discovered during the v12 migration are sourced from their dedicated validation receipts.

\section{Current matrix}

{\footnotesize
\setlength{\tabcolsep}{3pt}
\begin{longtable}{p{0.17\textwidth}p{0.10\textwidth}p{0.19\textwidth}p{0.18\textwidth}p{0.25\textwidth}}
\toprule
Observable / sector & Claim Class & Timing & Verdict & Current evidence note \\
\midrule
\endfirsthead
\toprule
Observable / sector & Claim Class & Timing & Verdict & Current evidence note \\
\midrule
\endhead
CKM nine magnitudes & C & RETROSPECTIVE & COMPATIBLE & PDG-2026 V2: all nine displayed magnitudes within the cited compatibility range; largest listed pull about $1.71\sigma$. \\
$\delta_{\rm CKM}$ & C & RETROSPECTIVE & COMPATIBLE & PDG-2026 V2 strong compatibility; timing remains retrospective because the mixing-weight formulas retain historical provenance. \\
$J_{\rm CP}$ from current CKM matrix & C & RETROSPECTIVE & COMPATIBLE & Exact downstream quartet identity once $(s_{12},s_{23},s_{13},\delta)$ are fixed; the former independent $J_{\rm old}$ assignment is deprecated, while external-data timing remains retrospective. \\
\midrule
$\sin^2\theta_{12}$ & C & RETROSPECTIVE & COMPATIBLE & Solar-angle comparison remains compatible in the frozen 2026 matrix. \\
$\Delta m^2_{21}$ & C & RETROSPECTIVE & COMPATIBLE & Solar mass splitting remains compatible in the frozen 2026 matrix. \\
$\sin^2\theta_{23}$ & C & RETROSPECTIVE & OPEN & Octant-sensitive comparison; frozen matrix records ambiguous compatibility. \\
$\sin^2\theta_{13}=1/49$ & C & RETROSPECTIVE & TENSION & Frozen 2026 comparison gives approximately $-3.98$ to $-2.66\sigma$. \\
$\delta_{\rm CP}^{\rm PMNS}$ & C & RETROSPECTIVE & OPEN & Frozen 2026 matrix records the comparison as inconclusive. \\
Absolute neutrino masses & C & RETROSPECTIVE & OPEN & Model output; oscillation data in the frozen matrix do not independently determine the displayed absolute spectrum. \\
Neutrino absolute-action equation & D & RETROSPECTIVE & QUARANTINED & v12 HOUND audit found a printed formula/value mismatch in legacy Ch.~17; the diagnostic reconstruction is retained in the dedicated v12 receipt. \\
\midrule
Electron mass relation & C & RETROSPECTIVE & FAIL & PDG-2026 V2 precision failure. \\
Muon mass relation & C & RETROSPECTIVE & FAIL & PDG-2026 V2 precision failure. \\
Tau mass relation & C & RETROSPECTIVE & FAIL & PDG-2026 V2 precision failure. \\
\midrule
Baryon octet pattern & C & RETROSPECTIVE & QUARANTINED & PDG-2026 V2 records retrospective refinement against target masses; independent-evidence promotion is blocked pending provenance closure. \\
Baryon decuplet pattern & C & RETROSPECTIVE & COMPATIBLE & Candidate compatibility with explicit proton-linked anchor provenance. \\
Pion printed exponential & D & RETROSPECTIVE & FAIL & PDG-2026 V2 records an internal arithmetic/formula failure. \\
Kaon printed exponential & D & RETROSPECTIVE & FAIL & PDG-2026 V2 records an internal arithmetic/formula failure. \\
$\eta,\eta'$ mapping & C & RETROSPECTIVE & QUARANTINED & Provenance/derivation closure pending. \\
Heavy/vector meson table & C & RETROSPECTIVE & QUARANTINED & Frozen audit classifies the near-equality table as reproduction rather than independent validation. \\
\midrule
Electroweak scale $v$ & B & RETROSPECTIVE & OPEN & Close relative comparison uses a derived reference; structural transport to a precision observable remains open. \\
$\sin^2\theta_W$ & B & RETROSPECTIVE & TENSION & Renormalization scheme and scale are required for a precision comparison. \\
$\alpha^{-1}(0)$ relation & C & RETROSPECTIVE & FAIL & PDG-2026 V2 precision failure. \\
$M_W$ relation & C & RETROSPECTIVE & FAIL & PDG-2026 V2 precision failure. \\
$M_Z$ relation & C & RETROSPECTIVE & FAIL & PDG-2026 V2 precision failure. \\
$M_H$ relation & C & RETROSPECTIVE & FAIL & Frozen 2026 audit gives a precision-level failure/tension. \\
Quark mass map & B & -- & OPEN & Prime labels exist; a scheme- and scale-defined mass map remains a completion gate. \\
\midrule
Strong-CP $\theta_{\rm QCD}$ assignment & C & RETROSPECTIVE & OPEN & The frozen v11 assignment is retained as the parent of the neutron-EDM gate. \\
Strong-CP $\to$ neutron EDM & D & RETROSPECTIVE & FAIL & $d_n\simeq5.3299\times10^{-26}\,e\,\mathrm{cm}$ against the manuscript bound $1.8\times10^{-26}\,e\,\mathrm{cm}$; retained physical failure. \\
Hypercharge source map & B & -- & OPEN & The arithmetic reproduces the conventional one-generation assignments; derivation/uniqueness remains open because $Y_L$ is an external parent in the legacy source. \\
Standard local anomaly sums + Witten count & A & -- & PASS & v12 Chapter~15 reproduces $\mathcal A_{33Y}=\mathcal A_{22Y}=\mathcal A_{YYY}=\mathcal A_{\mathrm{grav}\,Y}=0$ and the even doublet count $N_{\rm doublet}=4$. \\
Legacy A5 baryon-current row & D & RETROSPECTIVE & QUARANTINED & The historical expression uses hypercharges; the corrected $SU(2)^2B$ coefficient is $1/2$ in the declared $T(\mathbf2)=1/2$ normalization. \\
Cosmological density formula & D & RETROSPECTIVE & QUARANTINED & v12 HOUND arithmetic audit finds $(2/7)^{224}=1.345\ldots\times10^{-122}$, inconsistent with the legacy printed $10^{-119.5}$ intermediate value; Chapter~18 owns the repair audit. \\
\bottomrule
\end{longtable}}

\section{Status normalization rules}

The v12 map preserves source meaning while removing competing vocabularies:
\begin{align*}
\texttt{PASS\_COMPATIBILITY}
&\longrightarrow \texttt{COMPATIBLE},\\
\texttt{PASS\_STRONG\_COMPATIBILITY}
&\longrightarrow \texttt{COMPATIBLE},\\
\texttt{TENSION}
&\longrightarrow \texttt{TENSION},\\
\texttt{SCHEME\_UNDEFINED\_TENSION}
&\longrightarrow \texttt{TENSION},\\
\texttt{FAIL\_PRECISION}
&\longrightarrow \texttt{FAIL},\\
\texttt{FORMULA\_BROKEN}
&\longrightarrow \texttt{FAIL},\\
\texttt{FAIL\_PRECISION\_TENSION}
&\longrightarrow \texttt{FAIL},\\
\texttt{AMBIGUOUS\_COMPATIBILITY}
&\longrightarrow \texttt{OPEN},\\
\texttt{INCONCLUSIVE}
&\longrightarrow \texttt{OPEN},\\
\texttt{NOT\_TESTABLE\_CURRENTLY}
&\longrightarrow \texttt{OPEN},\\
\texttt{PROVENANCE\_CONTAMINATED}
&\longrightarrow \texttt{QUARANTINED},\\
\texttt{PROVENANCE\_INCOMPLETE}
&\longrightarrow \texttt{QUARANTINED},\\
\texttt{REPRODUCTION\_NOT\_VALIDATION}
&\longrightarrow \texttt{QUARANTINED}.
\end{align*}

The Claim Class axis remains independent of the verdict. Thus an exact implementation can coexist with an empirical FAIL, and a retrospective numerical match can be COMPATIBLE while remaining Class C.

\section{Single-owner invariant}

For v12 publication purposes:
\[
\boxed{
\text{sector chapter}\to\text{formula/provenance},
\qquad
\text{Chapter 19}\to\text{current evidence verdict}.
}
\]
A changed external dataset, corrected formula or new validator updates this matrix and its machine-readable receipt first. Historical Ch.~30, Ch.~32 and earlier sector tables remain provenance snapshots.

% ---- END TIR/monograph/v12/chapters/ch19_unified_evidence_matrix.tex ----
\clearpage

% ---- BEGIN TIR/monograph/v12/chapters/ch20_prospective_predictions_falsification.tex ----
% !TEX root = ../../tir_monograph_v12.tex
\chapter{Prospective Predictions and Falsification}
\label{ch:v12_predictions_falsification}

\section{Prospective evidence contract}

Version 12 separates prospective evidence from retrospective sector comparison. A prospective test is represented by
\begin{equation}
\boxed{
\mathcal P
=(F,\mathcal O,\mathcal D,\mathcal R,\mathcal V),
}
\end{equation}
where $F$ is the frozen formula/version identity, $\mathcal O$ the assigned observable, $\mathcal D$ the post-freeze dataset gate, $\mathcal R$ the decision rule and $\mathcal V$ the no-refit/versioning policy.

The active candidate family was frozen on 29 July 2026. Its executable source is
\path{TIR/validation/separable_universal_candidate_family_v10_7.py}.

\section{Separable architecture}

The candidate family uses
\begin{equation}
\boxed{
\ln y_{f,g}=F(S_f)+D(G_g,R_g),
}
\label{eq:v12_prospective_separable}
\end{equation}
with generation-invariant sector coordinate $S_f$, generation coordinate $G_g$, and Ramanujan release $R_g$. The same pair of functions $(F,D)$ is applied across the charged-lepton, down-quark and up-quark sectors.

The frozen structural coordinates are
\begin{align}
S_\ell&=4.5,
&S_d&=3.0876164691516155,
&S_u&=2.920949802484949,\\
G_1&=0.8471633855025916,
&G_2&=0.6088811663290957,
&G_3&=0.23926393244694075,\\
R_1&=0.4754443238393299,
&R_2&=1.0262897039931138,
&R_3&=2.442442115365772.
\end{align}

\section{Exactly three frozen candidates}

\subsection{C1: linear separable action}
\begin{equation}
F_1(S)=-S,
\qquad
D_1(G,R)=-G+R.
\end{equation}

\subsection{C2: Collatz quarter-power action}
\begin{equation}
F_2(S)=-S^{3/4},
\qquad
D_2(G,R)=-G^{3/4}+R.
\end{equation}

\subsection{C3: inverse quarter-power information potential}
\begin{equation}
F_3(S)=\frac{S^{-3/4}}{L_3\kappa},
\qquad
D_3(G,R)=\frac{G^{-3/4}}{L_3\kappa}+R.
\end{equation}

All three candidates carry
\vTwelveStatus{E}{PROSPECTIVE}{OPEN}
until their assigned post-freeze likelihood is evaluated under the frozen rule.

\section{Orthogonal observable pair}

For charm and muon, the generation coordinate is shared. Equation~\eqref{eq:v12_prospective_separable} therefore gives
\begin{equation}
\boxed{
\ln\frac{y_c}{y_\mu}=F(S_u)-F(S_\ell),
}
\label{eq:v12_yc_ymu}
\end{equation}
so the generation-release term cancels identically. The active sector observable is the first qualifying post-29-July-2026 joint ATLAS/CMS direct charm-to-muon Higgs-coupling likelihood.

For charm and top, the sector coordinate is shared, giving
\begin{equation}
\boxed{
\ln\frac{y_c}{y_t}
=D(G_2,R_2)-D(G_3,R_3),
}
\label{eq:v12_yc_yt}
\end{equation}
so the sector term cancels identically. The active generation observable is the first qualifying post-29-July-2026 joint ATLAS/CMS direct charm-to-top Higgs-coupling likelihood.

The earlier charm-to-tau mapping is retained as historical provenance; the v10.7 orthogonalization above is the active v12 contract.

\section{Frozen numerical predictions}

\begin{center}
\begin{tabular}{lrr}
\toprule
Candidate & $y_c/y_\mu$ & $y_c/y_t$\\
\midrule
C1 & 4.850346751338371 & 0.16766796647328305\\
C2 & 2.3521800134268784 & 0.17147213462587316\\
C3 & 6.858021228826228 & $2.8101955040512466\times10^{-11}$\\
\bottomrule
\end{tabular}
\end{center}

The v12 prospective-contract validator independently reconstructs all six ratios from the frozen coordinates and formula family.

\section{Decision rule}

For each candidate and assigned observable, the frozen decision rule is
\begin{equation}
\boxed{
\text{PASS}
\iff
\text{frozen ratio lies inside the published 95\% confidence region}
}
\end{equation}
for the first qualifying post-freeze likelihood. A failed candidate receives verdict `FAIL` for that frozen test. A later formula or observable receives a new version identity and a new prospective contract.

The three-member family is interpreted with an explicit multiplicity-aware comparison policy.

\section{Falsification witnesses from the retrospective programme}

The prospective programme inherits the versioning discipline demonstrated by the retained retrospective witnesses:
\begin{itemize}
\item the frozen strong-CP $\to$ neutron-EDM chain carries a physical `FAIL`;
\item the charged-lepton relations carry precision `FAIL` verdicts;
\item PMNS $\sin^2\theta_{13}=1/49$ carries `TENSION`;
\item legacy pion/kaon exponentials carry formula `FAIL`;
\item neutrino, baryon-octet and cosmological formula surfaces carry explicit `QUARANTINED` diagnostics where derivation or conversion repair is active.
\end{itemize}
These rows remain addressable in Chapter~\ref{ch:v12_evidence_matrix} while future versions are tested independently.

\section{Prospective audit receipt}

The active contract is checked by
\path{TIR/validation/tir_v12_prospective_contract_audit_v0_1.py}.
Its receipt
\path{TIR/validation/TIR_V12_PROSPECTIVE_CONTRACT_AUDIT_V0_1.json}
reports `PASS`: exactly three candidates, the two orthogonal observables, all six frozen ratios, the 29 July 2026 freeze date and the no-refit rule are reproduced.

% ---- END TIR/monograph/v12/chapters/ch20_prospective_predictions_falsification.tex ----
\clearpage

% ---- BEGIN TIR/monograph/v12/chapters/ch21_open_theorems_completion_frontier.tex ----
% !TEX root = ../../tir_monograph_v12.tex
\chapter{Completion Frontier}
\label{ch:v12_completion_frontier}

\paragraph{v12.4 graph reconciliation.}
The v12.3 semantic frontier remains the audited 18 September 2026 snapshot. Version 12.4 does not rewrite that history. The current dependency authority is the repository-root \path{DEPENDENCY_EXPORT.json}; Appendix~\ref{app:v12_relational_zero_graph_reconciliation} records the post-v12.3 cross-repository delta and the corrected relational-zero root. The v12.3 machine-readable ledger and its GREMLIN audit remain historical provenance.

Every gate now has two independent axes:
\[
\boxed{
\text{formal / derivational status}
\quad\oplus\quad
\text{production / physical status}.
}
\]
A closed theorem does not imply a physical PASS; an open production input does
not turn completed mathematics back into an open derivation.

\section{Pinned evidence surface}

The v12.4 reconciliation begins from the following exact repository heads:
\begin{description}
\item[TIR:] \path{853032e019824b836de43634d949da5566396153}.
\item[IDT:] \path{e474f64a079e0d90a039e0abd19d2b482f7e53d2}.
\item[RFC:] \path{664349e2f16b57ee23bf6f5f21f5fe8ada17c5a5}.
\item[Secret-of-a-Half:] \path{492dc2e31546f6fdce798801549d544d8cbda3b4}.
\end{description}

Branch containment and semantic closure are distinct. A branch already
contained in main does not make an older status ledger inside main current.

\section{Closed or narrowed gates}

\subsection{Spatial and local geometric spine}

Gates A2--A4 remain PASS on their declared sectors, and A5 remains an
implemented three-manifold certifier. The unresolved A5 coordinate is the
source-owned production incidence carrier, not the certifier itself.

\subsection{Einstein and ADM dynamics}

The former Einstein constraint/evolution item is stale. Current RFC main records
\[
\boxed{
\begin{aligned}
\text{RF-E8 ADM assembly} &:\quad \mathrm{PASS},\\
\text{RF-E9 extrinsic curvature} &:\quad \mathrm{PASS},\\
\text{RF-E10 Einstein projections} &:\quad \mathrm{PASS},\\
\text{RF-E11 source typing} &:\quad \mathrm{PASS},\\
\text{RF-E12 ADM source constraints} &:\quad \mathrm{PASS},\\
\text{RF-E13 evolution + Bianchi propagation} &:\quad \mathrm{PASS}.
\end{aligned}}
\]
RF-E24 further closes the local Einstein form on its declared selection rules,
\[
\boxed{
G_{\mu\nu}+\Lambda g_{\mu\nu}
=\kappa_E T_{\mu\nu}.
}
\]
RF-E26 and GSC5A provide conditional local-to-global tensor-gluing certifiers.
Therefore the Einstein/ADM derivation is closed locally; the production global
carrier, source lineage and domain coverage remain open.
\vTwelveStatus{B}{--}{PASS}

\subsection{Coefficient magnitudes}

For an already selected parent packet, the magnitude valuation is exact.
The unresolved object is instead the coefficient-free selector
\[
\boxed{
\Sigma:\mathcal S_{\rm precoef}\to(P_h,P_a,P_b,P_c).
}
\]
The identifiability and composition no-go theorems prove that the current coarse
role/orientation selectors are insufficient. Thus magnitude evaluation is
CLOSED while the transition-parent selector remains OPEN.

\subsection{Hypercharge relative uniqueness}

On the declared one-generation, one-Higgs field content with the TIR
normalization anchor, the hypercharge theorem fixes the relative vector
\[
(q,u,d,\ell,e,h)
=q(1,N_c+1,1-N_c,-N_c,-2N_c,N_c).
\]
For \(N_c=3\) and \(q=1/6\),
\[
\boxed{
\left(\frac16,\frac23,-\frac13,-\frac12,-1,\frac12\right).
}
\]
The external \(Y_L\) parent is eliminated and the cubic anomaly follows after
the linear constraints. Relative uniqueness is therefore CLOSED on the declared
field content. Geometric inevitability of the field content or normalization
anchor, and extensions such as an added gauged right-handed neutrino, are
separate questions.
\vTwelveStatus{B}{--}{PASS}

\subsection{Neutrino action repair}

The repository-internal source conflict is resolved:
\[
\boxed{
S_1=S_{\rm bare}+dS,
\qquad
dS=\kappa A_{\rm face}(1-\kappa).
}
\]
The historical extra factor of \(\kappa\) after \(dS\) is quarantined as a
transcription defect. Hence the derivational repair is CLOSED. The empirical
absolute-neutrino mass scale remains OPEN and is not promoted by this repair.

\subsection{Collatz--Fubini--Study phase interface}

The TIR/IDT relational-phase interface is present, and RFC RF-E27 adds a typed
phase fibre over the existing spacetime carrier. The kinematic interface is
therefore CLOSED. RF-E27 explicitly keeps the physical identity of the
connection, coupling scale and any metric/source coupling OPEN.

\section{Production-global geometry frontier}

The global-geometry frontier is now principally a production-evidence problem.
TIR supplies the global spatial source/freeze contract, A5 certifier,
inter-leaf matching contract and the v0.2 physical-realization bundle
certifier. RFC supplies product-clock and matching-flow soldering, source
assembly of the shared atlas, RF-E25 atlas certification, and RF-E26/GSC5
globalization reductions.

The unresolved source-owned coordinates are
\[
\boxed{
\begin{aligned}
&\text{production spatial realization},\\
&\text{production matching/event-placement realization},\\
&\text{production shared atlas and patchwise local-solution receipts},\\
&\text{target-domain coverage and physical source provenance}.
\end{aligned}}
\]
No concrete source-owned physical realization bundle was found on the pinned
TIR main during this audit.
\vTwelveStatus{B}{--}{OPEN}

\section{Gauge and Standard-Model dynamical frontier}

\subsection{Continuum gauge normalization and running}

The typed \(W_{ij}\) carrier, gauge-invariant quark-link coupling and discrete
gauge-matter action have PASS surfaces. What remains open is specifically
\[
\boxed{
W_{ij}\to A_\mu\to F_{\mu\nu}\to S_{YM}
\to\text{physical continuum normalization and running}.
}
\]
\vTwelveStatus{B}{--}{OPEN}

\subsection{Electroweak common transport and Higgs action}

The current audit still reports an inconsistent raw structural triple and
requires one common
\[
\boxed{\mathcal R_{EW}(\mu,\mathrm{scheme})}
\]
rather than observable-specific retuning. The Higgs scalar/action binding is
downstream of that unresolved transport. Both remain OPEN.
\vTwelveStatus{B}{--}{OPEN}

\subsection{Quark mass map}

Hypercharge closure does not supply a quark-mass theorem. A physical
scheme- and scale-defined quark-mass map remains OPEN.
\vTwelveStatus{B}{--}{OPEN}

\section{Hadronic, CP and cosmological frontier}

\subsection{Meson absolute-action baseline}

No current canonical theorem or receipt was found that closes a source-derived
absolute meson action baseline. The legacy pion/kaon formulas remain unsuitable
as that source theorem.
\vTwelveStatus{B}{--}{OPEN}

\subsection{Strong CP}

The current audit separates arithmetic from physics:
\[
\boxed{
\text{legacy arithmetic}=\mathrm{PASS},
\qquad
\text{frozen neutron-EDM gate}=\mathrm{FAIL}.
}
\]
The source derivation of \(\theta_{\rm QCD}\) remains OPEN. Any replacement must
be derived from the holonomic/topological source axis and must not overwrite
the retained physical FAIL.
\vTwelveStatus{B}{--}{OPEN}

\subsection{Cosmological scale binding}

The dimensionless power arithmetic has been audited, but the unit-complete
physical bridge is not closed. The \(\rho_{\rm crit}\) and
\(\Omega_\Lambda\) chain remains quarantined/open. The genuine gate is the
dimensionful scale/critical-density binding, not the already-audited integer
power arithmetic.
\vTwelveStatus{B}{--}{OPEN}

\section{Critical-axis frontier}

Current Secret-of-a-Half main explicitly keeps the global analytic coordinates
\[
\boxed{
\begin{aligned}
&\text{native/full admissible Li--Weil positivity},\\
&\text{global strict critical-axis positivity/nondegeneracy},\\
&\mathrm{PF}_3\text{ and }\mathrm{PF}_\infty\text{ on the actual kernel},\\
&\text{zero-to-native-closure / reciprocal-orbit bridge},\\
&\text{Riemann Hypothesis}
\end{aligned}}
\]
OPEN. An RH-equivalent criterion is not a proof. Operationally,
\[
\boxed{\texttt{riemann\_hypothesis\_in\_closure=false}.}
\]
\vTwelveStatus{B}{--}{OPEN}

\section{Current unresolved core}

After removing stale or over-broad entries, the principal unresolved frontier
is:
\begin{enumerate}
\item source-owned production spatial, matching/event-placement and shared-atlas
      realization, including target-domain coverage and source provenance;
\item the coefficient-free transition-parent selector \(\Sigma\);
\item continuum gauge normalization and running;
\item one common electroweak scheme/scale transport;
\item Higgs scalar/action binding;
\item a scheme/scale-defined quark-mass map and any stronger geometric
      uniqueness extension of the hypercharge normalization anchor;
\item a source-derived meson absolute-action baseline;
\item a holonomic/topological strong-CP source theorem;
\item a unit-complete cosmological scale/\(\rho_{\rm crit}\) binding;
\item physical metric/source coupling of the Collatz--FS phase fibre;
\item native Li--Weil and global critical-axis positivity/nondegeneracy.
\end{enumerate}

The absolute neutrino-mass empirical status remains separately OPEN even though
the source-level neutrino action repair is CLOSED.

\section{Completion invariant}

The v12.3 frontier enforces
\[
\boxed{
\text{closed derivation}\not\equiv\text{physical PASS},
\qquad
\text{open production input}\not\equiv\text{missing mathematics}.
}
\]
This separation replaces the stale v12.1 headline frontier as the current
monograph completion contract.

% ---- END TIR/monograph/v12/chapters/ch21_open_theorems_completion_frontier.tex ----
\clearpage





\appendix
\pagestyle{plain}
% GENERATED LOCAL-BUILD FLAT PART — DO NOT MERGE

% ---- BEGIN TIR/monograph/v12/appendices/appA_formal_proofs_long_algebra.tex ----
% !TEX root = ../../tir_monograph_v12.tex
\chapter{Formal Proofs, Algebraic Firewalls and Legacy Repairs}
\label{app:v12_formal_proofs}

\paragraph{Purpose and provenance.}
This appendix consolidates long algebra that supports Chapters~4--10 and the
diagnostic no-go surfaces used later in the monograph.  Historical appendix
material is preserved through explicit source pointers rather than copied
verbatim.  GREMLIN/HOUND is used here as a bounded adversarial audit layer:
candidate inconsistencies are promoted into this appendix only after an exact
algebraic or deterministic check.

The principal legacy sources are
\path{TIR/monograph/appendices/appB_collatz_tables.tex},
\path{TIR/monograph/appendices/appC_su3_primer.tex},
\path{TIR/monograph/appendices/appD_poincare_disk.tex},
\path{TIR/monograph/appendices/appE_tetrahedron_algebra.tex},
\path{TIR/monograph/appendices/appJ_collatz_quarter_power_scaling.tex},
\path{TIR/monograph/appendices/appK_sector_holonomy_release.tex}, and
\path{TIR/monograph/appendices/appL_up_sector_common_baseline_no_go.tex}.

\section{Finite Collatz statements and the structural labels}

For positive integers define
\[
C(n)=
\begin{cases}
3n+1,&n\ \hbox{odd},\\
n/2,&n\ \hbox{even}.
\end{cases}
\]
The seed required by the canonical v12 label is finite and directly checkable:
\[
3\to10\to5\to16\to8\to4\to2\to1.
\]
There are seven steps before the first occurrence of \(1\), giving the typed
structural label
\[
\boxed{L_3=7.}
\]
The twin-prime source gives
\[
\boxed{L_4=5-3=2,\qquad L_5=5.}
\]
These finite statements are the complete arithmetic input required by the
v12 label layer.  Their publication status is inherited from
Chapter~\ref{ch:v12_discrete_labels}.

A second finite trajectory, historically used to motivate the quark-prime
candidate set, is
\[
7\to22\to11\to34\to17\to52\to26\to13\to40\to20\to10\to5
\to16\to8\to4\to2\to1.
\]
The primes encountered are
\[
7,11,17,13,5,3,
\]
whose sorted set is
\[
\mathcal P_Q=\{3,5,7,11,13,17\}.
\]
This is retained as a finite set-generation observation.  The assignment to
ordered quark slots is owned by the explicit typed map of
Chapter~\ref{ch:v12_discrete_labels}; the exhaustive \(6!=720\) permutation
audit leaves two arithmetic survivors before that ordering rule and one after
it.

\section{Corrected \(SU(3)\) representation algebra}

The three-flavour carrier is \(V_F\cong\mathbb C^3\), with mixing algebra
\(\mathfrak{su}(3)\).  Its real dimension is
\[
\dim_{\mathbb R}\mathfrak{su}(3)=3^2-1=8.
\]
For an \(SU(3)\) irreducible representation with Dynkin labels \((p,q)\),
\[
\boxed{
\dim(p,q)=\frac{(p+1)(q+1)(p+q+2)}{2}.
}
\]
Consequently,
\[
\dim(1,0)=3,\qquad
\dim(0,1)=3,\qquad
\dim(1,1)=8,\qquad
\dim(3,0)=10,\qquad
\dim(2,2)=27.
\]
The mesonic flavour nonet belongs to the reducible decomposition
\[
\boxed{
3\otimes\bar3=8\oplus1,
}
\]
rather than to the \((2,2)\) irrep.  The historical primer row that paired
\((2,2)\) with ``meson nonet'' is therefore retained as a repaired
representation-label mismatch.
\vTwelveStatus{D}{RETROSPECTIVE}{PASS}

For the canonical \(\kappa\) normalization, only the exact carrier facts
\[
N_F=3,\qquad \dim\mathfrak{su}(3)=8,\qquad N_{\rm mix}=3\times8=24
\]
are required.  The full normalization derivation remains owned by
Chapter~\ref{ch:v12_kappa_normalization} and
\path{TIR/foundations/TIR_KAPPA_FLAVOUR_MIXING_NORMALIZATION_V0_1.md}.

\section{Poincar\'e disk normalization repair}

For the unit disk \(\mathbb D=\{z\in\mathbb C:|z|<1\}\) with metric
\[
ds^2=\frac{4\,|dz|^2}{(1-|z|^2)^2},
\]
the curvature is \(K=-1\).  Define the disk cross-ratio modulus
\[
\chi(z_1,z_2)=
\left|
\frac{z_1-z_2}{1-\bar z_1z_2}
\right|.
\]
The geodesic distance compatible with the displayed metric is
\[
\boxed{
d_{\mathbb D}(z_1,z_2)
=2\,\operatorname{artanh}\chi(z_1,z_2).
}
\]
For \(z_1=0\) this reduces to
\[
d_{\mathbb D}(0,r)=\int_0^r\frac{2\,du}{1-u^2}
=2\,\operatorname{artanh}r,
\]
which fixes the normalization directly.  The historical appendix printed the
same expression without the factor \(2\); v12 records that line as a metric
normalization repair.
\vTwelveStatus{D}{RETROSPECTIVE}{PASS}

For a geodesic triangle with angles \(\alpha,\beta,\gamma\),
\[
A_{\mathbb D}=\pi-(\alpha+\beta+\gamma).
\]
A one-angle identification therefore supplies only partial cell data.  v12
assigns any proposed information-area equality to an explicit bridge theorem
with the complete triangle, curvature convention and area normalization stated
before comparison.

\section{Regular tetrahedral Gram algebra}

Let \(\mathbf n_a\in\mathbb R^3\), \(a=1,\ldots,4\), be unit vectors satisfying
\[
\mathbf n_a\cdot\mathbf n_b=-\frac13\qquad(a\ne b).
\]
Their Gram matrix is
\[
G_{ab}=
\begin{cases}
1,&a=b,\\
-1/3,&a\ne b.
\end{cases}
\]
Equivalently,
\[
G=\frac43I_4-\frac13J_4,
\]
where \(J_4\) is the all-ones matrix.  Since \(J_4\) has eigenvalues \(4,0,0,0\),
\[
\operatorname{spec}(G)=\left\{0,\frac43,\frac43,\frac43\right\},
\qquad
\operatorname{rank}G=3.
\]
The null eigenvector is \((1,1,1,1)^T\), hence
\[
\sum_{a=1}^4\mathbf n_a=0.
\]
Every squared edge length is
\[
|\mathbf n_a-\mathbf n_b|^2
=2-2\mathbf n_a\cdot\mathbf n_b
=\frac83,
\]
so the four points form a regular tetrahedron in the three-real-dimensional
carrier.  Permuting the four vertices preserves the Gram matrix, giving
\[
\operatorname{Aut}(\Delta^3)\cong S_4,
\qquad |S_4|=24.
\]
This is the finite-symmetry crosscheck used after the independent
flavour-mixing count \(3\times8=24\).

\section{Legacy operator-tetrahedron firewall}

The historical operator-labelled tetrahedron introduced six sector-specific
edge numbers and then wrote
\[
\cos\theta_{ij}
=
\frac{e_{ik}^2+e_{jk}^2-e_{ij}^2}{2e_{ik}e_{jk}}.
\]
That expression is the planar law of cosines for the angle of a triangle built
from three edge lengths.  A tetrahedral dihedral angle instead requires the
normals of two faces, or an equivalent Cayley--Menger/Gram construction using
the full tetrahedral metric data.  v12 therefore assigns the historical
``dihedral'' line to the diagnostic ledger rather than to the regular
tetrahedral closure theorem.
\vTwelveStatus{D}{RETROSPECTIVE}{QUARANTINED}

The same historical appendix states a volume proportionality to \(\kappa^3\)
without an independently derived Gram determinant and scale map.  v12 assigns
volume/PMNS coupling to the sector-model provenance of Chapters~12--13, while
the geometric tetrahedron theorem is owned by Chapter~6.

\section{Accelerated Collatz quarter-power arithmetic}

For odd \(n\), define the accelerated map
\[
T(n)=\frac{3n+1}{2^{a(n)}},
\qquad
a(n)=\nu_2(3n+1).
\]
Under the standard uniform odd-residue calculation,
\[
\Pr(a=k)=2^{-k},\qquad k\ge1,
\]
so
\[
\mathbb E[a]=\sum_{k\ge1}k2^{-k}=2.
\]
The associated asymptotic geometric-mean multiplier is then
\[
\boxed{
\rho_C
=\exp\!\left(\ln3-\mathbb E[a]\ln2\right)
=\frac34.
}
\]
This arithmetic result is distinct from the particle-model bridge that uses
\(3/4\) as an exponent.  The archived one-anchor charged-fermion trace keeps its
retrospective physical-spectrum failure, while the arithmetic multiplier
remains a technical result.

\section{Common-baseline translation no-go}

Let \(r_g\) be logarithmic residuals for a fixed three-channel up-sector trace,
and apply one common additive baseline \(B\):
\[
r'_g=r_g+B.
\]
For every pair \(g,h\),
\[
r'_g-r'_h=r_g-r_h,
\]
so the residual spread is translation invariant.

\paragraph{Proposition.}
There exists a common \(B\) satisfying
\[
|r_g+B|\le\varepsilon
\quad\text{for every }g
\]
exactly when
\[
\max_g r_g-\min_g r_g\le2\varepsilon.
\]
The smallest achievable uniform error is
\[
\boxed{
\varepsilon_{\min}
=\frac{\max r-\min r}{2}.
}
\]

\paragraph{Proof.}
Each channel requires
\[
B\in[-r_g-\varepsilon,-r_g+\varepsilon].
\]
All intervals intersect exactly when the largest lower endpoint is no greater
than the smallest upper endpoint, which is equivalent to the stated spread
condition.  Centering the two extreme residuals around zero attains half the
spread.

For the archived v10.2 residual vector
\[
(r_u,r_c,r_t)
=(3.1857238256,\,0.0818052322,\,0.2293790515),
\]
the spread is
\[
\Delta r=3.1039185934
\]
and therefore
\[
\varepsilon_{\min}=1.5519592967,
\qquad
e^{\varepsilon_{\min}}=4.7207104\ldots .
\]
The common-baseline architecture is thus excluded for the fixed relative trace
at any target envelope narrower than this minimax bound.
\vTwelveStatus{D}{RETROSPECTIVE}{PASS}

\section{Appendix-A audit ownership}

The deterministic appendix integration gate is
\path{TIR/validation/tir_v12_appendix_integration_audit_v0_1.py}.
Its machine-readable receipt is
\path{TIR/validation/TIR_V12_APPENDIX_INTEGRATION_AUDIT_V0_1.json}.
The gate checks the corrected \(SU(3)\) representation identity, the factor-two
Poincar\'e distance normalization, the regular tetrahedral Gram spectrum, the
legacy dihedral quarantine marker and the common-baseline minimax identity.

% ---- END TIR/monograph/v12/appendices/appA_formal_proofs_long_algebra.tex ----
\clearpage

% ---- BEGIN TIR/monograph/v12/appendices/appB_numerical_tables_reproducibility.tex ----
% !TEX root = ../../tir_monograph_v12.tex
\chapter{Numerical Tables, Receipts and Reproducibility}
\label{app:v12_numerics_reproducibility}

\paragraph{Purpose.}
This appendix is the numerical companion to Chapter~\ref{ch:v12_evidence_matrix}.
Sector chapters own formulas; Chapter~19 owns current evidence verdicts; this
appendix owns reproducible arithmetic, validator locations and frozen receipts.
Historical calculation tables remain available under
\path{TIR/monograph/appendices/}, while v12 records their audited numerical
status rather than reproducing target-matched tables as fresh evidence.

\section{Canonical numerical constants}

The v12 constant surface used by the migrated particle-sector calculations is
\[
L_3=7,\qquad L_4=2,\qquad L_5=5,
\]
and
\[
\boxed{
\kappa=\frac{\ln2}{24\pi}
=0.009193\ldots .
}
\]
The exact ownership chain is
\[
\mathbb C^3
\to SU(3)_F
\to \dim\mathfrak{su}(3)=8
\to N_{\rm mix}=24
\to\Phi_{\rm mix}=24\pi
\to\kappa.
\]
The deterministic source is
\path{TIR/validation/tir_kappa_flavour_mixing_normalization_v0_1.py}.

\section{Active validator index}

The current publication build executes the following v12-specific numerical and
source validators in addition to the foundational geometry suite:
\begin{description}
\item[Discrete labels:]
\path{TIR/validation/tir_v12_discrete_labels_audit_v0_1.py}.
\item[Coefficient forcing:]
\path{TIR/validation/tir_coefficient_role_orientation_forcing_v0_1.py}.
\item[Coefficient selector composition no-go:]
\path{TIR/validation/tir_coefficient_selector_composition_nogo_v0_3.py}.
\item[Neutrino action consistency:]
\path{TIR/validation/tir_v12_neutrino_action_consistency_v0_1.py}.
\item[Neutrino source repair:]
\path{TIR/validation/tir_neutrino_absolute_action_source_repair_v0_1.py}.
\item[CKM/PMNS internal consistency:]
\path{TIR/validation/tir_v12_flavour_mixing_internal_consistency_v0_1.py}.
\item[CKM Jarlskog invariant closure:]
\path{TIR/validation/tir_ckm_jarlskog_invariant_closure_v0_1.py}.
\item[CKM 600-cell incidence:]
\path{TIR/validation/tir_ckm_600cell_flag_incidence_candidate_v0_1.py}.
\item[CKM local A2/operator bridge:]
\path{TIR/validation/tir_ckm_600cell_a2_selector_bridge_candidate_v0_2.py}.
\item[CKM projective even-sector selector:]
\path{TIR/validation/tir_ckm_projective_even_sector_selector_candidate_v0_3.py}.
\item[CKM clean noncommuting family pair:]
\path{TIR/validation/tir_ckm_600cell_clean_noncommuting_family_pair_candidate_v0_4.py}.
\item[CKM lambda0 nested-fraction provenance:]
\path{TIR/validation/tir_ckm_lambda0_nested_fraction_decomposition_v0_5.py}.
\item[CKM half-kernel kappa selector:]
\path{TIR/validation/tir_ckm_half_kernel_kappa_selector_candidate_v0_6.py}.
\item[CKM product-flow uniqueness audit:]
\path{TIR/validation/tir_ckm_half_kernel_product_flow_uniqueness_audit_v0_7.py}.
\item[CKM angle--sine generator firewall:]
\path{TIR/validation/tir_ckm_su3f_angle_sine_generator_firewall_v0_8.py}.
\item[Golden-mean primitive-orbit asymptotic:]
\path{TIR/validation/tir_golden_mean_prime_orbit_asymptotic_v0_1.py}.
\item[Hadron formula audit:]
\path{TIR/validation/tir_v12_hadron_formula_audit_v0_2.py}.
\item[Hypercharge/anomaly audit:]
\path{TIR/validation/tir_v12_hypercharge_anomaly_audit_v0_1.py}.
\item[Hypercharge relative uniqueness:]
\path{TIR/validation/tir_hypercharge_source_uniqueness_v0_1.py}.
\item[Electroweak/Higgs audit:]
\path{TIR/validation/tir_v12_electroweak_higgs_audit_v0_1.py}.
\item[Strong-CP/neutron-EDM audit:]
\path{TIR/validation/tir_v12_strong_cp_nedm_audit_v0_1.py}.
\item[Cosmology formula audit:]
\path{TIR/validation/tir_v12_cosmology_formula_audit_v0_1.py}.
\item[Evidence matrix consistency:]
\path{TIR/validation/tir_v12_evidence_matrix_consistency_v0_2.py}.
\item[Prospective contract:]
\path{TIR/validation/tir_v12_prospective_contract_audit_v0_1.py}.
\item[Source contract:]
\path{TIR/validation/tir_v12_source_contract_v0_7.py}.
\item[Semantic frontier audit:]
\path{TIR/validation/tir_v12_3_semantic_frontier_audit_v0_1.py}.
\item[Appendix integration:]
\path{TIR/validation/tir_v12_appendix_integration_audit_v0_1.py}.
\end{description}

A technical validator PASS certifies its declared arithmetic/source contract.
Observable-level empirical interpretation remains in
Chapter~\ref{ch:v12_evidence_matrix}.

\section{Charged-lepton arithmetic snapshot}

For the historical action coordinates
\[
\begin{aligned}
S_e&=\frac12-3\kappa+\frac{\kappa}{L_3}-\frac{\kappa^2}{2},\\
R_{e\mu}&=5\kappa+\frac{2\kappa}{L_3}
+\frac{L_3+1}{2}\kappa^2,\\
R_{\mu\tau}&=3\kappa-\frac{\kappa}{L_3}
-\frac{L_3}{2}\kappa^2,
\end{aligned}
\]
v12 reproduces
\[
S_e=0.473691600121\ldots,\quad
S_\mu=0.424761179774\ldots,\quad
S_\tau=0.398790835923\ldots .
\]
Under the archived exponential mass map, the corresponding values are
\[
m_e=0.511641997\ldots\ {\rm MeV},
\quad
m_\mu=104.831758\ldots\ {\rm MeV},
\quad
m_\tau=1767.503884\ldots\ {\rm MeV}.
\]
The frozen 2026 evidence audit assigns precision FAIL to all three rows.  The
numbers remain useful as a reproducible historical trace, not as a replacement
for the current evidence verdict.

\section{Neutrino formula/value reconciliation}

The legacy neutrino source defines
\[
S_{\rm bare}=\frac{1+L_4/L_3}{2}=\frac9{14},
\qquad
A_{\rm face}=\frac12\left(\frac{L_4}{L_3}\right)^2=\frac2{49},
\]
and
\[
\delta S=\kappa A_{\rm face}(1-\kappa).
\]
The printed source line uses
\[
S_1^{\rm literal}=S_{\rm bare}+\kappa\,\delta S,
\]
whereas the published mass triplet is reproduced by
\[
\boxed{
S_1^{\rm rec}=S_{\rm bare}+\delta S.
}
\]
The literal rule yields approximately
\[
(0.00521482,\ 0.0104296,\ 0.0521482)\ {\rm eV},
\]
while the reconciled rule yields
\[
(0.00500999,\ 0.01001999,\ 0.05009994)\ {\rm eV}.
\]
The source-repair theorem classifies the extra factor of $\kappa$ as a transcription/composition defect and restores the single-offset rule as the canonical repository derivation. The derivational repair is therefore closed; only physical absolute-mass validation remains open, as recorded in Chapter~21.

\section{Flavour-mixing snapshot}

The migrated PMNS formulas give
\[
\sin^2\theta_{13}=\frac1{49}=0.02040816\ldots,
\]
\[
\sin^2\theta_{12}
=\frac{L_4}{L_3+L_4}
+\left(\frac{L_4}{L_3}\right)^2
=0.30385488\ldots,
\]
\[
\sin^2\theta_{23}
=\frac12+\frac{L_4}{L_3^2}
=0.54081633\ldots,
\]
and
\[
\delta_{\rm PMNS}
=\pi\left[
1+\frac{L_4}{L_3}
+\left(\frac{L_4}{L_3}\right)^2
\right]
=246.122449^\circ\ldots .
\]
The corresponding frozen verdicts are COMPATIBLE, OPEN, TENSION and OPEN in the
order \(\theta_{12},\theta_{23},\theta_{13},\delta_{\rm PMNS}\).

For CKM,
\[
\lambda
=\frac{L_4}{L_3+L_4}
+\frac{L_4}{L_3}\kappa
=0.2248488365\ldots,
\]
\[
|V_{cb}|=\frac2{49}=0.0408163265\ldots,
\qquad
|V_{ub}|=\frac8{2205}=0.0036281179\ldots,
\]
and
\[
\delta_{\rm CKM}=\arccos\frac25=66.4218215^\circ\ldots .
\]
With
\[
s_{12}=\lambda,\qquad s_{23}=|V_{cb}|,\qquad s_{13}=|V_{ub}|,
\qquad c_{ij}=\sqrt{1-s_{ij}^2},
\]
the exact rephasing-invariant Jarlskog quantity of the current CKM matrix is
\[
\boxed{
J_{\rm CP}^{\rm CKM}
=s_{12}s_{23}s_{13}c_{12}c_{23}c_{13}^2\sin\delta_{\rm CKM}
=2.97106576668\times10^{-5}.
}
\]
The earlier independent structural assignment
\[
J_{\rm old}
=\kappa^2\frac{L_4}{L_5}
\left(1-\frac12\left(\frac{L_4}{L_5}\right)^2\right)
=3.110115459\times10^{-5}
\]
is retained only as historical provenance and is deprecated as an independent
Jarlskog invariant.  Its relative gap from the exact current CKM invariant is
\[
4.47088522\%.
\]
The closure validator also checks all nine nontrivial CKM quartets and the
matrix unitarity/determinant conditions.

\section{Hadron and meson reproduction firewall}

The baryon octet calculation is retained with provenance quarantine because
the historical coefficient construction was refined against target masses.
The decuplet trace remains a compatible anchor-linked candidate.  Their
current result status is defined in Chapter~19 rather than by the old appendix
summary percentages.

The legacy pion and kaon exponentials contain an immediate dimensional/arithmetic
witness.  With the displayed Planck-energy scale \(E_P\),
\[
E_Pe^{-6/\pi}
\]
remains of order \(10^{21}\) MeV rather than \(139.57\) MeV.  Likewise,
\[
E_Pe^{-(\pi^2/6-1)}
\]
remains of Planck order rather than \(493.68\) MeV.  The target-valued numbers
printed beside those expressions therefore carry formula FAIL in v12.

The historical heavy/vector meson tables list measured values as model outputs
without a complete executable derivation chain.  v12 keeps those rows
QUARANTINED as reproduction/provenance surfaces pending the meson absolute-action
baseline theorem.

\section{External-data ownership}

The historical
\path{TIR/monograph/appendices/appH_pdg_tables.tex}
is a dated reference snapshot.  The current particle-sector publication
comparison is owned by
\path{TIR/validation/pdg2026/TIR_PDG2026_VALIDATION_MATRIX_V2.md},
frozen 15 August 2026.  Values with renormalization-scheme or scale dependence
retain that typing in the current matrix.

This separation prevents an old numerical table from silently becoming the
current validation source after the theory layer changes.

\section{Exact-head publication reproducibility}

The authoritative v12 workflow is
\path{.github/workflows/compile-tir-monograph-v12.yml}.
For a pull request it checks out the exact PR-head SHA, verifies that checkout,
runs the deterministic validator suite, compiles the PDF and then applies a
publication preflight.

The hardened preflight requires:
\begin{enumerate}
\item zero unresolved references or citations and zero multiply-defined labels;
\item zero Hyperref PDF-string token warnings;
\item zero overfull boxes;
\item a successful \texttt{qpdf --check};
\item embedded fonts with zero Type-3 fonts;
\item correct title and author metadata;
\item an exact-head artifact carrying the actual PR-head SHA.
\end{enumerate}

The appendix integration receipt is
\path{TIR/validation/TIR_V12_APPENDIX_INTEGRATION_AUDIT_V0_1.json}.
Together with the per-sector receipts, it distinguishes formal arithmetic
reproducibility from the empirical verdict layer.

% ---- END TIR/monograph/v12/appendices/appB_numerical_tables_reproducibility.tex ----
\clearpage

% ---- BEGIN TIR/monograph/v12/appendices/appC_publication_protocol_prospective_gates.tex ----
% !TEX root = ../../tir_monograph_v12.tex
\chapter{Publication Protocol and Prospective Gates}
\label{app:v12_publication_protocol}

\paragraph{Purpose and source ownership.}
This appendix gives the operational detail behind
Chapter~\ref{ch:v12_predictions_falsification}.  It consolidates the frozen
v10.7 separable candidate family and the publication protocol into one
preregistration-style contract.  Historical source lineage is retained in
\path{TIR/monograph/appendices/appM_prospective_observable_identifiability.tex},
\path{TIR/monograph/appendices/appN_separable_candidate_family.tex}, and
\path{TIR/monograph/appendices/appO_publication_protocol.tex}.

\section{Prospective contract object}

A prospective test is represented by
\[
\boxed{
\mathcal P=(F,\mathcal O,\mathcal D,\mathcal R,\mathcal V),
}
\]
where \(F\) is the frozen formula/version identity, \(\mathcal O\) the assigned
observable, \(\mathcal D\) the post-freeze data gate, \(\mathcal R\) the decision
rule and \(\mathcal V\) the version/no-refit policy.

The active v12 family was frozen on 29 July 2026.  Its executable source is
\path{TIR/validation/separable_universal_candidate_family_v10_7.py}.
All three candidates carry
\vTwelveStatus{E}{PROSPECTIVE}{OPEN}
until the assigned qualifying post-freeze likelihood is evaluated.

\section{Shared separable architecture}

Every candidate has the same typed architecture
\[
\boxed{
\ln y_{f,g}=F(S_f)+D(G_g,R_g).
}
\]
The sector coordinates are
\[
S_\ell=4.5,\qquad
S_d=3.0876164691516155,\qquad
S_u=2.920949802484949,
\]
the generation coordinates are
\[
G_1=0.8471633855025916,\quad
G_2=0.6088811663290957,\quad
G_3=0.23926393244694075,
\]
and the Ramanujan release coordinates are
\[
R_1=0.4754443238393299,\quad
R_2=1.0262897039931138,\quad
R_3=2.442442115365772.
\]
The same functions \(F\) and \(D\) act across charged leptons, down-type quarks
and up-type quarks.

\section{Exactly three frozen candidates}

\paragraph{C1: linear separable action.}
\[
F_1(S)=-S,
\qquad
D_1(G,R)=-G+R.
\]

\paragraph{C2: Collatz quarter-power action.}
\[
F_2(S)=-S^{3/4},
\qquad
D_2(G,R)=-G^{3/4}+R.
\]

\paragraph{C3: inverse quarter-power information potential.}
\[
F_3(S)=\frac{S^{-3/4}}{L_3\kappa},
\qquad
D_3(G,R)=\frac{G^{-3/4}}{L_3\kappa}+R.
\]

The family size is itself frozen:
\[
\boxed{N_{\rm candidate}=3.}
\]
A later functional form receives a new version identity and a new prospective
contract.

\section{Orthogonal observable pair}

Charm and the muon share the second-generation coordinate.  Therefore
\[
\begin{aligned}
\ln y_c&=F(S_u)+D(G_2,R_2),\\
\ln y_\mu&=F(S_\ell)+D(G_2,R_2),
\end{aligned}
\]
so
\[
\boxed{
\ln\frac{y_c}{y_\mu}=F(S_u)-F(S_\ell).
}
\]
This observable isolates the sector contribution.

Charm and top share the up-type sector coordinate:
\[
\begin{aligned}
\ln y_c&=F(S_u)+D(G_2,R_2),\\
\ln y_t&=F(S_u)+D(G_3,R_3),
\end{aligned}
\]
hence
\[
\boxed{
\ln\frac{y_c}{y_t}
=D(G_2,R_2)-D(G_3,R_3).
}
\]
This observable isolates the generation/release contribution.

The active class-A mapping is therefore \(y_c/y_\mu\), while \(y_c/y_t\) is the
generation/release test.  The earlier \(y_c/y_\tau\) mapping remains historical
provenance from the pre-orthogonalization stage.

\section{Frozen numerical predictions}

\begin{center}
\begin{tabular}{lrr}
\toprule
Candidate & \(y_c/y_\mu\) & \(y_c/y_t\)\\
\midrule
C1 & 4.850346751338371 & 0.16766796647328305\\
C2 & 2.3521800134268784 & 0.17147213462587316\\
C3 & 6.858021228826228 & \(2.8101955040512466\times10^{-11}\)\\
\bottomrule
\end{tabular}
\end{center}

The deterministic v12 prospective validator reconstructs all six values from
the frozen coordinates and candidate definitions.

\section{Dataset admission gate}

For each observable, the admissible dataset is the first qualifying
post-29-July-2026 joint ATLAS/CMS direct likelihood that reports the relevant
coupling ratio or sufficient likelihood information to construct it on a common
scheme/scale surface.

Dataset admission is decided before reading the numerical location of the
likelihood maximum relative to the three predictions.  If a release lacks the
required joint information, the contract remains OPEN for that release.

\section{Decision rule}

For candidate \(C_i\) and assigned observable \(O\),
\[
\boxed{
\operatorname{PASS}(C_i,O)
\iff
r_i(O)\in{\rm CR}_{95\%}(O),
}
\]
where \(r_i(O)\) is the frozen predicted ratio and
\({\rm CR}_{95\%}(O)\) is the published 95\% confidence region of the admitted
likelihood.

A candidate outside the admitted confidence region receives FAIL for that
frozen test.  The formula identity remains fixed in the corresponding receipt.

\section{Multiplicity-aware interpretation}

There are three simultaneously frozen candidates and two orthogonal observables.
The publication report therefore presents the full \(3\times2\) result surface,
rather than selecting only the most favourable surviving candidate after the
data are visible.

Candidate ranking, if performed, must use the same admitted likelihood surface
for all candidates and must preserve the original three-member family in the
reported ledger.

\section{No-refit and versioning rule}

After freeze:
\begin{enumerate}
\item \(F_i\), \(D_i\), \(S_f\), \(G_g\), \(R_g\), \(L_3\) and \(\kappa\) remain fixed for this contract;
\item the observable identities \(y_c/y_\mu\) and \(y_c/y_t\) remain fixed;
\item an empirical FAIL is recorded without coefficient tuning;
\item a new formula, coordinate or observable receives a new version identity;
\item already inspected retrospective mass tables remain outside the prospective score.
\end{enumerate}

This creates an append-only evidence history: a failed frozen candidate can
coexist with a later new candidate version without rewriting the earlier test.

\section{External-input and scheme discipline}

Yukawa couplings and quark quantities are scheme- and scale-sensitive.
The qualifying experimental surface must therefore provide a common likelihood
or a declared translation to a common scheme before a precision decision is
made.  Ratios constructed from incompatible conventions remain outside the
frozen decision gate.

Upper limits, circular phases and dimensionful anchor quantities retain their
own statistical typing and are not pooled into one global percentage score.

\section{Machine-readable gate}

The prospective contract is checked by
\path{TIR/validation/tir_v12_prospective_contract_audit_v0_1.py}.
The appendix integration audit additionally checks that:
\begin{itemize}
\item exactly three candidate definitions are present;
\item the 29 July 2026 freeze date is present;
\item both orthogonal observables are present;
\item all six frozen ratios are present;
\item the active sector observable is \(y_c/y_\mu\);
\item the no-refit/version rule is explicit.
\end{itemize}

The appendix-level receipt is
\path{TIR/validation/TIR_V12_APPENDIX_INTEGRATION_AUDIT_V0_1.json}.

% ---- END TIR/monograph/v12/appendices/appC_publication_protocol_prospective_gates.tex ----
\clearpage

% ---- BEGIN TIR/monograph/v12/appendices/appD_cross_framework_interfaces.tex ----
% !TEX root = ../../tir_monograph_v12.tex
\chapter{Cross-Framework Interfaces and Dependency Direction}
\label{app:v12_cross_framework}

\paragraph{Purpose.}
This appendix records typed interfaces between TIR and sibling analytic or
dynamical frameworks.  Each interface preserves source direction: exact
mathematical identities remain exact, TIR structural identifications retain
their TIR ownership, and downstream physical bindings keep their own validation
gate.  GREMLIN may propose cross-domain isomorphisms on this surface; theorem or
validator receipts remain the promotion gate.

\section{Half coordinate and binary information}

For binary entropy
\[
H_2(p)=-p\ln p-(1-p)\ln(1-p),
\]
\[
H_2'(p)=\ln\frac{1-p}{p},
\qquad
H_2''(p)=-\frac1p-\frac1{1-p}<0.
\]
Therefore
\[
\boxed{
\operatorname*{arg\,max}_{0<p<1}H_2(p)=\frac12,
\qquad
H_2(1/2)=\ln2.
}
\]
The same midpoint is the unique fixed point of complement \(p\mapsto1-p\).
In the odds coordinate
\[
q=\frac{p}{1-p},
\]
it maps to \(q=1\), the positive fixed point of reciprocal inversion.

TIR uses this exact half/information pair as the numerator side of the canonical
flavour-mixing normalization.

\section{TIR ownership of the \(24\pi\) denominator}

The flavour carrier is
\[
V_F\cong\mathbb C^3,
\qquad
U_F\in SU(3)_F.
\]
Hence
\[
N_F=3,\qquad
\dim_{\mathbb R}\mathfrak{su}(3)_F=8.
\]
The generator--flavour incidence count is
\[
N_{\rm mix}=3\times8=24.
\]
The half coordinate on a full \(2\pi\) angular closure supplies the primitive
half-turn
\[
\Delta\phi_{1/2}=\pi,
\]
so
\[
\Phi_{\rm mix}=24\pi.
\]
Combining the balanced binary information with the mixing-phase measure gives
\[
\boxed{
\kappa
=\frac{H_2(1/2)}{\Phi_{\rm mix}}
=\frac{\ln2}{24\pi}.
}
\]
This is the canonical TIR-internal normalization theorem.  The regular
tetrahedral symmetry \(|S_4|=24\) supplies an independent finite-symmetry
crosscheck of the same integer while Chapter~9 retains sole publication
ownership of the complete derivation.

\section{Secret-of-a-Half negative-inverse interface}

Define
\[
\Omega(s)=\frac{s}{1-s},
\qquad
z_L(s)=1-\frac1s=-\frac1{\Omega(s)}.
\]
Writing \(s=\sigma+it\),
\[
|z_L(s)|^2
=
\frac{(\sigma-1)^2+t^2}{\sigma^2+t^2}.
\]
Therefore
\[
|z_L(s)|=1
\iff
(\sigma-1)^2=\sigma^2
\iff
\boxed{\sigma=\frac12}.
\]
Thus the critical axis is exactly the unit-circle locus of the
negative-inverse coordinate:
\[
\boxed{
\Re s=\frac12
\iff
|z_L(s)|=1.
}
\]

The current cross-framework source is
\path{TIR/interfaces/TIR_SOH_GLOBAL_LI_DOMINATION_BRIDGE_V0_1.md}.
For one reciprocal-conjugate zero quartet with
\[
z_\rho=Re^{i\phi},
\]
its local Li contribution is
\[
L_n(Q_\rho)
=4-2(R^n+R^{-n})\cos(n\phi).
\]
The interface document gives a conditional global-domination theorem candidate:
under the standard symmetric Li representation and classical zero-counting
estimate, an off-axis zero creates an extremal radial shell whose recurrent
phase subsequence dominates the subextremal remainder and yields negative Li
coefficients along a subsequence.

The remaining analytic promotion target in that branch is the independently
derived native positivity implication
\[
\boxed{
\text{native arithmetic closure}
\Longrightarrow
\lambda_n\ge0\quad\forall n.
}
\]
This target remains OPEN in Chapter~21.
\vTwelveStatus{B}{--}{OPEN}

\section{Self-dual fixed point versus extremality}

The half-axis interface also carries a useful diagnostic boundary.  Reciprocal
symmetry by itself fixes a self-dual locus, while dynamical preference requires
an additional variational or positivity condition.

In a finite reciprocal family \(N_n(q)=N_n(1/q)\), the fixed point \(q=1\) may
coexist with off-centre reciprocal extrema.  v12 therefore treats
\[
\boxed{
\text{reciprocal symmetry}
+\text{explicit extremality condition}
\to
\text{extremum claim}
}
\]
as the admissible dependency order.  This prevents symmetry from silently
serving as an unproved stability theorem.

\section{TIR--IDT phase-clock scale interface}

The projective quantities
\[
a_{FS},\qquad
\int\mathcal F_B,\qquad
|\langle\psi_a|\psi_b\rangle|^2,\qquad
\kappa
\]
are dimensionless.  The IDT phase-clock interface supplies the length carrier
\[
\boxed{
\ell_\varphi
=\frac{c}{|\omega_t|}
=c\left|\frac{dt}{d\varphi}\right|
=\frac{\hbar c}{E}.
}
\]
The exact dimensional typing therefore permits the physicalization candidate
\[
\boxed{
ds_{\rm rel}^2
=\ell_\varphi^2\,ds_{FS}^2.
}
\]
For
\[
ds_{FS}^2
=\frac14(d\theta^2+\sin^2\theta\,d\varphi^2),
\]
the associated area candidate is
\[
d\mathcal A_{\rm rel}
=\ell_\varphi^2\,da_{FS}.
\]
With spin-\(1/2\) Berry curvature
\[
\mathcal F_B=\pm2\,da_{FS},
\]
one obtains
\[
d\mathcal A_{\rm rel}
=\pm\frac{\ell_\varphi^2}{2}\mathcal F_B.
\]

The source contract is
\path{TIR/integration/TIR_IDT_INFORMATION_AREA_CURVATURE_V0_2.md},
whose status separates exact dimensional closure from the physical
phase-clock scale-binding candidate.

\section{Polyhedral refinement export}

For a CP1 polyhedral cell complex \(P\) with facewise phase-clock scale,
\[
\mathcal A_{\rm rel}^{(P)}
=
\sum_{f\in P}\ell_{\varphi,f}^2\,a_{FS}(f).
\]
For a uniform cell rate,
\[
\mathcal A_{\rm rel}^{(P)}
=
\frac{c^2}{\omega_P^2}a_{FS}^{(P)}.
\]
The refinement target is
\[
\boxed{
\mathcal A_{\rm rel}^{(P_n)}
\longrightarrow
\mathcal A_{\rm rel}
}
\]
under a controlled refinement sequence.  This target is the dimensional bridge
feeding the TIR--IDT spatial/temporal join.

\section{GR/RFC dependency export}

The v12 completion DAG keeps the relativistic line ordered as
\[
\boxed{
\begin{aligned}
\text{discrete solder/torsion}
&\to\text{Cartan refinement}\\
&\to T^a=0\text{ spatial sector}\\
&\to\text{TIR--IDT ADM join}\\
&\to\text{Einstein constraint/evolution closure}.
\end{aligned}
}
\]
The cross-framework appendix therefore exports typed parent surfaces rather
than importing downstream gravitational equations as foundational premises.

The information-area curvature candidate
\[
\Xi_I^{(P)}
=
\frac{\mathcal J_\pi}{a_{FS}^{(P)}}
\left(\frac{\omega_P}{c}\right)^2
\]
is available as a dimensionally closed interface.  RFC owns the downstream
binding of this quantity to its dynamical curvature/action surface.

\section{Software and GREMLIN boundary}

The runtime/solver layer may consume the normalized cycle, phase and relation
surfaces exported by TIR.  Runtime behavior is not a proof parent for the
mathematical theorem graph; instrument or runtime tests enter through explicit
evidence receipts.

GREMLIN is used as a bounded candidate-generation and adversarial-audit layer:
SPIDER may propose dependency isomorphisms, HOUND may search for contradictions,
MOLE may propose derivation routes and MANTIS may detect duplicate ownership.
A GREMLIN candidate enters the v12 theorem graph only after an independent
deterministic proof/validator gate.

\section{Interface status table}

\begin{center}
\begin{tabular}{p{0.32\textwidth}p{0.23\textwidth}p{0.32\textwidth}}
\toprule
Interface & Status & Owner\\
\midrule
\(1/2\leftrightarrow\ln2\) & exact & binary information layer\\
\(3\times8\times\pi\to24\pi\) & TIR structural PASS & Chapter~9\\
\(\Re s=1/2\leftrightarrow|z_L|=1\) & exact & SOH interface\\
off-axis \(\to\) negative Li subsequence & conditional theorem candidate & SOH bridge\\
native closure \(\to\lambda_n\ge0\) & OPEN & SOH completion frontier\\
\(\ell_\varphi=c/|\omega_t|\) typing & exact dimensional closure & IDT interface\\
\(ds_{\rm rel}^2=\ell_\varphi^2ds_{FS}^2\) & scale-binding candidate & TIR--IDT interface\\
Cartan/ADM/Einstein closure & OPEN & Chapter~21 / RFC join\\
\bottomrule
\end{tabular}
\end{center}

The deterministic appendix integration validator checks the exact half-axis
identity, canonical \(\kappa\) ownership markers, the IDT length-scale equation,
the Li-positivity frontier marker and the bounded GREMLIN promotion rule.

% ---- END TIR/monograph/v12/appendices/appD_cross_framework_interfaces.tex ----
\clearpage




% GENERATED LOCAL-BUILD FLAT PART — DO NOT MERGE

% ---- BEGIN TIR/monograph/v12/appendices/appE_v12_1_repository_state_sync.tex ----
% !TEX root = ../../tir_monograph_v12.tex
\chapter{Version 12.1 Repository-State Synchronization}
\label{app:v12_1_repository_sync}

\paragraph{Scope.}
This appendix synchronizes the dependency-ordered v12 publication surface with
TIR theorem and validation work merged after the original v12 migration build.
The v12.0 appendices A--D remain frozen provenance surfaces.  The baseline for
this audit is \texttt{main@0f63c823961a58770de22d14b5f53e2bfa38a9b3}.
No physical verdict is promoted merely because a later mathematical theorem or
software contract exists.

\section{Post-v12 spatial geometry closure}

The original v12 text ended the spatial line at the exact discrete solder and
Universal-Loop source identity.  The current repository contains four further
stages.

\paragraph{Gate A2: Cartan refinement.}
For an admitted regular smooth shrinking family, the discrete solder object and
rotational holonomy have the asymptotic forms
\[
\mathcal T_{\triangle}
=\frac12 T^a{}_{\mu\nu}\Sigma^{\mu\nu}_{\triangle}\sigma_a
+O(\ell^3),
\]
and
\[
R_{\triangle}
=I+\frac12\Omega_{\mu\nu}\Sigma^{\mu\nu}_{\triangle}
+O(\ell^3).
\]
Consequently the area-normalized translational and rotational channels converge
to Cartan torsion and curvature on the declared refinement surface.  This is a
conditional local/refining theorem; existence of the production global
relational complex is not supplied by this statement.
\vTwelveStatus{B}{--}{PASS}

Canonical source:
\path{TIR/foundations/TIR_CARTAN_CONTINUUM_REFINEMENT_V0_1.md}.

\paragraph{Gate A3: torsion-free spatial GR sector.}
When the direct and composed descriptions are admitted as representations of
the same primitive ordered endpoint relation, intrinsic affine-displacement
uniqueness selects
\[
\mathcal C_{xyz}=0.
\]
Together with
\[
\mathcal T_{xyz}=-\mathcal C_{xyz},
\]
this gives the endpoint-compatible discrete sector
\[
\mathcal T_{xyz}=0.
\]
Gate A2 transfers this condition to a regular refinement as
\[
\boxed{T^a=0}.
\]
The induced spatial frame transport is in \(SO(3)\), hence metric compatible;
the standard uniqueness theorem for a torsion-free metric-compatible
connection then selects
\[
\boxed{D=D^{\rm LC}}.
\]
This does not set curvature to zero and does not remove the broader typed
Cartan/torsional branch.
\vTwelveStatus{B}{--}{PASS}

Canonical source:
\path{TIR/foundations/TIR_ZERO_TORSION_LEVI_CIVITA_SELECTION_V0_1.md}.

\paragraph{Gate A4: leading-loop locality.}
For a regular shrinking loop with \(A_C=\Theta(\epsilon^2)\),
\[
\frac{R_C-I}{A_C}\longrightarrow\Omega.
\]
Derivative-of-curvature terms enter at higher refinement order.  Under the TIR
Leading Refinement Rule, the lowest finite nonzero covariant normalized loop
coefficient is selected as the primitive continuum carrier.  On the A3
Levi-Civita sector, this bounds the leading metric operator to
\[
\mathcal E_{\mu\nu}
=\mathcal E_{\mu\nu}(g,\partial g,\partial^2g),
\]
while higher curvature jets remain available as extended/correction sectors.
\vTwelveStatus{B}{--}{PASS}

Canonical source:
\path{TIR/foundations/TIR_LEADING_LOOP_LOCALITY_METRIC_JET_V0_1.md}.

\paragraph{Gate A5: global three-manifold certifier.}
For a supplied finite pure tetrahedral complex, the executable certificate
checks the closed combinatorial three-manifold conditions through face
incidence and spherical vertex links.  Once that certificate passes, standard
three-dimensional topology supplies the compatible smoothing bridge, and the
local \(SO(3)\)-related metrics and Levi-Civita restrictions glue globally.
The theorem/certifier is implemented, but the actual production TIR relational
complex remains an input gate.
\vTwelveStatus{B}{--}{OPEN}

Canonical source:
\path{TIR/foundations/TIR_GLOBAL_3MANIFOLD_SMOOTH_CERTIFICATE_V0_1.md}.

The precise interpretation of the last status is
\[
\boxed{
\text{certifier/theorem PASS}
\quad\text{and}\quad
\text{production relational-complex input OPEN}.
}
\]

\section{Production input contracts}

The previously vague phrase ``global smooth refinement'' has been split into
explicit source-owned data gates.

The global spatial-complex contract defines source capture, deterministic
canonical freeze, provenance hashes and the A5 handoff.  It is executable and
fail closed, while the production incidence dataset remains absent:
\[
\boxed{
\text{GSC-1 contract PASS}
\quad/\quad
\text{production spatial complex OPEN}.
}
\]
Source:
\path{TIR/foundations/TIR_GLOBAL_SPATIAL_COMPLEX_INPUT_CONTRACT_V0_1.md}.

The inter-leaf matching contract similarly fixes the packet carrying
\(\beta_{\rm match}\), its coordinate convention, provenance and digest.  The
handoff consistency condition mirrors the downstream matching law, and the
\(x^0=ct\) conversion exports the normalized ADM-shift packet.  The contract is
implemented; the production matching field is still open:
\[
\boxed{
\text{matching contract PASS}
\quad/\quad
\text{production }\beta_{\rm match}\text{ OPEN}.
}
\]
Source:
\path{TIR/foundations/TIR_INTERLEAF_MATCHING_FIELD_INPUT_CONTRACT_V0_1.md}.

\section{Platonic closure of the L-constants}

The v12 arithmetic provenance is retained, but it is no longer the only
structural path to \((L_3,L_4,L_5)\).

For convex regular polyhedra, the Schlaefli condition
\[
\frac1p+\frac1q>\frac12,
\qquad p,q\ge3,
\]
has exactly the five solutions
\[
\{3,3\},\ \{3,4\},\ \{3,5\},\ \{4,3\},\ \{5,3\}.
\]
The triangular branch has orientation-preserving rotation groups
\[
A_4,\qquad S_4,\qquad A_5,
\]
with orders \(12,24,60\).  Taking the tetrahedral rotational group as the
common root gives the exact subgroup indices
\[
\boxed{[S_4:A_4]=2},
\qquad
\boxed{[A_5:A_4]=5}.
\]
TIR defines the complete non-self-dual extension carrier as
\[
X_3^{\rm closure}
:=S_4/A_4\sqcup A_5/A_4,
\]
so
\[
\boxed{
L_4=2,
\qquad L_5=5,
\qquad L_3=|X_3^{\rm closure}|=2+5=7.
}
\]
Hence
\[
\boxed{(L_3,L_4,L_5)=(7,2,5)}.
\]
The derivation is exact conditional on the explicit TIR root-extension closure
rule.  The finite-group result does not, by itself, establish a physical law or
a derivation from a specific Kaehler geometry.
\vTwelveStatus{B}{--}{PASS}

Canonical source:
\path{TIR/foundations/TIR_PLATONIC_L_CONSTANTS_CLOSURE_V0_1.md}.

Validator:
\path{TIR/validation/tir_platonic_l_constants_closure_v0_1.py}.

The older arithmetic route remains an independent crosscheck:
\[
\operatorname{depth}_{\rm Collatz}(3)=7,
\qquad 5-3=2,
\qquad 5=L_5.
\]

\section{Collatz--Fubini--Study phase interface}

The merged interface imports the conditional discrete projective phase map
\[
q(Cn)=2q(n)\pmod1,
\qquad
\boxed{\zeta_C(Cn)=\zeta_C(n)^2},
\]
and permits an explicit relational coordinate
\[
R_{ij}=R_{ij}(S_i,S_j;\zeta_{ij}),
\qquad \zeta_{ij}=e^{i\phi_{ij}}.
\]
Projective \(2\pi\) phase and an optional spinorial \(4\pi\) lift remain typed
separately.  The CP1 carrier is a mathematical interface.  No validated law yet
identifies this phase itinerary with elapsed time, energy, mass, interaction
strength, atomic transition probability, spectroscopy or gravity.
\vTwelveStatus{B}{--}{OPEN}

Source:
\path{TIR/integration/TIR_COLLATZ_FS_RELATIONAL_PHASE_INTERFACE_V0_1.md}.

\section{Current unresolved physical and dynamical gates}

The current completion frontier is not a flat list of missing formulas.  The
principal unresolved gates are:

\begin{enumerate}
\item a frozen production global spatial-complex incidence input and A5 PASS on
that exact dataset;
\item a frozen production inter-leaf matching field and the global TIR--IDT--RFC
spacetime/ADM join;
\item Einstein constraint/evolution closure downstream of the admitted global
spacetime carrier and normalization/action gates;
\item target-independent extraction of the four coefficient magnitudes
\(|h|,|a|,|b|,|c|\);
\item continuum gauge normalization and running from the typed \(W_{ij}\)
transport carrier;
\item a common electroweak renormalization/scheme-scale transport;
\item scalar/Higgs action binding;
\item a source/uniqueness theorem for hypercharge and a scheme/scale-defined
quark-mass map;
\item a repaired neutrino absolute-action derivation;
\item a source-derived meson absolute-action baseline;
\item a holonomic/topological strong-CP source theorem before any replacement
of the failed frozen neutron-EDM route;
\item a unit-complete cosmological scale/critical-density binding;
\item a physical binding for the Collatz--FS relational phase coordinate;
\item native Li/Weil positivity and the remaining global critical-axis
nondegeneracy/strict-positivity premises.
\end{enumerate}

\section{Retained empirical failures and tensions}

This synchronization does not rescore the observable matrix.  In particular,
the current publication retains charged-lepton precision failures, the PMNS
reactor-angle tension, electroweak/fine-structure/Higgs precision failures or
tensions according to Chapter~19, the pion and kaon printed-formula failures,
and the neutron-EDM failure from the frozen strong-CP mapping.

For the latter,
\[
d_n\simeq5.3299\times10^{-26}\,e\,\mathrm{cm}
\]
remains approximately \(2.96\) times the manuscript bound used by the frozen
comparison.  A new structural derivation must receive a new version and evidence
record; it does not overwrite this failure.

\section{Critical-axis firewall}

The active critical-axis stack contains exact identities, equivalence
reductions and conditional positive corridors, but the globally quantified
strict-positivity/nondegeneracy premises remain open.  The integrated solver
contract explicitly requires that the Riemann hypothesis is not present in the
current closure.  Thus
\[
\boxed{\mathrm{RH}\ \text{remains OPEN}.}
\]
Failed candidate constructions and numerical negative controls remain evidence
rather than being removed from the project record.

\section{v12.1 synchronization invariant}

The publication rule for this revision is
\[
\boxed{
\text{merged theorem/contract PASS}
\not\Rightarrow
\text{physical observable PASS}.
}
\]
Every exact or conditional theorem keeps its assumptions.  Production-input
gates remain open until an exact frozen source dataset is supplied.  Existing
empirical failures remain visible.  Main promotion of this corrected monograph
is a separate repository action.

% ---- END TIR/monograph/v12/appendices/appE_v12_1_repository_state_sync.tex ----
\clearpage

% ---- BEGIN TIR/monograph/v12/appendices/appF_gremlin_pass4_repository_consolidation.tex ----
% !TEX root = ../../tir_monograph_v12.tex
\chapter{GREMLIN Pass-4 Repository Consolidation}
\label{app:v12_gremlin_pass4}

\paragraph{Purpose.}
Version 12.2 records a bounded repository-consolidation pass performed under the
GREMLIN candidate/audit runtime.  The pass does not create scientific promotion
authority.  Its function is to reduce branch ambiguity, preserve unique work,
synchronize the publication surface and keep the claim-status firewall intact.

\section{Branch sweep}

The baseline repository state was
\[
\boxed{\texttt{main@45b54c0c8aa60ec92607e6132b69aca9ab766d75}.}
\]
Exact comparison of all 81 non-main branches against that baseline gave
\[
\boxed{
81=77_{\rm contained}+4_{\rm unique}.
}
\]
The 77 contained branches had \texttt{ahead\_by=0}; their payload was already
reachable from main.  No branch deletion is performed by this publication
pass.

The four unique-payload sources were:
\begin{enumerate}
\item \path{feat/golden-mean-prime-orbit-asymptotic-v0.1};
\item \path{jarlskog-closure-v01};
\item \path{policy/no-ai-training-without-license-20260917};
\item \path{research/xf9-preprint-spider-bib-20260918}.
\end{enumerate}
The complete branch inventory is retained in
\path{TIR/monograph/v12/GREMLIN_PASS4_BRANCH_MONOGRAPH_AUDIT_20260918.md}.

\section{Integration rule}

Clean or additive unique files are preserved byte-for-byte.  The Jarlskog branch
was substantially behind the current repository state, so its old copies of
current monograph/workflow files were not merged wholesale.  Instead, unique
theorem/validator surfaces were preserved and the publication changes were
transplanted onto the current v12 surface.

This rule prevents a stale branch from reverting later theorem closures while
still preserving every unique research artifact.

\section{Scientific effects}

Pass-4 changes the canonical narrative only where the repository state changed
the mathematics:
\begin{itemize}
\item the Jarlskog invariant is exact downstream of the retained CKM matrix
      parameters and the former independent \(J_{\rm old}\) assignment is
      deprecated;
\item the CKM 600-cell/McKay/A2/half-kernel chain is preserved as an
      exact/conditional selector programme, with physical readout uniqueness
      still open;
\item the golden-mean subsystem gains an exact primitive dynamical-orbit
      asymptotic with explicit error bounds and an arithmetic-prime/RH
      firewall;
\item XF-9 gains the closure-force, pole-free Gram-kernel and de Branges
      crosswalk together with explicit prior/concurrent-art provenance.
\end{itemize}

No retained empirical FAIL, TENSION, OPEN or QUARANTINED verdict is overwritten
by these structural updates.

\section{GREMLIN authority boundary}

The live GREMLIN runtime for this pass reports
\[
\boxed{
\texttt{authority=CANDIDATE\_ONLY},
\qquad
\texttt{canon\_allowed=false}.
}
\]
GREMLIN routes candidate searches, dependency checks, adversarial audits and
redundancy analysis.  Canonical repository changes still require explicit
source edits, deterministic validators, publication preflight and normal
repository promotion.

\section{Destructive-cleanup boundary}

This appendix certifies branch classification and payload preservation only.
Remote-branch deletion is a separate destructive action and requires explicit
authorization plus a fresh post-integration containment check.

% ---- END TIR/monograph/v12/appendices/appF_gremlin_pass4_repository_consolidation.tex ----
\clearpage

% ---- BEGIN TIR/monograph/v12/appendices/appG_semantic_frontier_reconciliation.tex ----
% !TEX root = ../../tir_monograph_v12.tex
\chapter{Semantic Frontier Reconciliation}
\label{app:v12_semantic_frontier_reconciliation}

\paragraph{Purpose.}
This appendix records the v12.3 correction of the completion frontier after a
cross-repository theorem/validator/receipt audit. The previous v12.1 frontier
is retained as historical provenance but is not used as the current semantic
source of truth.

\section{Status model}

Every frontier coordinate is split into
\[
\boxed{
\text{formal/derivational status}
\quad\oplus\quad
\text{production/physical status}.
}
\]
This prevents two opposite errors: a software PASS cannot create a physical
PASS, and an open physical source input cannot relabel a completed derivation
as mathematically open.

\section{Corrected former headline gates}

\begin{longtable}{p{0.07\textwidth}p{0.36\textwidth}p{0.47\textwidth}}
\toprule
No. & Former headline & v12.3 status \\
\midrule
1 & Production global spatial complex &
A5/capture/binding certifiers closed; source-owned production realization open. \\
2 & Inter-leaf matching and global spacetime/ADM join &
Matching-flow, shared-atlas and globalization certifiers closed conditionally;
production packet, event placement and coverage open. \\
3 & Einstein constraint/evolution closure &
Local ADM constraints, evolution, Bianchi propagation and local Einstein form
closed; production global carrier/coverage remains open. \\
4 & Four coefficient magnitudes &
Magnitude-parent evaluation closed; transition-parent selector open. \\
5 & Continuum gauge normalization/running &
Discrete structural gauge layer closed; continuum normalization/running open. \\
6 & Electroweak scheme-scale transport &
Open. \\
7 & Higgs scalar/action binding &
Open. \\
8 & Hypercharge uniqueness and quark masses &
Hypercharge relative uniqueness closed on declared field content and anchor;
quark-mass map and stronger geometric-anchor uniqueness open. \\
9 & Neutrino absolute-action repair &
Source conflict and derivational repair closed; empirical absolute-mass scale
open. \\
10 & Meson absolute-action baseline &
Open. \\
11 & Strong-CP source theorem &
Open; retained neutron-EDM physical FAIL remains visible. \\
12 & Cosmological scale/critical density &
Dimensionless arithmetic audited; dimensionful and rho-critical binding open. \\
13 & Collatz--FS physical binding &
Kinematic phase-fibre interface closed; metric/source coupling and scale open. \\
14 & Li--Weil/global critical-axis positivity &
Open; RH remains open. \\
\bottomrule
\end{longtable}

\section{Pinned repositories}

The reconciliation pins:
\begin{itemize}
\item TIR main: \texttt{45b54c0c8aa60ec92607e6132b69aca9ab766d75};
\item IDT main: \texttt{58f453a4725bf0304417a5d07730f2dc1765dea5};
\item RFC main: \texttt{ce4af9ddd480ea40dfb5a3ee26ed396a58cb0e54};
\item Secret-of-a-Half main:
      \texttt{d99545aa447ef86bc253d8241a3fee9adb8a42c1}.
\end{itemize}

The machine-readable ledger is
\path{TIR/monograph/v12/SEMANTIC_FRONTIER_V12_3.json}.
The human-readable provenance audit is
\path{TIR/monograph/v12/GREMLIN_PASS5_SEMANTIC_FRONTIER_AUDIT_20260918.md}.

\section{Promotion boundary}

GREMLIN remains candidate/audit authority only. This reconciliation changes the
status description only where an existing theorem, validator or receipt already
supports the change. It does not create new physical evidence and does not
authorize a merge to main or destructive branch cleanup.

% ---- END TIR/monograph/v12/appendices/appG_semantic_frontier_reconciliation.tex ----
\clearpage

% ---- BEGIN TIR/monograph/v12/appendices/appH_relational_zero_graph_reconciliation.tex ----
% !TEX root = ../../tir_monograph_v12.tex
\chapter{Relational Zero and Current-Graph Reconciliation}
\label{app:v12_relational_zero_graph_reconciliation}

\section{Purpose}

Version 12.4 adds two corrections that must be kept separate:

\begin{enumerate}
\item the primitive zero layer is made explicitly relational and pre-object;
\item the publication dependency map is reconciled against repository state after the
18 September 2026 v12.3 semantic-frontier snapshot.
\end{enumerate}

The v12.3 semantic-frontier files remain historical provenance. They are not deleted or
retroactively rewritten.

\section{Relational-zero root}

TIR now begins from the relational presentation
\[
\mathfrak R=(X,\mathcal D),
\]
with support set $X$ and realized distinction/relation set $\mathcal D$. The zero layer is
\[
\boxed{
\mathfrak Z_{\rm rel}=(\varnothing,\varnothing).
}
\]
It is not a singleton and not an object. Once TIR defines existential admissibility
relationally, the absence of an existent object in this empty presentation is a direct
definitional consequence.

Leaving zero first supplies the minimal nonempty support
\[
\boxed{\mathcal P=\{p\}},
\]
and only then a minimal nontrivial distinction. Since a nontrivial partition contains
more than one nonempty block, the minimal one is binary:
\[
\boxed{\Pi_1=\{N,S\}}.
\]
Exchange symmetry is imposed after binarity has been established,
\[
J:N\leftrightarrow S,\qquad J^2=\mathrm{id},
\]
and normalized exchange invariance then gives
\[
\boxed{(w_N,w_S)=\left(\frac12,\frac12\right)},
\qquad
\boxed{H_2(1/2)=\ln2}.
\]

The corrected primitive dependency order is therefore
\[
\boxed{
0_{\rm rel}
\to
\text{point support}
\to
\text{minimal binary distinction}
\to
\text{exchange symmetry}
\to
\frac12
\to
\ln2.
}
\]

The quantum lift remains downstream and conditional on the TIR quantum-point postulate:
\[
\mathbb C^2\to\mathbb{CP}^1\cong S^2.
\]
No quantum, spacetime or empirical claim is inferred from relational zero alone.

\section{Reconciliation base heads}

The v12.4 reconciliation began from the following repository heads:

\begin{description}
\item[TIR:] \path{853032e019824b836de43634d949da5566396153}.
\item[IDT:] \path{e474f64a079e0d90a039e0abd19d2b482f7e53d2}.
\item[RFC:] \path{664349e2f16b57ee23bf6f5f21f5fe8ada17c5a5}.
\item[Secret-of-a-Half:] \path{492dc2e31546f6fdce798801549d544d8cbda3b4}.
\end{description}

The live machine-readable authority for TIR dependencies is the repository-root
\path{DEPENDENCY_EXPORT.json}. Appendix~\ref{app:v12_semantic_frontier_reconciliation}
and \path{SEMANTIC_FRONTIER_V12_3.json} retain the 18 September state only.

\section{Post-v12.3 dependency delta}

The current repository state contains theorem or candidate surfaces that did not exist in
the v12.3 publication snapshot. Their presence does not promote their physical interpretation.

\subsection{TIR gravity and event-source interfaces}

The post-v12.3 graph includes the fractal/orbital gravity kinematic controls, conditional
flow-coframe/ADM spherical-vacuum derivations, event-indexed metric-rate source bridge,
information-scalar acceleration criterion, temporal-$U(1)$ metric-response candidate, and
phase-cell vacuum-integration route.

The publication firewall is:

\[
\boxed{
\text{exact or conditional mathematical control}
\not\Rightarrow
\text{production physical binding}.
}
\]

Production event--spatial realization, absolute information normalization, nontrivial
holonomy persistence, and the relevant source-to-response identifications remain separately
typed open inputs where the repository says they are open.

\subsection{IDT clock and representation additions}

The companion IDT graph additionally contains, after the v12.3 pin, the periodic
$6\pi/C_3/F_3$/Pauli representation bridge, production event-complex and GSC2 source
contracts, temporal-trace/clock-calibration candidates, the EB/BEC conditional bridge,
and Lean proof gates for selected discrete statements.

These are cross-repository dependencies or candidates; they do not silently become TIR
physical claims.

\section{Graph/publication invariant}

The publication invariant for v12.4 is:

\[
\boxed{
\text{monograph statement}
\longleftrightarrow
\text{typed graph node/edge}
\longleftrightarrow
\text{named source surface}.
}
\]

If a graph export and prose snapshot disagree in date, the newer explicitly generated graph
is the dependency authority and the older prose remains historical provenance. If the graph
marks an edge \texttt{CANDIDATE\_ONLY}, the monograph must not describe that edge as a
canonical physical derivation.

\section{Validation boundary}

The deterministic relational-zero validator is

\path{TIR/validation/tir_relational_zero_axiom_v0_1.py}.

It checks the empty zero support, minimal nonempty support cardinality, minimal binary
distinction, exchange-fixed half, and the Shannon value $\ln2$. It does not validate a
physical ontology, a quantum realization, spacetime, or experiment.

% ---- END TIR/monograph/v12/appendices/appH_relational_zero_graph_reconciliation.tex ----
\clearpage




\markboth{References}{References}
\pagestyle{headings}

% GENERATED LOCAL-BUILD FLAT PART — DO NOT MERGE

% ---- BEGIN TIR/monograph/references_expanded_v10_8.tex ----
% Expanded bibliography and source map for the Metatime monograph.
% Version v10.8, 2026-07-29.
% The grouping distinguishes external scientific context from original TIR claims.

\renewcommand{\bibname}{Bibliography and Source Map}
\begin{thebibliography}{999}
\small
\sloppy

\item[]\textit{Editorial note.} The references below document the established
mathematical, physical, experimental, and methodological context used throughout
this monograph. Inclusion of a source does not imply that the source endorses the
Metatime/TIR construction, nor that the original claims of this monograph follow
from the cited literature. Numerical particle data should be read against the
current Particle Data Group listings and their published updates.

\item[]\textbf{A. Reference data, precision measurements, and major experiments}

\bibitem{PDG2024}
S.~Navas \emph{et al.} (Particle Data Group),
``Review of Particle Physics,''
Phys.\ Rev.\ D \textbf{110}, 030001 (2024), with the 2025 online update.

\bibitem{CODATA2022}
E.~Tiesinga, P.~J.~Mohr, D.~B.~Newell, and B.~N.~Taylor,
``The 2022 CODATA Recommended Values of the Fundamental Physical Constants,''
NIST Standard Reference Database 121, Web Version 9.0 (2024),
\url{https://physics.nist.gov/constants}.

\bibitem{Planck2018}
N.~Aghanim \emph{et al.} (Planck Collaboration),
``Planck 2018 results. VI. Cosmological parameters,''
Astron.\ Astrophys.\ \textbf{641}, A6 (2020).

\bibitem{DESI2024}
A.~G.~Adame \emph{et al.} (DESI Collaboration),
``DESI 2024 VI: Cosmological Constraints from the Measurements of Baryon
Acoustic Oscillations,'' arXiv:2404.03002 (2024).

\bibitem{ATLAS2012Higgs}
G.~Aad \emph{et al.} (ATLAS Collaboration),
``Observation of a new particle in the search for the Standard Model Higgs
boson with the ATLAS detector at the LHC,''
Phys.\ Lett.\ B \textbf{716}, 1--29 (2012).

\bibitem{CMS2012Higgs}
S.~Chatrchyan \emph{et al.} (CMS Collaboration),
``Observation of a new boson at a mass of 125 GeV with the CMS experiment at
the LHC,'' Phys.\ Lett.\ B \textbf{716}, 30--61 (2012).

\bibitem{ATLAS2022Charm}
G.~Aad \emph{et al.} (ATLAS Collaboration),
``Direct constraint on the Higgs--charm coupling from a search for Higgs boson
decays into charm quarks with the ATLAS detector,''
Eur.\ Phys.\ J.\ C \textbf{82}, 717 (2022).

\bibitem{SuperK1998}
Y.~Fukuda \emph{et al.} (Super-Kamiokande Collaboration),
``Evidence for oscillation of atmospheric neutrinos,''
Phys.\ Rev.\ Lett.\ \textbf{81}, 1562--1567 (1998).

\bibitem{SNO2002}
Q.~R.~Ahmad \emph{et al.} (SNO Collaboration),
``Direct evidence for neutrino flavor transformation from neutral-current
interactions in the Sudbury Neutrino Observatory,''
Phys.\ Rev.\ Lett.\ \textbf{89}, 011301 (2002).

\bibitem{KamLAND2003}
K.~Eguchi \emph{et al.} (KamLAND Collaboration),
``First results from KamLAND: Evidence for reactor antineutrino disappearance,''
Phys.\ Rev.\ Lett.\ \textbf{90}, 021802 (2003).

\bibitem{DayaBay2012}
F.~P.~An \emph{et al.} (Daya Bay Collaboration),
``Observation of electron-antineutrino disappearance at Daya Bay,''
Phys.\ Rev.\ Lett.\ \textbf{108}, 171803 (2012).

\bibitem{Abel2020nEDM}
C.~Abel \emph{et al.},
``Measurement of the permanent electric dipole moment of the neutron,''
Phys.\ Rev.\ Lett.\ \textbf{124}, 081803 (2020).

\bibitem{MuonG2Final2025}
D.~P.~Aguillard \emph{et al.} (Muon $g-2$ Collaboration),
``Measurement of the positive muon anomalous magnetic moment to 127 ppb,''
Phys.\ Rev.\ Lett.\ \textbf{135} (2025), DOI: 10.1103/7clf-sm2v.

\bibitem{Riess1998}
A.~G.~Riess \emph{et al.},
``Observational evidence from supernovae for an accelerating universe and a
cosmological constant,'' Astron.\ J.\ \textbf{116}, 1009--1038 (1998).

\bibitem{Perlmutter1999}
S.~Perlmutter \emph{et al.},
``Measurements of $\Omega$ and $\Lambda$ from 42 high-redshift supernovae,''
Astrophys.\ J.\ \textbf{517}, 565--586 (1999).

\item[]\textbf{B. Quantum mechanics, quantum field theory, and the Standard Model}

\bibitem{Noether1918}
E.~Noether,
``Invariante Variationsprobleme,''
Nachr.\ Ges.\ Wiss.\ G\"ottingen, Math.-Phys.\ Kl. 235--257 (1918).

\bibitem{Dirac1928}
P.~A.~M.~Dirac,
``The quantum theory of the electron,''
Proc.\ R.\ Soc.\ Lond.\ A \textbf{117}, 610--624 (1928); \textbf{118},
351--361 (1928).

\bibitem{Weyl1929}
H.~Weyl,
``Elektron und Gravitation. I,''
Z.\ Phys.\ \textbf{56}, 330--352 (1929).

\bibitem{Tomonaga1946}
S.~Tomonaga,
``On a relativistically invariant formulation of the quantum theory of wave
fields,'' Prog.\ Theor.\ Phys.\ \textbf{1}, 27--42 (1946).

\bibitem{Schwinger1948}
J.~Schwinger,
``Quantum electrodynamics. I. A covariant formulation,''
Phys.\ Rev.\ \textbf{74}, 1439--1461 (1948).

\bibitem{Feynman1949}
R.~P.~Feynman,
``Space-time approach to quantum electrodynamics,''
Phys.\ Rev.\ \textbf{76}, 769--789 (1949).

\bibitem{Dyson1949}
F.~J.~Dyson,
``The radiation theories of Tomonaga, Schwinger, and Feynman,''
Phys.\ Rev.\ \textbf{75}, 486--502 (1949).

\bibitem{YangMills1954}
C.~N.~Yang and R.~L.~Mills,
``Conservation of isotopic spin and isotopic gauge invariance,''
Phys.\ Rev.\ \textbf{96}, 191--195 (1954).

\bibitem{Glashow1961}
S.~L.~Glashow,
``Partial-symmetries of weak interactions,''
Nucl.\ Phys.\ \textbf{22}, 579--588 (1961).

\bibitem{Weinberg1967}
S.~Weinberg,
``A model of leptons,'' Phys.\ Rev.\ Lett.\ \textbf{19}, 1264--1266 (1967).

\bibitem{Salam1968}
A.~Salam,
``Weak and electromagnetic interactions,'' in
\emph{Elementary Particle Theory}, edited by N.~Svartholm
(Almqvist and Wiksell, Stockholm, 1968), pp.~367--377.

\bibitem{EnglertBrout1964}
F.~Englert and R.~Brout,
``Broken symmetry and the mass of gauge vector mesons,''
Phys.\ Rev.\ Lett.\ \textbf{13}, 321--323 (1964).

\bibitem{Higgs1964}
P.~W.~Higgs,
``Broken symmetries and the masses of gauge bosons,''
Phys.\ Rev.\ Lett.\ \textbf{13}, 508--509 (1964).

\bibitem{GHK1964}
G.~S.~Guralnik, C.~R.~Hagen, and T.~W.~B.~Kibble,
``Global conservation laws and massless particles,''
Phys.\ Rev.\ Lett.\ \textbf{13}, 585--587 (1964).

\bibitem{tHooftVeltman1972}
G.~'t~Hooft and M.~Veltman,
``Regularization and renormalization of gauge fields,''
Nucl.\ Phys.\ B \textbf{44}, 189--213 (1972).

\bibitem{GrossWilczek1973}
D.~J.~Gross and F.~Wilczek,
``Ultraviolet behavior of non-Abelian gauge theories,''
Phys.\ Rev.\ Lett.\ \textbf{30}, 1343--1346 (1973).

\bibitem{Politzer1973}
H.~D.~Politzer,
``Reliable perturbative results for strong interactions?''
Phys.\ Rev.\ Lett.\ \textbf{30}, 1346--1349 (1973).

\bibitem{Wilson1974}
K.~G.~Wilson,
``Confinement of quarks,'' Phys.\ Rev.\ D \textbf{10}, 2445--2459 (1974).

\bibitem{PeskinSchroeder1995}
M.~E.~Peskin and D.~V.~Schroeder,
\emph{An Introduction to Quantum Field Theory}
(Addison-Wesley, Reading, 1995).

\bibitem{WeinbergQFT1995}
S.~Weinberg,
\emph{The Quantum Theory of Fields, Vols. I--III}
(Cambridge University Press, Cambridge, 1995--2000).

\item[]\textbf{C. Flavor, hadrons, neutrinos, and symmetry classification}

\bibitem{gellmann1961}
M.~Gell-Mann,
``The eightfold way: A theory of strong interaction symmetry,''
Caltech Synchrotron Laboratory Report CTSL-20 (1961).

\bibitem{Neeman1961}
Y.~Ne'eman,
``Derivation of strong interactions from a gauge invariance,''
Nucl.\ Phys.\ \textbf{26}, 222--229 (1961).

\bibitem{GellMann1962}
M.~Gell-Mann,
``Symmetries of baryons and mesons,''
Phys.\ Rev.\ \textbf{125}, 1067--1084 (1962).

\bibitem{okubo1962}
S.~Okubo,
``Note on unitary symmetry in strong interactions,''
Prog.\ Theor.\ Phys.\ \textbf{27}, 949--966 (1962).

\bibitem{Cabibbo1963}
N.~Cabibbo,
``Unitary symmetry and leptonic decays,''
Phys.\ Rev.\ Lett.\ \textbf{10}, 531--533 (1963).

\bibitem{GellMann1964}
M.~Gell-Mann,
``A schematic model of baryons and mesons,''
Phys.\ Lett.\ \textbf{8}, 214--215 (1964).

\bibitem{Zweig1964}
G.~Zweig,
``An SU(3) model for strong interaction symmetry and its breaking,''
CERN Reports TH-401 and TH-412 (1964).

\bibitem{kobayashi1973}
M.~Kobayashi and T.~Maskawa,
``CP-violation in the renormalizable theory of weak interaction,''
Prog.\ Theor.\ Phys.\ \textbf{49}, 652--657 (1973).

\bibitem{Wolfenstein1983}
L.~Wolfenstein,
``Parametrization of the Kobayashi--Maskawa matrix,''
Phys.\ Rev.\ Lett.\ \textbf{51}, 1945--1947 (1983).

\bibitem{Jarlskog1985}
C.~Jarlskog,
``Commutator of the quark mass matrices in the Standard Electroweak Model and
a measure of maximal CP nonconservation,''
Phys.\ Rev.\ Lett.\ \textbf{55}, 1039--1042 (1985).

\bibitem{Maki1962}
Z.~Maki, M.~Nakagawa, and S.~Sakata,
``Remarks on the unified model of elementary particles,''
Prog.\ Theor.\ Phys.\ \textbf{28}, 870--880 (1962).

\bibitem{Pontecorvo1957}
B.~Pontecorvo,
``Mesonium and antimesonium,''
Zh.\ Eksp.\ Teor.\ Fiz.\ \textbf{33}, 549--551 (1957)
[Sov.\ Phys.\ JETP \textbf{6}, 429 (1958)].

\bibitem{Wolfenstein1978}
L.~Wolfenstein,
``Neutrino oscillations in matter,''
Phys.\ Rev.\ D \textbf{17}, 2369--2374 (1978).

\bibitem{MikheyevSmirnov1985}
S.~P.~Mikheyev and A.~Y.~Smirnov,
``Resonance amplification of oscillations in matter and spectroscopy of solar
neutrinos,'' Sov.\ J.\ Nucl.\ Phys.\ \textbf{42}, 913--917 (1985).

\item[]\textbf{D. Geometric phase, differential geometry, topology, and holonomy}

\bibitem{Pancharatnam1956}
S.~Pancharatnam,
``Generalized theory of interference and its applications,''
Proc.\ Indian Acad.\ Sci.\ A \textbf{44}, 247--262 (1956).

\bibitem{Simon1983}
B.~Simon,
``Holonomy, the quantum adiabatic theorem, and Berry's phase,''
Phys.\ Rev.\ Lett.\ \textbf{51}, 2167--2170 (1983).

\bibitem{berry1984}
M.~V.~Berry,
``Quantal phase factors accompanying adiabatic changes,''
Proc.\ R.\ Soc.\ Lond.\ A \textbf{392}, 45--57 (1984).

\bibitem{WilczekZee1984}
F.~Wilczek and A.~Zee,
``Appearance of gauge structure in simple dynamical systems,''
Phys.\ Rev.\ Lett.\ \textbf{52}, 2111--2114 (1984).

\bibitem{AharonovAnandan1987}
Y.~Aharonov and J.~Anandan,
``Phase change during a cyclic quantum evolution,''
Phys.\ Rev.\ Lett.\ \textbf{58}, 1593--1596 (1987).

\bibitem{AnandanAharonov1990}
J.~Anandan and Y.~Aharonov,
``Geometry of quantum evolution,''
Phys.\ Rev.\ Lett.\ \textbf{65}, 1697--1700 (1990).

\bibitem{ProvostVallee1980}
J.~P.~Provost and G.~Vall\'ee,
``Riemannian structure on manifolds of quantum states,''
Commun.\ Math.\ Phys.\ \textbf{76}, 289--301 (1980).

\bibitem{Kahler1933}
E.~K\"ahler,
``\"Uber eine bemerkenswerte Hermitesche Metrik,''
Abh.\ Math.\ Sem.\ Univ.\ Hamburg \textbf{9}, 173--186 (1933).

\bibitem{Hopf1931}
H.~Hopf,
``\"Uber die Abbildungen der dreidimensionalen Sph\"are auf die Kugelfl\"ache,''
Math.\ Ann.\ \textbf{104}, 637--665 (1931).

\bibitem{Chern1946}
S.-S.~Chern,
``Characteristic classes of Hermitian manifolds,''
Ann.\ Math.\ \textbf{47}, 85--121 (1946).

\bibitem{AmbroseSinger1953}
W.~Ambrose and I.~M.~Singer,
``A theorem on holonomy,''
Trans.\ Am.\ Math.\ Soc.\ \textbf{75}, 428--443 (1953).

\bibitem{pontryagin1955}
L.~S.~Pontryagin,
``Smooth manifolds and their applications in homotopy theory,''
Trudy Mat.\ Inst.\ Steklov \textbf{45}, 1--139 (1955).

\bibitem{KobayashiNomizu1963}
S.~Kobayashi and K.~Nomizu,
\emph{Foundations of Differential Geometry}, Vols.~I--II
(Wiley-Interscience, New York, 1963--1969).

\bibitem{BottTu1982}
R.~Bott and L.~W.~Tu,
\emph{Differential Forms in Algebraic Topology}
(Springer, New York, 1982).

\bibitem{Nakahara2003}
M.~Nakahara,
\emph{Geometry, Topology and Physics}, 2nd ed.
(Institute of Physics Publishing, Bristol, 2003).

\bibitem{BengtssonZyczkowski2006}
I.~Bengtsson and K.~\.{Z}yczkowski,
\emph{Geometry of Quantum States}
(Cambridge University Press, Cambridge, 2006).

\item[]\textbf{E. Information theory, statistical geometry, and computation}

\bibitem{shannon1948}
C.~E.~Shannon,
``A mathematical theory of communication,''
Bell Syst.\ Tech.\ J.\ \textbf{27}, 379--423 and 623--656 (1948).

\bibitem{vonNeumann1932}
J.~von Neumann,
\emph{Mathematische Grundlagen der Quantenmechanik}
(Springer, Berlin, 1932).

\bibitem{Fisher1925}
R.~A.~Fisher,
``Theory of statistical estimation,''
Proc.\ Camb.\ Philos.\ Soc.\ \textbf{22}, 700--725 (1925).

\bibitem{Rao1945}
C.~R.~Rao,
``Information and the accuracy attainable in the estimation of statistical
parameters,'' Bull.\ Calcutta Math.\ Soc.\ \textbf{37}, 81--91 (1945).

\bibitem{Jaynes1957}
E.~T.~Jaynes,
``Information theory and statistical mechanics,''
Phys.\ Rev.\ \textbf{106}, 620--630 (1957); \textbf{108}, 171--190 (1957).

\bibitem{Landauer1961}
R.~Landauer,
``Irreversibility and heat generation in the computing process,''
IBM J.\ Res.\ Dev.\ \textbf{5}, 183--191 (1961).

\bibitem{Bures1969}
D.~Bures,
``An extension of Kakutani's theorem on infinite product measures to the
tensor product of semifinite $W^*$-algebras,''
Trans.\ Am.\ Math.\ Soc.\ \textbf{135}, 199--212 (1969).

\bibitem{Wootters1981}
W.~K.~Wootters,
``Statistical distance and Hilbert space,''
Phys.\ Rev.\ D \textbf{23}, 357--362 (1981).

\bibitem{Bennett1982}
C.~H.~Bennett,
``The thermodynamics of computation---a review,''
Int.\ J.\ Theor.\ Phys.\ \textbf{21}, 905--940 (1982).

\bibitem{BraunsteinCaves1994}
S.~L.~Braunstein and C.~M.~Caves,
``Statistical distance and the geometry of quantum states,''
Phys.\ Rev.\ Lett.\ \textbf{72}, 3439--3443 (1994).

\bibitem{AmariNagaoka2000}
S.-I.~Amari and H.~Nagaoka,
\emph{Methods of Information Geometry}
(American Mathematical Society, Providence, 2000).

\bibitem{NielsenChuang2000}
M.~A.~Nielsen and I.~L.~Chuang,
\emph{Quantum Computation and Quantum Information}
(Cambridge University Press, Cambridge, 2000).

\bibitem{Lloyd2000}
S.~Lloyd,
``Ultimate physical limits to computation,''
Nature \textbf{406}, 1047--1054 (2000).

\bibitem{CoverThomas2006}
T.~M.~Cover and J.~A.~Thomas,
\emph{Elements of Information Theory}, 2nd ed.
(Wiley, Hoboken, 2006).

\item[]\textbf{F. Synchronization, nonlinear dynamics, and collective phase models}

\bibitem{Kuramoto1975}
Y.~Kuramoto,
``Self-entrainment of a population of coupled nonlinear oscillators,'' in
\emph{International Symposium on Mathematical Problems in Theoretical Physics},
Lecture Notes in Physics Vol.~39 (Springer, Berlin, 1975), pp.~420--422.

\bibitem{Feigenbaum1978}
M.~J.~Feigenbaum,
``Quantitative universality for a class of nonlinear transformations,''
J.\ Stat.\ Phys.\ \textbf{19}, 25--52 (1978).

\bibitem{Strogatz2000}
S.~H.~Strogatz,
``From Kuramoto to Crawford: Exploring the onset of synchronization in
populations of coupled oscillators,''
Physica D \textbf{143}, 1--20 (2000).

\bibitem{Acebron2005}
J.~A.~Acebr\'on, L.~L.~Bonilla, C.~J.~P.~Vicente, F.~Ritort, and R.~Spigler,
``The Kuramoto model: A simple paradigm for synchronization phenomena,''
Rev.\ Mod.\ Phys.\ \textbf{77}, 137--185 (2005).

\bibitem{OttAntonsen2008}
E.~Ott and T.~M.~Antonsen,
``Low dimensional behavior of large systems of globally coupled oscillators,''
Chaos \textbf{18}, 037113 (2008).

\bibitem{Strogatz2015}
S.~H.~Strogatz,
\emph{Nonlinear Dynamics and Chaos}, 2nd ed.
(Westview Press, Boulder, 2015).

\item[]\textbf{G. Collatz dynamics, Ramanujan analysis, primes, and zeta structure}

\bibitem{Terras1976}
R.~Terras,
``A stopping time problem on the positive integers,''
Acta Arith.\ \textbf{30}, 241--252 (1976).

\bibitem{Lagarias1985}
J.~C.~Lagarias,
``The $3x+1$ problem and its generalizations,''
Am.\ Math.\ Mon.\ \textbf{92}, 3--23 (1985).

\bibitem{Wirsching1998}
G.~J.~Wirsching,
\emph{The Dynamical System Generated by the $3n+1$ Function}
(Springer, Berlin, 1998).

\bibitem{Tao2022Collatz}
T.~Tao,
``Almost all orbits of the Collatz map attain almost bounded values,''
Forum Math.\ Pi \textbf{10}, e12 (2022).

\bibitem{Ramanujan1918}
S.~Ramanujan,
``On certain trigonometrical sums and their applications in the theory of
numbers,'' Trans.\ Cambridge Philos.\ Soc.\ \textbf{22}, 259--276 (1918).

\bibitem{HardyRamanujan1918}
G.~H.~Hardy and S.~Ramanujan,
``Asymptotic formulae in combinatory analysis,''
Proc.\ Lond.\ Math.\ Soc.\ \textbf{17}, 75--115 (1918).

\bibitem{Rademacher1937}
H.~Rademacher,
``On the partition function $p(n)$,''
Proc.\ Natl.\ Acad.\ Sci.\ USA \textbf{23}, 78--84 (1937).

\bibitem{HardyLittlewood1923}
G.~H.~Hardy and J.~E.~Littlewood,
``Some problems of `Partitio Numerorum'; III: On the expression of a number as
a sum of primes,'' Acta Math.\ \textbf{44}, 1--70 (1923).

\bibitem{Riemann1859}
B.~Riemann,
``\"Uber die Anzahl der Primzahlen unter einer gegebenen Gr\"osse,''
Monatsber.\ K\"onigl.\ Preuss.\ Akad.\ Wiss.\ Berlin 671--680 (1859).

\bibitem{Selberg1956}
A.~Selberg,
``Harmonic analysis and discontinuous groups in weakly symmetric Riemannian
spaces with applications to Dirichlet series,''
J.\ Indian Math.\ Soc.\ \textbf{20}, 47--87 (1956).

\bibitem{BerryKeating1999}
M.~V.~Berry and J.~P.~Keating,
``The Riemann zeros and eigenvalue asymptotics,''
SIAM Rev.\ \textbf{41}, 236--266 (1999).

\bibitem{Connes1999}
A.~Connes,
``Trace formula in noncommutative geometry and the zeros of the Riemann zeta
function,'' Selecta Math.\ \textbf{5}, 29--106 (1999).

\bibitem{Zhang2014}
Y.~Zhang,
``Bounded gaps between primes,''
Ann.\ Math.\ \textbf{179}, 1121--1174 (2014).

\bibitem{Polymath2014}
D.~H.~J.~Polymath,
``New equidistribution estimates of Zhang type,''
Algebra Number Theory \textbf{8}, 2067--2199 (2014).

\bibitem{Maynard2015}
J.~Maynard,
``Small gaps between primes,''
Ann.\ Math.\ \textbf{181}, 383--413 (2015).

\bibitem{Apostol1976}
T.~M.~Apostol,
\emph{Introduction to Analytic Number Theory}
(Springer, New York, 1976).

\bibitem{HardyWright2008}
G.~H.~Hardy and E.~M.~Wright,
\emph{An Introduction to the Theory of Numbers}, 6th ed., revised by
D.~R.~Heath-Brown and J.~H.~Silverman
(Oxford University Press, Oxford, 2008).

\item[]\textbf{H. Anomalies, CP violation, strong CP, and electric dipole moments}

\bibitem{Adler1969}
S.~L.~Adler,
``Axial-vector vertex in spinor electrodynamics,''
Phys.\ Rev.\ \textbf{177}, 2426--2438 (1969).

\bibitem{BellJackiw1969}
J.~S.~Bell and R.~Jackiw,
``A PCAC puzzle: $\pi^0\to\gamma\gamma$ in the sigma model,''
Nuovo Cim.\ A \textbf{60}, 47--61 (1969).

\bibitem{tHooft1976}
G.~'t~Hooft,
``Symmetry breaking through Bell--Jackiw anomalies,''
Phys.\ Rev.\ Lett.\ \textbf{37}, 8--11 (1976).

\bibitem{PecceiQuinn1977}
R.~D.~Peccei and H.~R.~Quinn,
``CP conservation in the presence of pseudoparticles,''
Phys.\ Rev.\ Lett.\ \textbf{38}, 1440--1443 (1977).

\bibitem{Weinberg1978Axion}
S.~Weinberg,
``A new light boson?'' Phys.\ Rev.\ Lett.\ \textbf{40}, 223--226 (1978).

\bibitem{Wilczek1978Axion}
F.~Wilczek,
``Problem of strong $P$ and $T$ invariance in the presence of instantons,''
Phys.\ Rev.\ Lett.\ \textbf{40}, 279--282 (1978).

\bibitem{crewther1979}
R.~J.~Crewther, P.~Di~Vecchia, G.~Veneziano, and E.~Witten,
``Chiral estimate of the electric dipole moment of the neutron in quantum
chromodynamics,'' Phys.\ Lett.\ B \textbf{88}, 123--127 (1979).

\bibitem{witten1982}
E.~Witten,
``An SU(2) anomaly,'' Phys.\ Lett.\ B \textbf{117}, 324--328 (1982).

\bibitem{pospelov2005}
M.~Pospelov and A.~Ritz,
``Electric dipole moments as probes of new physics,''
Ann.\ Phys.\ (N.Y.) \textbf{318}, 119--169 (2005).

\item[]\textbf{I. Gravitation, cosmology, horizons, and information}

\bibitem{Einstein1915}
A.~Einstein,
``Die Feldgleichungen der Gravitation,''
Sitzungsber.\ Preuss.\ Akad.\ Wiss.\ Berlin 844--847 (1915).

\bibitem{Friedmann1922}
A.~Friedmann,
``\"Uber die Kr\"ummung des Raumes,''
Z.\ Phys.\ \textbf{10}, 377--386 (1922).

\bibitem{Lemaitre1927}
G.~Lema\^itre,
``Un univers homog\`ene de masse constante et de rayon croissant, rendant
compte de la vitesse radiale des n\'ebuleuses extra-galactiques,''
Ann.\ Soc.\ Sci.\ Bruxelles A \textbf{47}, 49--59 (1927).

\bibitem{Hubble1929}
E.~Hubble,
``A relation between distance and radial velocity among extra-galactic
nebulae,'' Proc.\ Natl.\ Acad.\ Sci.\ USA \textbf{15}, 168--173 (1929).

\bibitem{Penrose1965}
R.~Penrose,
``Gravitational collapse and space-time singularities,''
Phys.\ Rev.\ Lett.\ \textbf{14}, 57--59 (1965).

\bibitem{Bekenstein1973}
J.~D.~Bekenstein,
``Black holes and entropy,'' Phys.\ Rev.\ D \textbf{7}, 2333--2346 (1973).

\bibitem{Hawking1975}
S.~W.~Hawking,
``Particle creation by black holes,''
Commun.\ Math.\ Phys.\ \textbf{43}, 199--220 (1975).

\bibitem{HawkingEllis1973}
S.~W.~Hawking and G.~F.~R.~Ellis,
\emph{The Large Scale Structure of Space-Time}
(Cambridge University Press, Cambridge, 1973).

\bibitem{Wald1984}
R.~M.~Wald,
\emph{General Relativity}
(University of Chicago Press, Chicago, 1984).

\bibitem{Weinberg1989CC}
S.~Weinberg,
``The cosmological constant problem,''
Rev.\ Mod.\ Phys.\ \textbf{61}, 1--23 (1989).

\bibitem{Jacobson1995}
T.~Jacobson,
``Thermodynamics of spacetime: The Einstein equation of state,''
Phys.\ Rev.\ Lett.\ \textbf{75}, 1260--1263 (1995).

\bibitem{Carroll2001}
S.~M.~Carroll,
``The cosmological constant,''
Living Rev.\ Relativ.\ \textbf{4}, 1 (2001).

\bibitem{Verlinde2011}
E.~P.~Verlinde,
``On the origin of gravity and the laws of Newton,''
J.\ High Energy Phys.\ \textbf{04}, 029 (2011).

\item[]\textbf{J. Scientific methodology, falsifiability, preregistration, and reproducibility}

\bibitem{Popper1959}
K.~R.~Popper,
\emph{The Logic of Scientific Discovery}
(Hutchinson, London, 1959).

\bibitem{Lakatos1978}
I.~Lakatos,
\emph{The Methodology of Scientific Research Programmes}
(Cambridge University Press, Cambridge, 1978).

\bibitem{Ioannidis2005}
J.~P.~A.~Ioannidis,
``Why most published research findings are false,''
PLoS Med.\ \textbf{2}, e124 (2005).

\bibitem{Sandve2013}
G.~K.~Sandve, A.~Nekrutenko, J.~Taylor, and E.~Hovig,
``Ten simple rules for reproducible computational research,''
PLoS Comput.\ Biol.\ \textbf{9}, e1003285 (2013).

\bibitem{Wilson2014}
G.~Wilson \emph{et al.},
``Best practices for scientific computing,''
PLoS Biol.\ \textbf{12}, e1001745 (2014).

\bibitem{Munafo2017}
M.~R.~Munaf\`o \emph{et al.},
``A manifesto for reproducible science,''
Nat.\ Hum.\ Behav.\ \textbf{1}, 0021 (2017).

\bibitem{Nosek2018}
B.~A.~Nosek, C.~R.~Ebersole, A.~C.~DeHaven, and D.~T.~Mellor,
``The preregistration revolution,''
Proc.\ Natl.\ Acad.\ Sci.\ USA \textbf{115}, 2600--2606 (2018).

\bibitem{Stark2018}
P.~B.~Stark,
``Before reproducibility must come preproducibility,''
Nature \textbf{557}, 613 (2018).



\item[]\textbf{K. Model selection, likelihoods, and reporting standards}

\bibitem{Akaike1974}
H.~Akaike,
``A new look at the statistical model identification,''
IEEE Trans. Automat. Control \textbf{19}, 716--723 (1974).

\bibitem{Schwarz1978}
G.~Schwarz,
``Estimating the dimension of a model,''
Ann. Stat. \textbf{6}, 461--464 (1978).

\bibitem{Wilks1938}
S.~S.~Wilks,
``The large-sample distribution of the likelihood ratio for testing composite
hypotheses,'' Ann. Math. Stat. \textbf{9}, 60--62 (1938).

\bibitem{Cowan2011}
G.~Cowan, K.~Cranmer, E.~Gross, and O.~Vitells,
``Asymptotic formulae for likelihood-based tests of new physics,''
Eur. Phys. J. C \textbf{71}, 1554 (2011).

\bibitem{Humphreys1990}
J.~E.~Humphreys,
\emph{Reflection Groups and Coxeter Groups}
(Cambridge University Press, Cambridge, 1990).

\end{thebibliography}

% ---- END TIR/monograph/references_expanded_v10_8.tex ----


\vfill
\noindent\small\textit{Source code and version history:}
\url{https://github.com/AdrianLipa90/The-Fundamental-Theory-of-Informational-Relations}
\par\vspace{2cm}
\backmatter
\printindex
\end{document}
